% !TEX root = ../../tree_notes.tex
% =====================================================================
%  Chapter 24 -- Centroid Decomposition                          WRITTEN
%  (folder ch18_centroid_decomposition; label prefix c18; PRINTS as
%  Chapter 24; folder number 18; build with:  python tools/build.py --only 18)
%
%  NUMBERING TRAP (the Ch.22/Ch.23 trap, mirrored): \chref{c18} prints
%  "Chapter 24" (THIS chapter).  Do NOT confuse the label with the printed
%  number: \chref{c14} prints "Chapter 18" (Fenwick), \chref{c24} prints
%  "Chapter 5" (Euler tour), \chref{c16} prints "Chapter 22" (HLD),
%  \chref{c29} prints "Chapter 23" (Li Chao), \chref{c23} prints "Chapter 25"
%  (link-cut), \chref{c21} prints "Chapter 27" (Mo), \chref{c09} prints
%  "Chapter 12" (Height/Diameter/Centroid), \chref{c28} prints "Chapter 17"
%  (DSU/sack).  Every cite of THIS chapter is \chref{c18}.
%
%  Owns CENTROID DECOMPOSITION -- the remove-the-centroid-and-recurse engine,
%  the centroid TREE (parent par_c, height <= ceil(log2 n)), counting/aggregate
%  over ALL paths by splitting each path at the highest centroid on it, the
%  "subtract the same-child-subtree pairs" correction, the per-centroid DISTANCE
%  TABLE that indexes a whole component from a single touched node, and the
%  divide-and-conquer-on-trees template beyond plain path sums.
%  It STARTS where Ch.12 12.3 ENDS: Ch.12 owns the centroid DEFINITION, the
%  existence proof (thm:c09-centexist) and the linear find (thm:c09-centfind,
%  lst:c09-cent); this chapter never re-derives those -- it back-cites them and
%  begins at "remove it and recurse."
%  Sections:
%    * 24.1 building the centroid tree -- the recursive solve(S): find centroid
%      of a component, mark it deleted, recurse on each surviving piece; why
%      every level <= n/2 gives O(log n) levels and O(n log n) build; par_c and
%      cd-depth; the deleted-marker discipline (never re-enter a removed node).
%      Hand trace decomposes the shared tree T (15 nodes, top centroid 4).
%      Solved CF 321C "Ciel the Commander" (2100, L4) -- output the centroid tree:
%      rank each node by its centroid depth ('A' = top centroid); the build IS the solve.
%    * 24.2 counting paths by distance -- every path is claimed by the highest
%      centroid it passes through; count pairs at distance <= K / == K with a
%      distance histogram per subtree; the subtract-same-child correction; the
%      sort-and-two-pointer merge in O(size log size).  DISCHARGES the forward-
%      cites from Ch.22 (05_boundary "path COUNTING by centroid"), Ch.23 (99
%      "over all paths is centroid's job") and Ch.16 LCA.  CF 161D "Distance in
%      Tree" (=K) is OWNED by Ch.10 as an O(nK) DP -- PROSE ONLY, never re-linked.
%      Solved CF 715C "Digit Tree" (count ordered pairs, distance as a base-10
%      number divisible by m) -- the centroid pair-count template + modular filter.
%    * 24.3 the centroid tree as an index -- one distance table per centroid;
%      touch a node = O(log n) table writes up its centroid ancestors; answer a
%      nearest / farthest / aggregate query by walking the same ancestors. The
%      dynamic nearest-marked pattern (Xenia and Tree) is taught here in prose; its
%      URL is OWNED by Ch.12 12.3 (342/E), so it is named, never re-linked.  Solved
%      CF 293E "Close Vertices" (count pairs whose unique path has at most l edges
%      AND total weight at most w) -- the per-centroid (length, weight) distance
%      table, merged by a sort-on-length sweep with a Fenwick over weights.
%    * 24.4 divide and conquer on trees beyond paths -- combine a Fenwick / sorted
%      structure inside each centroid bucket (back-cites Ch.18 Fenwick, c14);
%      aggregates over all paths that a single tag cannot fold.  Solved CF 833D
%      "Red-Black Cobweb" (count colour-balanced paths via centroid + an arm
%      two-pointer) -- the general "recurse at the centroid, merge arms, prune" template.
%    * 24.5 the engine boundary -- when centroid decomposition is the WRONG tool:
%      path RANGE queries on a fixed tree -> HLD (c16); subtree multiset merge ->
%      DSU-on-tree/sack (c28); a state per node no distance needs -> tree DP /
%      rerooting (c20/c10); offline range-of-node-set queries -> Mo (c21); dynamic
%      topology -> link-cut (c23).  Capstone PICK table def:c18-map.  WORKED
%      examples only -- NO solvedproblem (in-section count stays 4).
%  Chapter-end Summit: CF 757G "Can Bash Save the Day?" (3400) -- online sum of
%  distances from v to a permutation-index RANGE under adjacent swaps; centroid
%  decomposition + a Fenwick per centroid over positions with the child-
%  subtraction, reusing Ch.18 Fenwick (c14).  Sample hand-replayed (XOR-decoded
%  queries -> 23, 37, 28).
%
%  Diagram-first: this chapter is TREE based.  The running example is the shared
%  15-node tree T (defined once, authoritatively, in 00_opener; every section re-
%  uses it exactly).  Edges (1-based):
%    (4,1) (4,2) (4,3) (1,5) (1,6) (6,11) (2,7) (2,8) (8,12)
%    (3,9) (3,10) (9,13) (10,14) (14,15)
%  The authoritative centroid decomposition (hand-verified in 24.1):
%    whole tree n=15 -> centroid 4 (remove4 pieces sizes {4,4,6}, max 6 <= 7);
%    A={1,5,6,11} -> centroid 1 (tie with 6, pick 1); 5->5, {6,11}->6->11;
%    B={2,7,8,12} -> centroid 2 (tie with 8, pick 2); 7->7, {8,12}->8->12;
%    C={3,9,10,13,14,15} -> centroid 3 (tie with 10, pick 3); {9,13}->9->13;
%       {10,14,15} path -> centroid 14, then 10 and 15 as its children;
%    centroid tree (par_c, 0 = root):  par_c(4)=0, par_c(1..3)=4,
%       par_c(5)=1, par_c(6)=1, par_c(11)=6, par_c(7)=2, par_c(8)=2,
%       par_c(12)=8, par_c(9)=3, par_c(13)=9, par_c(14)=3, par_c(10)=14,
%       par_c(15)=14;  centroid depth = 0(4),1(1,2,3),2(5,6,7,8,9,14),
%       3(11,12,13,10,15); max depth 3 <= ceil(log2 15)=4.
%
%  Anti-repetition contract (teach once, cross-reference by label):
%    * Ch.12 12.3 (\chref{c09}) owns the centroid DEFINITION, existence and
%      linear find; 24.1 REUSES thm:c09-centfind as a black box -- the "walk to
%      the heavy side" is never re-derived here.
%    * Ch.13 (\chref{c09}) diameter is cited (never recomputed) where "farthest"
%      via centroid tables needs it.
%    * Ch.18 Fenwick (\chref{c14}) owns the BIT; 24.4/Summit bolt one onto each
%      centroid bucket and back-cite it -- the BIT is never re-taught.
%    * Ch.22 HLD (\chref{c16}) owns PATH ranges; 24.2/24.3 handle COUNTS over all
%      paths and node-SET distance indexes, a different job; the boundary is drawn
%      in 24.5 and both back-cite \secref{c16-boundary}.
%    * Ch.23 Li Chao (\chref{c29}) is cited where a centroid bucket needs a max
%      over lines; not required by this chapter's four solved problems.
%    * Explicitly OUT (prose + forward-cite only, NO code): dynamic topology ->
%      link-cut Ch.25 (\chref{c23}); online path-RANGE aggregates on a fixed tree
%      -> HLD Ch.22 (\chref{c16}); subtree multiset union -> DSU-on-tree Ch.17
%      (\chref{c28}).
%
%  Known URL collisions (confirmed by grep of chapters/; prose-cite only,
%  NEVER re-link -- urlsweep must print an empty DUPLICATES section):
%    * CF 342E "Xenia and Tree" (2400) -> Ch.12 (\chref{c09}) 12.3 practice.  The
%      canonical dynamic nearest-marked centroid task; PROSE + pattern only here.
%    * CF 161D "Distance in Tree" (1800) -> Ch.10 (\chref{c10}) tree-DP.  The =K
%      counting warm-up; re-treated as "the centroid upgrade" in PROSE, never re-linked.
%    * CF 914D, 1601C, 1743F and the whole Ch.19/20/21/22/23 rosters above are
%      untouched by c18.  Ch.23 owns the Li Chao/ODT family; Ch.12 owns 592D/337D/
%      1092F/219D; none is re-linked here.
%
%  Roster (every title/rating/tag/sample live-verified against the Codeforces
%  problemset API / problem page BEFORE a trace was written; Codeforces-forward):
%    solved: CF 321C Ciel the Commander          (2100, L4) (24.1)
%            CF 715C Digit Tree                (2700, L5) (24.2)
%            CF 293E Close Vertices            (2700, L5) (24.3)
%            CF 833D Red-Black Cobweb          (2800, L5) (24.4)
%    worked: the centroid decomposition of T (24.1); a distance-k pair count on
%            T (24.2); a within-(length, weight) pair count on T (24.3); the colour-
%            balanced arm merge (24.4); the engine pick table (24.5).  CF 342E and
%            161D are discussed in PROSE only (their URLs live in Ch.12 and Ch.10).
%    summit: CF 757G Can Bash Save the Day?    (3400, L6) -- weighted tree +
%            permutation a; query 1 l r v = sum over i in [l,r] of dist(a_i,v);
%            query 2 x = swap a_x,a_{x+1}; queries XOR-encoded against the
%            previous answer (online).  Centroid decomposition with a Fenwick per
%            centroid over positions (child-subtree subtraction).  Sample
%            (n=5,q=5, a=4 5 1 3 2) hand-decoded and replayed -> 23, 37, 28.
%    mixed (10, ascending by live-verified rating, every ID absent from the owned
%            lists above -- locked at the 99_chapter_end write): see that file.  CF
%            990G GCD Counting (2400, centroid + Mobius pair count) is a mixed-set entry.
% =====================================================================
\bookchapter{Centroid Decomposition}{c18}
\chapterinfo{CF 2100--3400}{Written}

% =====================================================================
%  Chapter 24 opener -- "Chapter at a Glance", three reading paths + the
%  ladder staircase, THE running example (the shared 15-node tree T, top
%  centroid 4), the who-owns-what comparison against Ch.12 (centroid find),
%  Ch.13 (=K DP), Ch.18 Fenwick, Ch.22 HLD and Ch.17 sack, the recognition
%  guide, a constraint->technique table, and notation.
%  Exports fig:c18-tree, fig:c18-taxonomy.
%  No \term{centroid} or \term{centroid decomposition} here -- both are
%  indexed by Ch.12 (\chref{c09}); \term the NEW object "centroid tree" only.
% =====================================================================

\begin{chaptertheorybox}{Chapter at a Glance}
Chapter~12 taught you to \emph{find} a centroid: the node whose removal leaves
no piece bigger than \(n/2\), located in one linear walk toward whatever side is
heavy (\(\thmref{c09-centfind}\)).  That chapter stopped on purpose at "one
centroid, and why it is safe to cut here."  This chapter cashes that promise.  Do
not stop at one cut: cut at the centroid, then \emph{recurse} on every leftover
piece, and you get a divide-and-conquer whose recursion tree has only
\(O(\log n)\) levels -- because each cut leaves pieces at most half the size.  That
recursion is \emph{centroid decomposition}, and the tree it produces (this
chapter's new object, the \term{centroid tree}) is the load-bearing structure of
the whole chapter.

\textbf{You will learn to}
\begin{itemize}[nosep]
  \item build the \term{centroid tree} of a tree in \(\bigO{n\log n}\) -- find the
        centroid of a component, mark it deleted, recurse on the surviving
        pieces; prove the height is \(\le\lceil\log_2 n\rceil\) because every level
        at least halves a component (\(\secref{c18-tree}\), solved CF~321C); \hfill\lv{3}--\lv{4}
  \item \emph{claim every path by its top centroid} and count pairs at distance
        \(\le K\) (or \(=K\)) with a per-subtree distance histogram and the
        "subtract the same-child pairs" correction (\(\secref{c18-count}\),
        solved CF~715C); \hfill\lv{4}--\lv{5}
  \item use one centroid as a \emph{distance table} for its whole component: touch
        a node in \(O(\log n)\) writes up its centroid ancestors, then answer
        nearest / farthest / aggregate queries by walking the same ancestors
        (\(\secref{c18-index}\), solved CF~293E); \hfill\lv{5}
  \item run divide-and-conquer on trees \emph{beyond} plain path sums -- merge
        "arms" at the centroid, prune by an invariant, and let each bucket carry a
        Fenwick or sorted array (back-citing Ch.~18 \chref{c14})
        (\(\secref{c18-beyond}\), solved CF~833D); \hfill\lv{5}
  \item read a specification and decide whether the job is centroid's at all, or
        belongs to HLD / DSU-on-tree / tree DP / Mo / link-cut
        (\(\secref{c18-boundary}\), the pick table). \hfill\lv{5}
\end{itemize}

\textbf{Prerequisites.} Chapter~12 (\chref{c09}) -- the centroid definition and the
linear find; this chapter treats \(\thmref{c09-centfind}\) as a black box and never
re-derives it.  Chapter~13 (\chref{c10}) -- the \(O(nK)\) subtree DP that counts
distance-\(K\) pairs, the rival this chapter beats to \(O(n\log n)\) for large
\(K\).  Chapter~16 (\chref{c11}) -- \(dist(u,v)\) by depth and LCA, used to turn a
path count into a distance count.  Chapter~18 (\chref{c14}) -- the Fenwick tree
that rides inside a centroid bucket (\(\secref{c18-beyond}\) and the Summit).
\end{chaptertheorybox}

% ---------------- Reading paths + ladder position ----------------
\begin{chaptertheorybox}{How to Read This Chapter (three paths)}
Level badges \(\lv3\)--\(\lv6\): \(\lv3\) is decomposing the running tree \(T\) by
hand, cutting at the centroid and recursing; \(\lv{4}\) is the \(O(n\log n)\) build
written with a correct deleted-marker discipline, or a distance-\(K\) histogram with
the same-subtree subtraction; \(\lv{5}\) is a per-centroid distance table driving an
online nearest/aggregate query, or an arm-merge with a pruning invariant; \(\lv{6}\)
is the Summit, a centroid decomposition with a Fenwick living inside every bucket.

\textbf{Build-and-count path (\(\lv3\)--\(\lv{5}\), "count over all paths").}
\secref{c18-tree} in full -- the recursion, the \(\bigO{\log n}\)-level proof, solved
CF~321C -- then \secref{c18-count}, where paths are claimed by their top centroid and
counted by distance histogram (CF~715C), capped by the Summit \(\probref{c18-summit}\).

\textbf{Distance-index path (\(\lv{4}\)--\(\lv{6}\), "one node touches \(\log n\) tables").}
\secref{c18-index} -- the per-centroid distance table, used offline to count pairs
within a distance bound (solved CF~293E: \(\le l\) edges and \(\le w\) total weight) and,
with updates, for the online nearest-marked pattern (the classic Xenia and Tree, whose
statement Ch.~12 previews, named only);
then \secref{c18-beyond}, where a bucket carries a Fenwick (Ch.~18 \chref{c14}) --
solved CF~833D, and the Summit.

\textbf{Boundary path (\(\lv{4}\)--\(\lv{5}\), "is this centroid's job?").}
\secref{c18-boundary}: the pick table that hands path \emph{range} queries to HLD
(\chref{c16}), subtree multiset unions to DSU-on-tree (\chref{c28}), per-node states
to tree DP (\chref{c10}), offline node-set ranges to Mo (\chref{c21}), and mutable
topology to link-cut (\chref{c23}).
\end{chaptertheorybox}

\begin{figure}[H]
\centering
\begin{tikzpicture}[diagram, x=8.5mm, y=6.5mm]
  \foreach \st/\lab in {1/lv1,2/lv2,3/lv3,4/lv4,5/lv5,6/lv6} {
    \fill[\lab!22, draw=\lab, line width=0.8pt] (\st-1,0) rectangle (\st,\st*0.52);
    \node[font=\scriptsize\bfseries, text=\lab] at (\st-0.5,\st*0.52-0.42) {L\st};}
  \node[annot, anchor=south west, font=\scriptsize]
    at (0.05,0.62) {prereq (Ch.12): find one centroid, linear walk};
  \node[annot, anchor=south west, font=\scriptsize]
    at (1.05,1.14) {24.1 cut at 4, recurse; O(n log n); parent array, depth log n};
  \node[annot, anchor=south west, font=\scriptsize]
    at (2.05,1.66) {24.1 solved CF 321C (output centroid tree); 24.2 distance histogram, subtract correction};
  \node[annot, anchor=south west, font=\scriptsize]
    at (3.05,2.18) {24.2 solved CF 715C; 24.3 per-centroid distance tables};
  \node[annot, anchor=south west, font=\scriptsize, text=lv5]
    at (4.05,2.70) {24.3 solved CF 293E; 24.4 colour/arm aggregates, solved CF 833D};
  \node[annot, anchor=south west, font=\scriptsize, text=lv6]
    at (5.05,3.22) {24.5 the pick table -- is it centroid's job; Summit CF 757G};
\end{tikzpicture}
\caption*{Chapter~24 on the six-step ladder: decomposing \(T\) by hand is
\(\lv3\); a correct build and a distance-\(K\) count are \(\lv{4}\); per-centroid
distance tables and an arm-merge with a Fenwick bucket are \(\lv{5}\); the Summit
(centroid decomposition with a Fenwick inside every centroid) is step~6.}
\end{figure}

% ---------------- The running example ----------------
\subsection*{The Running Example: One Tree, Decomposed Over and Over}

Every section in this chapter runs on the same unweighted tree \(T\) on
\(n=15\) nodes, fixed once here and authoritative (if a later trace disagrees,
this figure wins).  Edges, written as pairs:
\[
  (4,1)\ (4,2)\ (4,3)\ (1,5)\ (1,6)\ (6,11)\ (2,7)\ (2,8)\ (8,12)\ (3,9)\ (3,10)\ (9,13)\ (10,14)\ (14,15).
\]
Its centroid is \(4\): removing it leaves three pieces \(\{1,5,6,11\}\),
\(\{2,7,8,12\}\), \(\{3,9,10,13,14,15\}\) of sizes \(4,4,6\), all \(\le 7\); every
other node leaves a piece of \(9\) or more, so \(4\) is the \emph{unique} top
centroid.  Recurse into those three pieces and you get the \term{centroid tree}
drawn in \(\secref{c18-tree}\).

\begin{figure}[H]
\centering
\begin{tikzpicture}[diagram, scale=0.9, transform shape,
   tnode/.style={circle, draw=primaryblue!70, fill=primaryblue!8, line width=0.7pt,
                 minimum size=6.5mm, font=\scriptsize\bfseries, inner sep=1pt},
   cen/.style={circle, draw=warningred, fill=warningred!18, line width=1.1pt,
               minimum size=7mm, font=\footnotesize\bfseries, inner sep=1pt}]
  \node[cen] (n4) at (0,0) {4};
  \node[tnode] (n1) at (-5,1.7) {1};
  \node[tnode] (n2) at (0,2.6) {2};
  \node[tnode] (n3) at (5,1.7) {3};
  \node[tnode] (n5) at (-6.6,3.2) {5};
  \node[tnode] (n6) at (-4.2,3.2) {6};
  \node[tnode] (n11) at (-4.2,4.7) {11};
  \node[tnode] (n7) at (-1.3,4.1) {7};
  \node[tnode] (n8) at (1.3,4.1) {8};
  \node[tnode] (n12) at (1.3,5.6) {12};
  \node[tnode] (n9) at (3.4,3.2) {9};
  \node[tnode] (n13) at (3.4,4.7) {13};
  \node[tnode] (n10) at (6.6,3.2) {10};
  \node[tnode] (n14) at (6.6,4.7) {14};
  \node[tnode] (n15) at (8.1,5.6) {15};
  \draw (n4)--(n1) (n4)--(n2) (n4)--(n3);
  \draw (n1)--(n5) (n1)--(n6) (n6)--(n11);
  \draw (n2)--(n7) (n2)--(n8) (n8)--(n12);
  \draw (n3)--(n9) (n9)--(n13);
  \draw (n3)--(n10) (n10)--(n14) (n14)--(n15);
  \draw[warningred, line width=1.3pt, densely dashed] (-7.2,2.35) -- (-2.0,2.35);
  \draw[warningred, line width=1.3pt, densely dashed] (-2.6,3.05) -- (2.6,3.05);
  \draw[warningred, line width=1.3pt, densely dashed] (2.0,2.35) -- (7.6,2.35);
\end{tikzpicture}
\caption{The running tree \(T\) (\(n=15\)) and its centroid \(4\) (red).  Cutting
at \(4\) -- the dashed cuts -- splits \(T\) into the three pieces of sizes
\(4,4,6\); each of size \(\le \lfloor 15/2\rfloor\).  Those three pieces are
recursed into independently in \(\secref{c18-tree}\), producing the centroid tree.
Left piece \(A=\{1,5,6,11\}\), middle \(B=\{2,7,8,12\}\), right
\(C=\{3,9,10,13,14,15\}\).}
\label{fig:c18-tree}
\end{figure}

% ---------------- Which job, cost chart ----------------
\subsection*{Which Job? (the question-shape chart)}

\begin{center}
\small
\renewcommand{\arraystretch}{1.25}
\setlength{\tabcolsep}{3pt}
\begin{tabular}{@{}p{0.30\linewidth} p{0.17\linewidth} p{0.29\linewidth} >{\raggedright\arraybackslash}p{0.14\linewidth}@{}}
\toprule
\textbf{Question} & \textbf{Technique} & \textbf{Why not a single DFS/DP} & \textbf{Where} \\
\midrule
count pairs of nodes at distance \(=K\)/\(\le K\) & top-centroid claim + histogram & a rooted DP is \(O(nK)\); centroid is \(O(n\log n)\) & \secref{c18-count} \\
for every node, nearest node with property \(P\) & per-centroid distance table & no single root holds all-pairs distances & \secref{c18-index} \\
online mark a node / query nearest marked & centroid tree index & the marked set changes; \(O(\log n)\) tables per touch & \secref{c18-index} \\
derive the centroid tree / rank every node by its centroid depth & build by cut-and-recurse & it is the decomposition itself, not one rooted DP & \secref{c18-tree} \\
sum of distances from \(v\) to a set / range & centroid bucket + Fenwick & the set is dynamic; \(\log n\) buckets touched per node & \secref{c18-beyond} \\
aggregate over \emph{every} simple path & \emph{divide and conquer on trees} & the path, not the node, is the unit of work & \secref{c18-count} \\
path \emph{range} update/query, fixed tree & \emph{go to HLD} & chains answer range-on-path, not all-pairs & \chref{c16} \\
\bottomrule
\end{tabular}
\end{center}

\begin{insightbox}[One promise decides everything: split where the tree balances]
A rooted DP fixes a root and combines children; it cannot see a path that wanders
up through the root and back down without counting it once per root, and it has no
handle on "the set of all paths."  Centroid decomposition changes the unit of work
from \emph{node} to \emph{path}: a path is charged to the highest centroid it
touches, and at that centroid the path has a fixed midpoint, so its length is
just \(depth_u + depth_v\) over two independent arms.  Because every cut leaves
pieces \(\le n/2\), the recursion is shallow (\(O(\log n)\)) yet each node
participates in only \(O(\log n)\) centroid buckets -- which is what turns an
\(O(n^2)\) all-pairs job into \(O(n\log^2 n)\), and lets a whole component be
\emph{indexed} from one node (\(\secref{c18-index}\)).
\end{insightbox}

% ---------------- Taxonomy chart ----------------
\begin{figure}[H]
\centering
\begin{tikzpicture}[diagram, scale=0.95, transform shape,
    mnode/.style={circle, draw=#1, line width=0.9pt, fill=#1!8, minimum size=8mm,
                  font=\scriptsize\bfseries, text=#1!60!black}]
  \node[mnode=primaryblue] (a) at (0,0)   {cut};
  \node[annot, font=\scriptsize, anchor=west, text width=9.6cm, align=left] at (1.4,0) {find a centroid (Ch.12's find), delete it, recurse on each piece; the parent array records the parent centroid, depth is log-bounded (24.1)};
  \node[mnode=successgreen] (b) at (0,-1) {claim};
  \node[annot, font=\scriptsize, anchor=west, text width=9.6cm, align=left] at (1.4,-1) {a path is charged to the highest centroid on it; there it is depth(u)+depth(v) across two arms -- count by distance histogram (24.2, CF 715C)};
  \node[mnode=difforange]   (c) at (0,-2) {index};
  \node[annot, font=\scriptsize, anchor=west, text width=9.6cm, align=left] at (1.4,-2) {each centroid keeps a distance table for its component; touch a node = O(log n) writes up its centroid ancestors, walk them to answer nearest/farthest (24.3, CF 293E)};
  \node[mnode=warningred]   (d) at (0,-3) {merge};
  \node[annot, font=\scriptsize, anchor=west, text width=9.6cm, align=left] at (1.4,-3) {beyond counting: merge arms at the centroid, prune by an invariant, and let a bucket hold a Fenwick or sorted array (Ch.18) (24.4, CF 833D)};
  \node[mnode=softpurple]   (e) at (0,-4) {pick};
  \node[annot, font=\scriptsize, anchor=west, text width=9.6cm, align=left] at (1.4,-4) {is it centroid's job? all-pairs/path-statistic yes; path ranges -> HLD; subtree unions -> sack; per-node states -> DP; offline sets -> Mo; dynamic edges -> link-cut (24.5)};
\end{tikzpicture}
\caption{The chapter's five-state taxonomy.  \emph{cut} is the decomposition
itself; \emph{claim} and \emph{index} are the two ways it pays off (count paths by
their top centroid; index a component by its centroid distance table); \emph{merge}
is the general divide-and-conquer-on-trees template with a data structure in the
bucket; \emph{pick} is the boundary choosing between centroid and its neighbours.
Each solved problem sits on one row: CF~321C on \emph{cut}, CF~715C on \emph{claim},
CF~293E on \emph{index}, CF~833D on \emph{merge}, the Summit CF~757G on
\emph{index}+\emph{merge}.}
\label{fig:c18-taxonomy}
\end{figure}

% ---------------- Who owns what ----------------
\subsection*{Who Owns What (so nothing is learned twice)}

\begin{center}
\small
\renewcommand{\arraystretch}{1.25}
\begin{tabular}{@{}p{0.30\linewidth} p{0.24\linewidth} p{0.40\linewidth}@{}}
\toprule
\textbf{Idea} & \textbf{Owner} & \textbf{This chapter's share} \\
\midrule
centroid definition + linear find & Ch.~12 (\chref{c09}) & treated as a black box; 24.1 begins at "delete it and recurse", never re-derives existence \\
count distance-\(K\) pairs by subtree DP & Ch.~13 (\chref{c10}) & the \(O(nK)\) rival (\(\probref{c10-harvest-solved}\)); 24.2 upgrades it to \(O(n\log n)\) and explains when the DP still wins \\
diameter / farthest by two DFS & Ch.~12 (\chref{c09}) & cited, never recomputed, where a centroid farthest-table needs a baseline \\
Fenwick / prefix sums & Ch.~18 (\chref{c14}) & 24.4 and the Summit place one BIT per centroid bucket and reuse it verbatim \\
path = \(O(\log n)\) chain ranges & Ch.~22 (\chref{c16}) & contrasted in 24.5: HLD answers path \emph{ranges}, centroid answers counts over \emph{all} paths \\
subtree multiset union (sack) & Ch.~17 (\chref{c28}) & forward-cited rival for "merge the children's sets"; centroid is for path/whole-component work \\
\(dist(u,v)\) by depth and LCA & Ch.~16 (\chref{c11}) & 24.2 uses it to turn "distance \(K\)" into "sum of two depths at the top centroid" \\
dynamic topology (link/cut) & Ch.~25 (\chref{c23}) & explicitly out; centroid is built once on a fixed tree \\
\bottomrule
\end{tabular}
\end{center}

% ---------------- Recognition guide ----------------
\begin{recognitionbox}[Pattern card -- "a claim about paths, or about distance from a set"]
\reckey{statements that ask for something over \emph{all pairs} or \emph{all simple
paths}: ``count pairs of nodes at distance \(K\) (or \(\le K\), or divisible by
\(m\)'', ``find the longest/shortest path with a property that a subtree DP cannot
localise'', ``for each node, the distance to the nearest node of colour/type \(C\)'',
``maintain a dynamic set of marked nodes and answer nearest-marked after each
update'', ``sum of distances from \(v\) to all nodes in a set/range''.  The tell is
that the unit of work is a \emph{path} (two endpoints) or a \emph{whole-component
distance}, not a single rooted subtree.}
\recuse{
  \item count / aggregate over all paths by a property additive along the path -- claim by top centroid (\(\secref{c18-count}\));
  \item nearest / farthest node from a moving set, online -- per-centroid distance index (\(\secref{c18-index}\));
  \item a path statistic a rooted DP cannot fold (a path-gcd count, a colour-balanced path count) -- arm-merge D\&C on trees (\(\secref{c18-beyond}\)).
}
\recnotuse{
  \item a range \emph{update and query along one path} on a fixed tree -- HLD (\chref{c16});
  \item a value that is a \emph{state per node} with no distance dimension -- tree DP / rerooting (\chref{c10}, \chref{c20});
  \item \emph{merge the children's subtrees' sets} bottom-up -- DSU-on-tree/sack (\chref{c28});
  \item many offline queries over a node set with free order -- Mo on trees (\chref{c21});
  \item the tree's \emph{edges change} (link/cut) -- link-cut / Euler-tour trees (\chref{c23}).
}
One-line test: \emph{the problem quantifies over paths or needs a component-wide
distance reachable from a single touched node -- then centroid decomposition, not a
rooted DP and not a path-chain engine.}
\end{recognitionbox}

% ---------------- Constraint -> technique ----------------
\subsection*{Reading the Constraints (which technique the limits want)}

\begin{center}
\small
\renewcommand{\arraystretch}{1.25}
\begin{tabular}{@{}p{0.34\linewidth} p{0.40\linewidth} p{0.20\linewidth}@{}}
\toprule
\textbf{Constraint signature} & \textbf{What it rules in/out} & \textbf{Technique} \\
\midrule
count pairs at distance \(=K\), \(n\le2\cdot10^5\), \(K\) large & the \(O(nK)\) subtree DP times out; the path has no fixed root & \secref{c18-count} \\
for each node nearest same-colour node & all-pairs is \(O(n^2)\); a centroid table answers in \(O(\log n)\) probes & \secref{c18-index} \\
online paint black / query nearest black, \(Q\le10^5\) & each touch writes \(O(\log n)\) buckets; query walks them & \secref{c18-index} \\
longest path with a global invariant (gcd \(>1\)) & not a composable fold; merge arms at each centroid and prune & \secref{c18-beyond} \\
sum of distances from \(v\) to a permutation index range, swaps allowed & dynamic set per bucket; a Fenwick over positions per centroid & \secref{c18-beyond} \\
\bottomrule
\end{tabular}
\end{center}

% ---------------- Notation ----------------
\subsection*{Notation Used Only Here}

The tree is \(T=(V,E)\), \(|V|=n\), unweighted unless a section says otherwise;
\(dist(u,v)\) is the number of edges (or the weighted length) on their path.  A
\emph{component} is a connected set of live (not-yet-deleted) nodes; \emph{size} is
its node count.  \(\mathit{cen}(S)\) is a centroid of component \(S\): a node whose
removal leaves every piece of size \(\le |S|/2\) (\(\defref{c09-centroid}\), Ch.~12).
\(\mathit{par}_c(u)\) is the \term{centroid tree} parent of \(u\): the centroid of
the component one level up that split \(u\) away (\(\secref{c18-tree}\)); \(\mathit{cd}(u)\)
is \(u\)'s depth in that tree, \(0\le\mathit{cd}\le\lceil\log_2 n\rceil\).  An
\emph{arm} of a centroid \(c\) is the set of nodes in one piece left after removing
\(c\), viewed by its distance from \(c\); a distance histogram \(H_c[d]\) counts
nodes at distance \(d\) from \(c\).  A node is \emph{deleted} / \emph{blocked} once
it is used as a centroid and is skipped in every deeper recursion.  The \emph{top
centroid of a path} is the first centroid (highest in the centroid tree) lying on
that path -- the node that "claims" it (\(\secref{c18-count}\)).  Every section runs
on \(T\) above; if a table here and a trace there disagree, this opener is
authoritative.

% =====================================================================
%  24.1  Recursion at the Balance -- building the centroid tree.
%  Exports: sec:c18-tree  def:c18-decomp  thm:c18-height
%           fig:c18-centtree  lst:c18-build  prob:c18-321c  lst:c18-321c
%  Solved CF 321C "Ciel the Commander" (2100), live-verified via the
%  problemset API (tags constructive algorithms / dfs / divide and conquer
%  / greedy / trees).  Official samples read live:
%    S1: n=4, star centred at 1 (edges 1-2,1-3,1-4) -> "A B B B".
%    S2: n=10, the path 1-2-...-10 -> "D C B A D C B D C D" (SPECIAL JUDGE:
%        any valid rank assignment is accepted, so the centroid-depth answer
%        need not equal the printed sample).
%  Hand-verified here: S1 centroid = 1 (remove 1 -> three singleton pieces),
%  cd(1)=0 -> 'A', cd(2)=cd(3)=cd(4)=1 -> 'B' => "A B B B", matching S1.
%  Back-cites Ch.12 thm:c09-centfind / lst:c09-cent for the find; this
%  section NEVER re-derives existence -- it begins at "delete it and recurse".
% =====================================================================
\section{Recursion at the Balance}\label{sec:c18-tree}

\begin{previouslybox}[Ch.~12 handed you one centroid and stopped on purpose]
Chapter~12 proved that every tree has a node \(c\) whose removal leaves no piece
bigger than \(\lfloor n/2\rfloor\), and that one DFS finds it
(\(\thmref{c09-centfind}\), \(\coderef{c09-cent}\)). It then stopped: \emph{one}
centroid, \emph{one} cut. This chapter opens the door it left ajar. A centroid is
not a fact to learn once -- it is the pivot of a divide-and-conquer. Delete \(c\),
then do the \emph{same thing} inside every leftover piece. Because each piece has
at most \(n/2\) nodes, the recursion is shallow; the tree of "who split me from
whom" it leaves behind -- the \term{centroid tree} (indexed once in the opener) --
is the object every later section stands on.
\end{previouslybox}

\begin{topicintro}
The whole decomposition is three lines of intent. Given a live component with a
chosen entry node: (1) find its centroid \(c\) with Chapter~12's find; (2)
\emph{mark \(c\) deleted} so no deeper call can step on it; (3) recurse
independently on each neighbour of \(c\) that is still live -- each such neighbour
roots one surviving piece. The centroid found at a deeper call declares the
centroid one level up as its \term{centroid tree} parent, writing
\(\mathit{par}_c\). The recursion's shape is the payoff: removing a centroid caps
every piece at \(|S|/2\), so a component's size at least halves every level --
the centroid tree has height \(\le\lceil\log_2 n\rceil\) (\(\thmref{c18-height}\))
and the build is \(\bigO{n\log n}\).

\textbf{Prerequisites.} Chapter~12 (\chref{c09}) -- the centroid definition
(\(\defref{c09-centroid}\)), existence (\(\thmref{c09-centexist}\)) and the
\(\bigO{n}\) find (\(\thmref{c09-centfind}\)), all used here as a black box. The
subtree-size DFS and the "upward piece \(n-\size(u)\)" discipline are also
Chapter~12's; nothing about distances appears in this section.

\textbf{After this section you can}
\begin{itemize}[nosep]
  \item run the cut-and-recurse loop on a small tree by hand and read off
        \(\mathit{par}_c\) and the centroid depth \(\mathit{cd}\); \hfill\lv{3}
  \item write the recursive build with a correct deleted-marker discipline -- a
        removed centroid is never re-entered -- in \(\bigO{n\log n}\); \hfill\lv{4}
  \item prove the height bound \(\mathit{cd}\le\lceil\log_2 n\rceil\) from the
        "every level halves" fact; \hfill\lv{4}
  \item solve CF~321C: the centroid depth \emph{is} the rank, so the build prints
        the answer directly. \hfill\lv{4}
\end{itemize}
\end{topicintro}

\begin{recognitionbox}[Pattern card -- "a property true along the whole recursion"]
\reckey{statements whose answer is a per-node number you can only get by
recursively re-centring -- \emph{rank every node so that two nodes of the same
rank are watched by a higher one} (CF~321C), \emph{sum over all nodes of their
centroid-tree depth}, or \emph{any quantity defined by "look at the balance point,
solve each side the same way, combine"}. The tell is a global, all-pairs or
whole-tree property with no natural fixed root.}
\recuse{
  \item the decomposition \emph{itself} is the output -- build \(\mathit{par}_c\), read \(\mathit{cd}\) (\(\secref{c18-tree}\));
  \item a path statistic a single rooted DP cannot fold -- claim paths at the centroid (\(\secref{c18-count}\));
  \item an all-component distance index reached from one touched node (\(\secref{c18-index}\)).
}
\recnotuse{
  \item just "print a centroid" -- that is Chapter~12's \(\thmref{c09-centfind}\), no recursion needed;
  \item a range update/query along one fixed path -- HLD (\chref{c16});
  \item a value that is a state per node with no whole-tree/all-pairs dimension -- tree DP (\chref{c10}).
}
One-line test: \emph{the question survives "cut at the balance point and solve both
sides identically" -- then recurse on the centroid, don't root once and hope.}
\end{recognitionbox}

% ---------------------------------------------------------------
\subsection{The decomposition, and the tree it leaves behind}
% ---------------------------------------------------------------

\begin{definition}{Centroid decomposition and the centroid tree}{c18-decomp}
The \term{centroid decomposition} of \(T\) is the recursion
\(\mathit{solve}(S)\): for a live component \(S\), let \(c=\mathit{cen}(S)\)
(\(\defref{c09-centroid}\)); set \(\mathit{par}_c(c)\) to the centroid that
created \(S\); mark \(c\) deleted; recurse on each connected component of
\(S-c\). The \term{centroid tree} is the tree on the same \(n\) nodes whose
parent of \(u\) is \(\mathit{par}_c(u)\); \(\mathit{cd}(u)\) is \(u\)'s depth in
it, and the original tree \(T\) and the centroid tree share their vertex set but
almost never their edges.
\end{definition}

\begin{insightbox}[Why deleting is the whole discipline]
The recursion never rebuilds the tree; it \emph{marks the centroid dead} and lets
the size-DFS treat every dead node as a wall. The moment a centroid is skipped
rather than deleted, a deeper call routes \emph{through} it into a piece it has
already claimed, a component gets processed twice, and the \(\le n/2\) guarantee
-- hence the \(\bigO{\log n}\) depth -- evaporates. The single line
\(\mathit{dead}[c]=\)~\code{true} is the difference between an
\(\bigO{n\log n}\) build and a re-processing \(O(n^2)\) one; the Common Mistakes
station at the chapter end makes you name the symptom.
\end{insightbox}

The running tree \(T\) of \(\figref{c18-tree}\) (opener) decomposes as follows;
the cuts are exactly the top centroid \(4\) and the three pieces it strands.

\begin{tracebox}[Decomposing \(T\) by hand -- component, its centroid, the pieces it leaves]
\textbf{Whole tree}, \(|S|=15\): centroid \(4\) (remove it leaves \(4,4,6\), all
\(\le 7\)); \(\mathit{par}_c(4)=0\), \(\mathit{cd}(4)=0\). Recurse on the pieces
\(A,B,C\).\\
\textbf{\(A=\{1,5,6,11\}\)}, \(|A|=4\): edges \(1\!-\!5,1\!-\!6,6\!-\!11\); centroid
\(1\) (pieces \(1,2\le 2\); ties with \(6\), keep \(1\)); \(\mathit{par}_c(1)=4\),
\(\mathit{cd}=1\). Pieces: \(\{5\}\to\mathit{par}_c(5)=1\); \(\{6,11\}\to\) centroid
\(6\), \(\mathit{par}_c(6)=1\), then \(\{11\}\to\mathit{par}_c(11)=6\).\\
\textbf{\(B=\{2,7,8,12\}\)}, \(|B|=4\): centroid \(2\) (ties with \(8\));
\(\mathit{par}_c(2)=4\). \(\{7\}\to7\); \(\{8,12\}\to\) centroid \(8\), \(\{12\}\to
\mathit{par}_c(12)=8\).\\
\textbf{\(C=\{3,9,10,13,14,15\}\)}, \(|C|=6\): centroid \(3\) (pieces \(2,3\le 3\);
ties with \(10\)); \(\mathit{par}_c(3)=4\). \(\{9,13\}\to\) centroid \(9\),
\(\mathit{par}_c(13)=9\); \(\{10,14,15\}\) is the path \(10\!-\!14\!-\!15\), centroid
\(14\), \(\mathit{par}_c(14)=3\), then \(\{10\},\{15\}\) with parent \(14\).\\
\textbf{Read off:} depths \(0\!:\!\{4\}\); \(1\!:\!\{1,2,3\}\);
\(2\!:\!\{5,6,7,8,9,14\}\); \(3\!:\!\{10,11,12,13,15\}\). Max depth \(3\le
\lceil\log_2 15\rceil=4\).
\end{tracebox}

\begin{figure}[H]
\centering
\begin{tikzpicture}[diagram, scale=0.78, transform shape,
   ct/.style={circle, draw=primaryblue!70, fill=primaryblue!8, line width=0.7pt,
              minimum size=6mm, font=\scriptsize\bfseries, inner sep=1pt},
   root/.style={circle, draw=warningred, fill=warningred!18, line width=1.1pt,
                minimum size=7mm, font=\footnotesize\bfseries, inner sep=1pt}]
  \node[root] (n4) at (0,3) {4};
  \node[ct] (n1) at (-6,2) {1}; \node[ct] (n2) at (0,2) {2}; \node[ct] (n3) at (6,2) {3};
  \node[ct] (n5) at (-7.2,1) {5}; \node[ct] (n6) at (-4.8,1) {6};
  \node[ct] (n7) at (-1.2,1) {7}; \node[ct] (n8) at (1.2,1) {8};
  \node[ct] (n9) at (4.8,1) {9}; \node[ct] (n14) at (7.6,1) {14};
  \node[ct] (n11) at (-4.8,0) {11}; \node[ct] (n12) at (1.2,0) {12};
  \node[ct] (n13) at (4.8,0) {13}; \node[ct] (n10) at (6.6,0) {10}; \node[ct] (n15) at (8.4,0) {15};
  \draw (n4)--(n1) (n4)--(n2) (n4)--(n3);
  \draw (n1)--(n5) (n1)--(n6) (n6)--(n11);
  \draw (n2)--(n7) (n2)--(n8) (n8)--(n12);
  \draw (n3)--(n9) (n3)--(n14) (n9)--(n13) (n14)--(n10) (n14)--(n15);
  \node[annot, font=\scriptsize, text=darkgray, anchor=east] at (-8.6,3) {\(\mathit{cd}=0\)};
  \node[annot, font=\scriptsize, text=darkgray, anchor=east] at (-8.6,1) {\(\mathit{cd}=2\)};
\end{tikzpicture}
\caption{The \term{centroid tree} of \(T\), built by \(\secref{c18-tree}\)'s
recursion. Its edges are \(\mathit{par}_c\) relations, not \(T\)'s edges: the top
centroid \(4\) (red) roots the tree, \(1,2,3\) are its three piece-centroids, and
the remaining nodes fill depths 2 and 3. The height of \(3\) obeys
\(\thmref{c18-height}\); CF~321C prints exactly this tree, reading each rank as
\(\text{'A'}+\mathit{cd}\).}
\label{fig:c18-centtree}
\end{figure}

\begin{lstlisting}[caption={The recursive centroid-tree build over a fixed tree. The centroid finder is Ch.12's walk-to-heavy (thm:c09-centfind) used as a black box; the dead[] marker keeps pieces out of already-claimed nodes. ASCII only}, label={lst:c18-build}]
// Global: adj[], n. Outputs par_c[1..n] (centroid parent, 0 = top) and cd[1..n].
const int MAXN = 200005;
vector<int> adj[MAXN];
bool dead[MAXN];
int sz[MAXN], par_c[MAXN], cd[MAXN];

int dfs_size(int u, int p){                    // sizes inside this LIVE component
    sz[u] = 1;
    for(int v : adj[u])
        if(v != p && !dead[v]) sz[u] += dfs_size(v, u);
    return sz[u];
}
int find_centroid(int u, int p, int tot){      // Ch.12 thm:c09-centfind, walk-to-heavy
    bool moved = true;
    while(moved){
        moved = false;
        for(int v : adj[u])
            if(v != p && !dead[v] && sz[v] > tot/2){ p = u; u = v; moved = true; break; }
    }
    return u;                                  // here every piece is <= tot/2
}
void decompose(int entry, int parent_centroid){
    int tot = dfs_size(entry, -1);             // entry roots one live piece
    int c   = find_centroid(entry, -1, tot);   // Ch.12's find
    par_c[c] = parent_centroid;                // c's parent in the centroid tree
    cd[c]    = (parent_centroid == 0) ? 0 : cd[parent_centroid] + 1;
    dead[c]  = true;                           // DELETED: a wall for every deeper call
    for(int v : adj[c])                        // each live neighbour roots a surviving piece
        if(!dead[v]) decompose(v, c);
}
// Build the whole centroid tree with:  decompose(1, 0);
\end{lstlisting}

\begin{theorem}{The centroid tree has logarithmic height}{c18-height}
For any tree on \(n\ge 1\) nodes, the centroid tree produced by
\(\secref{c18-tree}\) has height \(\le\lceil\log_2 n\rceil\): every node satisfies
\(\mathit{cd}(u)\le\lceil\log_2 n\rceil\).
\end{theorem}
\begin{proofbox}
Follow the chain of components containing a fixed node \(u\): the whole tree has
\(n\) nodes; after the top centroid is deleted, \(u\)'s component has at most
\(\lfloor n/2\rfloor\) nodes because a centroid caps every piece at half
(\(\defref{c09-centroid}\)); each further level re-centres a component and again
caps \(u\)'s piece at half its parent's size. So the component holding \(u\) has
size \(\le n,\ \lfloor n/2\rfloor,\ \lfloor n/4\rfloor,\dots\) reaching \(1\) after
at most \(\lceil\log_2 n\rceil\) halvings -- and each halving corresponds to
exactly one centroid on \(u\)'s centroid-ancestor chain. Hence
\(\mathit{cd}(u)\le\lceil\log_2 n\rceil\). Two-centroid ties (\(\thmref{c09-centexist}\))
do not affect the bound: whichever of the pair is chosen, both cap every piece at
\(\le n/2\), which is all the induction uses. \(\Box\)
\end{proofbox}

\begin{complexitybox}[The build]
\textbf{Time} \(\bigO{n\log n}\): a node lies in exactly one live component per
centroid level, so it is visited by \(\mathit{dfs\_size}\)/\(\mathit{find\_centroid}\)
once per level, and there are \(\le\lceil\log_2 n\rceil\) levels
(\(\thmref{c18-height}\)). \textbf{Space} \(\bigO{n}\) for \(\mathit{par}_c\),
\(\mathit{cd}\), \(\mathit{sz}\), \(\mathit{dead}\). On a chain the naive
"recolor and recurse on a rooted subtree" would be \(\bigO{n^2}\); the
\(\le n/2\) cut is precisely what collapses it to \(\bigO{n\log n}\) -- the same
halving \(\thmref{c09-centexist}\)'s proof already promised.
\end{complexitybox}

% ---------------------------------------------------------------
\subsection{Problem Lab: CF 321C -- rank every node by its centroid depth}
% ---------------------------------------------------------------

\solvedproblem{c18-321c}{Ciel the Commander}{CF 321C}
  {https://codeforces.com/problemset/problem/321/C}{CF 2100 L4}
  {centroid decomposition; the answer is 'A' + centroid depth; special judge}

\begin{problemstatement}[Ciel the Commander (statement and samples verified against the live problem)]
A tree has \(n\) cities (\(2\le n\le 10^5\)). Assign each city an officer rank, a
letter \(\text{'A}\)–\(\text{'Z}\) ('A' topmost). Rule: for any two distinct cities
\(x,y\) of the \emph{same} rank, the simple path between them must contain a city
of \emph{strictly higher} rank. Output the \(n\) ranks (space-separated), or
\texttt{Impossible!} if none exists. \textbf{Special judge: any valid assignment
is accepted}, so your string may differ from the sample.

\textbf{Examples (official).}
\begin{verbatim}
Input:              Output:
4                   A B B B
1 2
1 3
1 4

Input:              Output:
10                  D C B A D C B D C D
1 2 ... 9 10        (path on 10 nodes; one of many valid answers)
\end{verbatim}
\end{problemstatement}

\begin{attemptgate}[Attempt it yourself first -- \lv{4}, give yourself 25 minutes]
Guess that the topmost rank 'A' goes to a centroid, then ask what rank 'B' can
occupy. \peek{which nodes may share a rank without violating the monitoring rule?}
\nudge{two same-rank nodes are fine if their path holds a higher rank; a centroid's
removal is exactly the thing that separates pieces from each other.}
\fullhint{Rank a node by its centroid-tree depth: \(\text{'A'}+\mathit{cd}(u)\)
(\(\coderef{c18-build}\)). If \(x,y\) have the same depth \(d\), their lowest common
ancestor in the \emph{centroid} tree has depth \(<d\), and that LCA centroid lies on
the original-tree path \(x\text{--}y\) -- the higher-rank monitor the rule wants.
For \(n\le 10^5\), \(\lceil\log_2 n\rceil\le 17<26\), so 'A'..'R' always suffice and
\texttt{Impossible!} never fires on valid input. Build the decomposition and print
\(\text{'A'}+\mathit{cd}[i]\).}
\end{attemptgate}

\solstep{Understanding}
The rank is a function you must place on vertices so that equal values are always
"bridged" by a larger one along their path. That is not a distance problem and not
a fixed-root DP; it is a statement about a \emph{recursive separation} of the tree
into halves -- which is the centroid decomposition. The insight is that
\(\mathit{cd}\) (not \(\size\), not depth in \(T\)) is the right coordinate, and
then the problem is \emph{literally} "run \(\secref{c18-tree}\)'s build and print
the depth."

\solstep{Key Observations}
\begin{observation}
Setting \(\mathrm{rank}(u)=\text{'A'}+\mathit{cd}(u)\) is valid. For same-depth
\(x\ne y\), let \(w\) be their centroid-tree LCA; \(w\) is a proper ancestor of
both so \(\mathit{cd}(w)<\mathit{cd}(x)\), i.e.\ higher rank; and \(w\) is the first
centroid whose cut separated \(x\) and \(y\), which forces \(w\) onto the
\(T\)-path between them.
\end{observation}
\begin{observation}
The bound is always satisfiable here: \(\thmref{c18-height}\) gives
\(\max\mathit{cd}\le\lceil\log_2 10^5\rceil=17\le 25\), so no rank exceeds 'Z' and
the \texttt{Impossible!} branch is unreachable for \(n\le 10^5\).
\end{observation}

\solstep{The Approach Ladder}
\begin{center}\renewcommand{\arraystretch}{1.25}\setlength{\tabcolsep}{3pt}
\begin{tabular}{@{}c p{0.26\linewidth} p{0.30\linewidth} l c@{}}
\toprule
\textbf{\#} & \textbf{Approach} & \textbf{Idea} & \textbf{Time} & \textbf{Verdict}\\
\midrule
1 & \lv{2}\ fixed-root depth & rank by depth in \(T\) & \(\bigO{n}\) & \textcolor{warningred}{\ding{55}\ \textbf{WA}} \\
2 & \lv{3}\ greedy leaf-peel & colour leaves 'Z' upward & \(\bigO{n^2}\) & \tle \\
3 & \lv{4}\ \textbf{centroid depth} & rank\(=\text{'A'}+\mathit{cd}\) (\(\coderef{c18-build}\)) & \(\bigO{n\log n}\) & \ac \\
\bottomrule
\end{tabular}\end{center}
Rung~1 is the seductive trap: depth in \(T\) under a bad root can be \(O(n)\) (a
chain), pushing ranks past 'Z', and it puts two same-rank leaves on one side of
every internal node -- no higher-rank monitor between them. Rung~3's centroid depth
is \(O(\log n)\) (\(\thmref{c18-height}\)) and separates by construction.

\solstep{Solution}
Run \(\mathit{decompose}(1,0)\) from \(\coderef{c18-build}\) to fill \(\mathit{cd}\);
print \(\text{'A'}+\mathit{cd}[i]\) for \(i=1..n\).

\solstep{Implementation}
\begin{lstlisting}[caption={CF 321C driver: build the centroid tree once, then print 'A' + centroid depth per node. Reuses decompose() from lst:c18-build verbatim. ASCII only}, label={lst:c18-321c}]
#include <bits/stdc++.h>
using namespace std;
// ... paste adj[], dead[], sz[], par_c[], cd[], dfs_size, find_centroid,
//     and decompose() exactly as in lst:c18-build ...
int main(){
    int n;  scanf("%d", &n);
    for(int i=0;i<n-1;i++){
        int a,b; scanf("%d %d",&a,&b);
        adj[a].push_back(b); adj[b].push_back(a);
    }
    decompose(1, 0);                       // fills cd[1..n]
    for(int i=1;i<=n;i++)
        printf("%c%c", 'A'+cd[i], i==n ? '\n' : ' ');
    // cd[i] <= 17 <= 25, so a rank is always a letter: never "Impossible!" here.
}
\end{lstlisting}

\solstep{Trace}
\begin{tracebox}[Sample 1: the star \(n=4\), edges \(1\!-\!2,1\!-\!3,1\!-\!4\)]
\(\mathit{decompose}(1,0)\): centroid of the whole star is \(1\) (remove it leaves
\(\{2\},\{3\},\{4\}\), each size \(1\le 2\)); \(\mathit{cd}(1)=0\). Recurse on each
leaf component: centroid of \(\{2\}\) is \(2\), \(\mathit{cd}=1\); likewise \(3,4\).
So \(\mathit{cd}=(0,1,1,1)\) for nodes \((1,2,3,4)\) and the printout is
\(\boxed{\texttt{A B B B}}\) -- matching the official output. On \(T\) of
\(\figref{c18-tree}\) the same code yields
\(1\!:\!\text{B},2\!:\!\text{B},3\!:\!\text{B},4\!:\!\text{A},5\!:\!\text{C},6\!:\!\text{C},
7\!:\!\text{C},8\!:\!\text{C},9\!:\!\text{C},10\!:\!\text{D},11\!:\!\text{D},12\!:\!\text{D},
13\!:\!\text{D},14\!:\!\text{C},15\!:\!\text{D}\) -- the depth layers read straight
off \(\figref{c18-centtree}\).
\end{tracebox}

\begin{warningbox}[Reusing sizes across the recursion -- the signature bug]
\(\mathit{sz}[\cdot]\) is meaningful \emph{only} inside the current component,
rooted at the current \(\mathit{entry}\), \emph{after} the ancestors are dead. If you
compute \(\mathit{sz}\) once at the top and reuse it in deeper calls (skipping the
per-call \(\mathit{dfs\_size}\)), the centroid test is wrong: a stale \(\mathit{sz}\)
counts nodes that belong to an already-deleted sibling. Symptom on a chain: the build
either picks a non-centroid (linear depth, ranks overflow 'Z', WA) or loops. Re-run
\(\mathit{dfs\_size}\) at every \(\mathit{decompose}\) call -- that is the whole
\(\bigO{n\log n}\) budget.
\end{warningbox}

\begin{hardeningbox}
\begin{itemize}[nosep,leftmargin=1.4em]
  \item \textbf{Deletion before recursion.} Set \(\mathit{dead}[c]=\)~\code{true}
        \emph{before} iterating \(c\)'s neighbours, and guard every edge walk with
        \(\mathit{!dead}[\cdot]\); otherwise the recursion steps back onto the pivot.
  \item \textbf{Stack budget.} The centroid recursion is \(O(\log n)\) deep, but the
        inner \(\mathit{dfs\_size}\) is as deep as the \emph{component} (up to \(n\) on
        a chain). For \(n=10^5\) use an explicit stack (Ch.~12's \(\coderef{c09-iter}\))
        or raise the recursion limit -- \(\secref{c18-tree}\)'s listings keep the
        recursive form for legibility.
  \item \textbf{Two-centroid ties.} When a component admits a pair
        (\(\thmref{c09-centexist}\)), any choice is legal -- the height proof uses only
        the \(\le n/2\) cap, so \(\mathit{cd}\) stays bounded and CF~321C stays valid.
  \item \textbf{Special judge.} Do not compare your string to the sample literally;
        verify the rule, or trust \(\thmref{c18-height}\) and the LCA-monitor
        observation.
\end{itemize}
\end{hardeningbox}

\begin{recallbox}[24.1 in one breath]
\emph{Find the centroid (Ch.~12, \(\thmref{c09-centfind}\)), delete it, recurse on
each surviving piece, and record \(\mathit{par}_c\) -- that is the centroid tree
(\(\defref{c18-decomp}\), \(\coderef{c18-build}\)). Every level at least halves a
component, so height \(\le\lceil\log_2 n\rceil\) (\(\thmref{c18-height}\)) and the
build is \(\bigO{n\log n}\). CF~321C is this build seen from the output side: rank
\(=\text{'A'}+\mathit{cd}\), because a same-rank pair's centroid-LCA is a higher-rank
monitor on their path.}
\selftest{
  \item Why is \(\mathit{par}_c\) almost never equal to \(T\)'s parent array? \emph{(\(\mathit{par}_c(u)\) is the centroid that split \(u\)'s component, not its tree parent (\(\figref{c18-centtree}\)).)}
  \item State the invariant that bounds the centroid-tree height. \emph{(Every level caps the component at \(\le\) half (\(\thmref{c18-height}\)).)}
  \item In CF~321C, if two 'B' cities exist, what guarantees a higher rank on their path? \emph{Their centroid-tree LCA has depth \(0\) (rank 'A') and lies on their path (Key Observation~1).)}
}
\end{recallbox}

% =====================================================================
%  24.2  Counting paths by distance -- claim each path at its top centroid.
%  Exports: sec:c18-count  def:c18-claim  lst:c18-count
%           prob:c18-715c  lst:c18-715c
%  Solved CF 715C "Digit Tree" (2700), live-verified via the problemset API
%  (tags dfs and similar / divide and conquer / dsu / trees).  Official
%  samples read live:
%    S1: n=6, M=7 -> 7  (note enumerates (0,4),(1,2),(1,5),(3,2),(2,5),(5,2),
%        (3,5) forming 14,21,217,91,7,7,917, all multiples of 7).
%    S2: n=5, M=11 -> 8 (six pairs give 33, two give 3333).
%  Discharges the forward-cites from Ch.22/Ch.23/Ch.16 ("path COUNTING /
%  distance stats -> centroid").  CF 161D (pairs at distance exactly K) is
%  OWNED by Ch.13 as the O(nK) subtree DP -- back-cited by name/label only,
%  never re-linked.  Hand-verified here: the distance-exactly-2 count on T is
%  17 (= sum_z C(deg(z),2)); the top centroid 4 claims 9 of them (6 with the
%  centroid itself + 3 pairs among its three neighbours); recursion claims 8.
% =====================================================================
\section{Counting Paths by Distance}\label{sec:c18-count}

\begin{previouslybox}[Ch.~13 counted distance-\(K\) pairs in \(\bigO{nK}\); here \(K\) is huge]
Chapter~13 harvested, for every node, a histogram of downward distances and merged
children, answering ``how many pairs sit at distance exactly \(K\)'' in
\(\bigO{nK}\) time (\(\probref{c10-harvest-solved}\), CF~161D). That is excellent
when \(K\le 500\) and hopeless when \(n=10^5\) and \(K\) can reach \(n\): the DP
budget is a product, not a bound. The centroid claim removes the product. A path
has two endpoints and, in a tree, a unique route; charge the path to the highest
centroid that lies on it. At that centroid the path splits into two independent
arms, so its length is a \emph{sum of two depths from one node} -- count those
sums with a histogram and you never touch a \(K\)-sized table.
\end{previouslybox}

\begin{topicintro}
Fix a centroid \(c\) of a live component. Every node \(x\) has a distance
\(d_c(x)\) from \(c\). A pair \(\{u,v\}\) whose path passes through \(c\) has
length \(d_c(u)+d_c(v)\), and \(c\) is precisely the first centroid on that path
when \(u\) and \(v\) live in \emph{different} child-components of \(c\). So at
\(c\): walk its child arms one at a time, and for each new node query a running
histogram for partners already seen (in earlier arms, or \(c\) itself), then merge
the new arm into the histogram. Because a partner is always from an
\emph{earlier} arm, same-child pairs are never counted at \(c\) -- they wait for a
deeper centroid. The recursion visits each node \(O(\log n)\) times, so a linear
per-level merge gives \(\bigO{n\log n}\) total, and the same skeleton answers
\(\le K\), \(=K\), or (with a different per-node value) far more exotic path
statistics.

\textbf{Prerequisites.} \(\secref{c18-tree}\) -- the decomposition, the deleted
marker, and \(\mathit{find\_centroid}\). Chapter~16 (\chref{c11}) -- the identity
\(\dist(u,v)=\depth(u)+\depth(v)-2\depth(\lca)\), which here collapses to
\(d_c(u)+d_c(v)\) because the path goes through \(c\).

\textbf{After this section you can}
\begin{itemize}[nosep]
  \item count unordered pairs at distance \(=K\) in \(\bigO{n\log n}\) by the
        query-then-merge histogram, and explain why same-arm pairs are excluded; \hfill\lv{4}
  \item state the top-centroid claim: each path is counted once, at its highest
        centroid (\(\defref{c18-claim}\)); \hfill\lv{4}
  \item recognise a path statistic that is a \emph{sum of two per-arm values} and
        route it through the same skeleton; \hfill\lv{5}
  \item solve CF~715C: the ``value'' is a decimal hash, and coprimality with 10
        turns the concatenation into a matched-key count. \hfill\lv{5}
\end{itemize}
\end{topicintro}

% ---------------------------------------------------------------
\subsection{The top-centroid claim and the distance histogram}
% ---------------------------------------------------------------

\begin{definition}{The top centroid of a path}{c18-claim}
For distinct nodes \(u,v\), the \emph{top centroid} of the path \(u\text{--}v\) is
the centroid lying on that path with the smallest \(\mathit{cd}\) (the highest in
the centroid tree). It exists because the first centroid ever selected inside a
component containing both \(u\) and \(v\) separates them into different child
components, and therefore sits on their path. We say that centroid \emph{claims}
the pair.
\end{definition}

\begin{insightbox}[Why the correction is free]
The textbook way is ``count every pair with \(d_c\) sum \(=K\), then subtract the
pairs that stayed inside one child component.'' The query-then-merge loop
(\(\coderef{c18-count}\)) does the same thing in one pass: a node only ever meets
partners from \emph{earlier} arms and from the centroid, so a same-arm pair is
never formed at \(c\) and there is nothing to subtract after the fact. The
subtraction and the incremental merge are the same idea; the incremental form is
harder to get wrong, so the listing uses it -- and \(\secref{c18-beyond}\)'s map
method reverts to explicit subtraction when the value is not a small integer
array.
\end{insightbox}

\begin{example}{Counting distance-2 pairs on \(T\)}
On the running tree, \(\sum_z\binom{\deg(z)}{2}=17\) pairs sit at distance exactly
\(2\) (any two neighbours of a node are \(2\) apart in a tree, uniquely). Run the
claim at the top centroid \(4\): distances from \(4\) are
\(\{1,2,3\}{=}1\), \(\{5,6,7,8,9,10\}{=}2\), \(\{11,12,13,14\}{=}3\), \(15{=}4\).
A pair through \(4\) with depth-sum \(2\) is either (centroid, a depth-2 node) or
(two depth-1 nodes): \(1\cdot 6 + \binom{3}{2}=6+3=9\). Those \(9\) are
\((4,5),(4,6),(4,7),(4,8),(4,9),(4,10),(1,2),(1,3),(2,3)\). Every other distance-2
pair lives inside one of \(4\)'s arms (\((5,6)\) under~\(1\), \((9,10)\) under~\(3\),
\((10,15)\) under~\(14\), and so on) and is claimed deeper. \(9+8=17\): the pieces
sum to the ground truth. \(\checkmark\)
\end{example}

\begin{lstlisting}[caption={Count unordered pairs at distance exactly K over a fixed tree by centroid decomposition: at each centroid, query-then-merge a depth histogram across the child arms so same-child pairs are never formed here. Reuses the deleted marker and find-centroid from lst:c18-build. ASCII only}, label={lst:c18-count}]
// K given; H[d] = how many seen nodes sit at distance d from the current centroid
const int MAXN = 200005;
vector<int> adj[MAXN];
bool dead[MAXN]; int sz[MAXN];
int K;
long long ans = 0;
int touched[MAXN], tcnt;            // depths written, so we can clear H

int H[MAXN];                        // zero-initialised

void collect_depth(int u, int p, int d, vector<int>& arm){
    arm.push_back(d);
    if(d <= K) touched[tcnt++] = d;
    for(int v : adj[u]) if(v != p && !dead[v]) collect_depth(v, u, d+1, arm);
}
void count_through(int c){
    H[0] = 1;                        // the centroid itself, at distance 0
    for(int v : adj[c]){
        if(dead[v]) continue;
        vector<int> arm;
        collect_depth(v, c, 1, arm);              // depths of this whole arm
        for(int d : arm)
            if(K - d >= 0) ans += H[K-d];         // pair with earlier arms + c
        for(int d : arm) H[d]++;                  // only THEN merge the arm
    }
    H[0] = 0;
    for(int i = 0; i < tcnt; i++) H[touched[i]] = 0;   // clear the histogram
    tcnt = 0;
}
// full recursion (same shape as lst:c18-build):
void solve_comp(int entry){
    int tot = dfs_size(entry, -1);
    int c = find_centroid(entry, -1, tot);
    count_through(c);
    dead[c] = true;
    for(int v : adj[c]) if(!dead[v]) solve_comp(v);
}
// call  solve_comp(1);  then  ans  = #pairs at distance K.
\end{lstlisting}

\begin{complexitybox}[The distance count]
Each centroid does work linear in its component (one depth-collection and one
query per node), and every node lies in \(O(\log n)\) components
(\(\thmref{c18-height}\)), so total time is \(\bigO{n\log n}\) with
\(\bigO{n}\) scratch -- independent of \(K\). That independence is the whole
point versus the \(\bigO{nK}\) subtree DP of Chapter~13: when \(K\) is a
constant \(\le 500\) the DP wins on small \(n\); when \(K\) is \(\Theta(n)\) the
centroid count is the only one that finishes.
\end{complexitybox}

% ---------------------------------------------------------------
\subsection{Problem Lab: CF 715C -- the number read along a path, mod M}
% ---------------------------------------------------------------

\solvedproblem{c18-715c}{Digit Tree}{CF 715C}
  {https://codeforces.com/problemset/problem/715/C}{CF 2700 L5}
  {centroid decomposition; forward/backward decimal hashes; gcd(M,10)=1 gives an inverse of 10; per-arm subtraction over hash maps}

\begin{problemstatement}[Digit Tree (statement and samples verified against the live problem)]
A tree on \(n\) vertices \(0,\dots,n-1\) (\(2\le n\le 10^5\)) has a single nonzero
digit \(1\)–\(9\) on each edge. A positive integer \(M\) (\(1\le M\le 10^9\)) is
given with \(\gcd(M,10)=1\). An \textbf{ordered} pair of distinct vertices
\((u,v)\) is \emph{interesting} if the decimal number obtained by writing the edge
digits along the shortest path from \(u\) to \(v\), \emph{in the order they are
encountered}, is divisible by \(M\). Count the interesting pairs.

\textbf{Examples (official).}
\begin{verbatim}
Input:              Output:
6 7                 7
0 1 2
4 2 4
2 0 1
3 0 9
2 5 7

Input:              Output:
5 11                8
1 2 3
2 0 3
3 0 3
4 3 3
\end{verbatim}
Sample 1's note lists the seven pairs
\((0,4),(1,2),(1,5),(3,2),(2,5),(5,2),(3,5)\) giving \(14,21,217,91,7,7,917\) --
all multiples of \(7\); \((2,5)\) and \((5,2)\) count separately (ordered).
\end{problemstatement}

\begin{attemptgate}[Attempt it yourself first -- \lv{5}, give yourself 45 minutes]
Write the path number as ``\(u\) to \(c\), then \(c\) to \(v\)'' for the top
centroid \(c\) and try to separate the two halves. \peek{the digits from \(u\) up
to \(c\) are the reverse of those from \(c\) down to \(u\); how does appending a
block of digits combine two numbers?}
\nudge{the value read along \(u\to v\) is \(\mathit{back}[u]\cdot 10^{\mathit{dep}[v]} + \mathit{for}[v]\), where \(\mathit{back}[u]\) is the number read \(u\to c\) and \(\mathit{for}[v]\) the number read \(c\to v\). You need that \(\equiv 0 \pmod M\).}
\fullhint{Divide by \(10^{\mathit{dep}[v]}\), legal since \(\gcd(M,10)=1\): the condition
becomes back[u] \(\equiv -\,\text{for}[v]\cdot 10^{-\mathit{dep}[v]}\pmod M\). Store each
node twice -- as a start by key back[u], as an end by key
\(-\,\text{for}[v]\cdot\) inv10\([\mathit{dep}[v]]\) -- and count matched keys with a hash
map. Count all ordered matches over the component, then \emph{subtract} the
matches inside each single child arm (those paths do not pass through \(c\)); the
centroid, with for \(=\) back \(= 0\), pairs correctly with every node whose
one-sided value is \(0\pmod M\). Recurse; total \(\bigO{n\log n}\) expected.}
\end{attemptgate}

\solstep{Understanding}
The unit of work is still a path claimed at its top centroid, but the per-node
value is no longer an integer distance: it is a residue that depends on the
\emph{digits} and their \emph{order}, so a path through \(c\) is a concatenation of
an upward part and a downward part. Concatenation is `\(A\cdot 10^{\mathit{len}(B)}+B\)',
and coprimality with 10 makes \(10^{\mathit{len}(B)}\) invertible, which decouples the two
sides into an equality of single-node keys. That decoupling is exactly what lets
the \(\secref{c18-count}\) skeleton -- count by a matched key, subtract one-arm
pairs -- survive the move from distances to decimal numbers.

\solstep{Key Observations}
\begin{observation}
For a path \(u\to v\) through centroid \(c\), let \(\mathit{back}[u]\) be the value
of the digits read \(u\to c\) and \(\mathit{for}[v]\) the value read \(c\to v\),
with \(\mathit{dep}[\cdot]\) the edge counts. Then
\(\mathrm{val}(u,v)=\mathit{back}[u]\cdot 10^{\mathit{dep}[v]}+\mathit{for}[v]\).
\end{observation}
\begin{observation}
Since \(\gcd(M,10)=1\), \(10^{\mathit{dep}[v]}\) has an inverse mod \(M\), so
\(\mathrm{val}(u,v)\equiv 0 \iff
\mathit{back}[u]\equiv -\,\mathit{for}[v]\cdot 10^{-\mathit{dep}[v]}\pmod M\): a
pure key-equality, with \(u\)'s key independent of \(v\).
\end{observation}
\begin{observation}
A pair is claimed by \(c\) only if its two endpoints are in different child arms
(or one endpoint is \(c\)). Matches with both endpoints in one arm are subtracted;
they are counted at a deeper centroid, exactly as in \(\secref{c18-count}\).
\end{observation}

\solstep{The Approach Ladder}
\begin{center}\renewcommand{\arraystretch}{1.25}\setlength{\tabcolsep}{3pt}
\begin{tabular}{@{}c p{0.28\linewidth} p{0.30\linewidth} l c@{}}
\toprule
\textbf{\#} & \textbf{Approach} & \textbf{Idea} & \textbf{Time} & \textbf{Verdict}\\
\midrule
1 & \lv{2}\ all ordered pairs & DFS a number per \((u,v)\) & \(\bigO{n^2}\) & \tle \\
2 & \lv{4}\ rooted DP on digits & fix a root, fold & -- & \textcolor{warningred}{\ding{55}\ \textbf{WA}} \\
3 & \lv{5}\ \textbf{centroid + hash keys} & back / for keys, subtract arms (\(\coderef{c18-715c}\)) & \(\bigO{n\log n}\) & \ac \\
\bottomrule
\end{tabular}\end{center}
Rung~2 fails because a single root cannot see a path that climbs to the top
centroid and back down with the correct digit order in both directions; the
reverse-direction hash \(\mathit{back}[\cdot]\) is not derivable from
\(\mathit{for}[\cdot]\) without the \(10^{-\mathit{dep}}\) normalisation the centroid
framing supplies.

\solstep{Solution}
Precompute \(10^k\bmod M\) and its inverses. Build the centroid decomposition
(\(\coderef{c18-build}\)); at each centroid \(c\) collect, for every node of the
component, the pair of keys \((\mathit{back},\mathit{key})\); add
\(\mathit{pairs}(\text{all})\) and subtract \(\mathit{pairs}(\text{each arm})\),
where \(\mathit{pairs}\) counts ordered distinct matches of the start-key against
the end-key.

\solstep{Implementation}
\begin{lstlisting}[caption={CF 715C: centroid decomposition with a forward/backward decimal hash per node and a per-arm subtraction over hash maps; gcd(M,10)=1 supplies inv(10, M). M=1 is short-circuited. ASCII only}, label={lst:c18-715c}]
#include <bits/stdc++.h>
using namespace std;
typedef long long ll;
const int MAXN = 100005;
int n; ll M;
vector<pair<int,int>> adj[MAXN];        // (neighbour, digit 1..9)
bool dead[MAXN]; int sz[MAXN];
ll p10[MAXN], ip10[MAXN];               // 10^k mod M, and its inverse
ll ans = 0;

ll egcd(ll a, ll b, ll &x, ll &y){
    if(!b){ x=1; y=0; return a; }
    ll g, x1, y1; g = egcd(b, a % b, x1, y1);
    x = y1; y = x1 - (a / b) * y1; return g;
}
ll invmod(ll a){ ll x, y; egcd(a, M, x, y); return (x % M + M) % M; }
ll norm(ll x){ x %= M; if(x < 0) x += M; return x; }

int dfs_size(int u, int p){
    sz[u] = 1;
    for(auto &e : adj[u]) if(e.first != p && !dead[e.first])
        sz[u] += dfs_size(e.first, u);
    return sz[u];
}
int walk_cent(int u, int p, int tot){
    bool moved = true;
    while(moved){
        moved = false;
        for(auto &e : adj[u]){
            int v = e.first;
            if(v != p && !dead[v] && sz[v] > tot/2){ p = u; u = v; moved = true; break; }
        }
    }
    return u;
}

struct Info { ll b, k; };               // b = back[u]; k = -for*inv10[dep]
void collect(int u, int p, int dep, ll fwd, ll bwd, vector<Info>& out){
    out.push_back({ norm(bwd), norm(-fwd * ip10[dep]) });
    for(auto &e : adj[u]){
        int v = e.first; ll w = e.second;
        if(v == p || dead[v]) continue;
        collect(v, u, dep+1, norm(fwd*10 + w), norm(w * p10[dep] + bwd), out);
    }
}
ll pairs(const vector<Info>& a){         // ordered distinct (u,v), back[u]==key[v]
    unordered_map<ll,ll> fB, fK;
    for(auto &x : a){ fB[x.b]++; fK[x.k]++; }
    ll cnt = 0;
    for(auto &kv : fB){ auto it = fK.find(kv.first); if(it != fK.end()) cnt += kv.second * it->second; }
    for(auto &x : a) if(x.b == x.k) cnt--;        // drop u == v
    return cnt;
}
void process(int c){
    vector<Info> all; all.push_back({ 0, 0 });     // the centroid itself
    vector<vector<Info>> arms;
    for(auto &e : adj[c]){
        int v = e.first; ll w = e.second;
        if(dead[v]) continue;
        vector<Info> arm;
        collect(v, c, 1, norm(w), norm(w), arm);
        all.insert(all.end(), arm.begin(), arm.end());
        arms.push_back(arm);
    }
    ans += pairs(all);
    for(auto &a : arms) ans -= pairs(a);           // subtract same-child pairs
}
void decompose(int entry){
    int tot = dfs_size(entry, -1);
    int c = walk_cent(entry, -1, tot);
    process(c);
    dead[c] = true;
    for(auto &e : adj[c]) if(!dead[e.first]) decompose(e.first);
}
int main(){
    scanf("%d %lld", &n, &M);
    if(M == 1){ printf("%lld\n", (ll)n * (n - 1)); return 0; }   // every pair works
    for(int i = 0; i < n-1; i++){
        int u, v, w; scanf("%d %d %d", &u, &v, &w);
        adj[u].push_back({v, w}); adj[v].push_back({u, w});
    }
    ll inv10 = invmod(10 % M);
    p10[0] = 1 % M; ip10[0] = 1 % M;
    for(int i = 1; i <= n; i++){ p10[i] = p10[i-1]*10 % M; ip10[i] = ip10[i-1]*inv10 % M; }
    decompose(0);
    printf("%lld\n", ans);
}
\end{lstlisting}

\solstep{Trace}
\begin{tracebox}[Sample 2: \(n=5\), \(M=11\), every edge digit \(=3\)]
A run of \(k\) threes is the number \(3,33,333,\dots\); mod \(11\), \(33\equiv 0\)
and \(3333\equiv 0\) while \(3\not\equiv 0\) and \(333\equiv 3\), so a path reads to
\(0\pmod{11}\) exactly when it uses an \emph{even} number of edges. The four edges
\(1\!-\!2,2\!-\!0,3\!-\!0,4\!-\!3\) form the path \(1\!-\!2\!-\!0\!-\!3\!-\!4\)
(order \(1,2,0,3,4\)); a pair is interesting iff their positions differ by an even
number. Position gaps of \(2\) give \(\{1,0\},\{2,3\},\{0,4\}\) and a gap of \(4\)
gives \(\{1,4\}\): four unordered pairs, eight ordered ones --
\((4,0),(0,4),(3,2),(2,3),(0,1),(1,0),(4,1),(1,4)\), the official eight, six reading
\(33\) and two reading \(3333\). The centroid decomposition returns \(\boxed{8}\)
whatever centroid is chosen first, since each pair is claimed at exactly one top
centroid.
\end{tracebox}

\begin{warningbox}[Ordered pairs, and the modulo must be normalised]
This problem counts \emph{ordered} pairs: \((2,5)\) and \((5,2)\) are different, and
they can even carry different numbers (\(7\) vs \(7\) here, but not in general), so
\(\mathit{pairs}\) matches the start-key against the end-key over all nodes in both
directions -- never halve the result as you would for an unordered distance count.
And every key must be pushed into \([0,M)\) with \(\mathit{norm}\): the raw
\(-\,\text{for}\cdot\mathit{inv10}\) is negative, and comparing a negative residue
against a stored positive one silently under-counts -- a classic \textbf{WA} that
passes the small samples and dies on the large ones.
\end{warningbox}

\begin{hardeningbox}
\begin{itemize}[nosep,leftmargin=1.4em]
  \item \textbf{\(\gcd(M,10)=1\) is load-bearing.} The \(10^{-\mathit{dep}}\) normalisation
        needs an inverse of 10; if the input had \(2\mid M\) or \(5\mid M\) the key
        split is invalid. Here it is promised. Guard \(M=1\) separately (every
        ordered pair works: \(n(n-1)\)).
  \item \textbf{Overflow.} \(M\le 10^9\) so products like \(\mathit{fwd}\cdot
        \mathit{inv10}\) reach \(10^{18}\) -- use \code{long long} and take \(\bmod M\)
        immediately; \(\mathit{ans}\) can be \(\sim n^2\), also \code{long long}.
  \item \textbf{Map hygiene.} Building two \code{unordered\_map}s per centroid is
        \(O(\text{component})\); reserve them or use a sorted vector if the judge
        anti-hashes. The subtraction is over the \emph{same} \(\mathit{pairs}\) helper
        -- do not forget to drop \(u=v\) in both the \(\mathit{all}\) and each arm call.
  \item \textbf{Digit vs edge.} The single digit lives on the \emph{edge}, so a path
        of \(t\) edges reads \(t\) digits and has length \(t\); \(\mathit{dep}\) counts edges,
        matching the exponent of 10 exactly.
\end{itemize}
\end{hardeningbox}

\begin{recallbox}[24.2 in one breath]
\emph{Claim every path at its top centroid (\(\defref{c18-claim}\)); there its length
is a sum of two per-arm depths, counted by a query-then-merge histogram in
\(\bigO{n\log n}\) (\(\coderef{c18-count}\)) -- the upgrade of Ch.~13's \(\bigO{nK}\)
DP for CF~161D. The same skeleton survives any path value that factors into two
independent per-arm pieces: CF~715C reads the path as a decimal number, and
coprimality with 10 turns `concatenate and test divisible' into `match a forward
hash against a normalised backward hash, then subtract the same-child pairs'
(\(\coderef{c18-715c}\)).}
\selftest{
  \item Why does query-then-merge across arms already exclude same-child pairs? \emph{(A node only meets partners from strictly earlier arms and the centroid, so an intra-arm pair is never formed here (\(\secref{c18-count}\)).)}
  \item In CF~715C what makes the concatenation separable into single-node keys? \emph{\(\gcd(M,10)=1\), giving \(10^{-\mathit{dep}[v]}\) to move the length factor onto one side (\(\coderef{c18-715c}\)).)}
  \item Which pairs does the per-arm subtraction remove? \emph{Those with both endpoints in one child component of \(c\); they are claimed at a deeper centroid.}
}
\end{recallbox}

% =====================================================================
%  24.3  The centroid tree as a distance index -- one table per centroid.
%  Exports: sec:c18-index  def:c18-index  thm:c18-nn  lst:c18-index
%           prob:c18-293e  lst:c18-293e
%  Solved CF 293E "Close Vertices" (2700), live-verified via the problemset
%  API (tags data structures / divide and conquer / trees).  Statement:
%  count unordered pairs (u,v) whose UNIQUE path uses at most l edges AND has
%  total edge weight at most w.  Official samples read live and replayed:
%    S1: n=4, l=4, w=6, star centred at 1 (edges 1-2=3,1-3=4,1-4=3) -> 4.
%        Centroid 1 claims: (1,2),(1,3),(1,4),(2,4); (2,3),(3,4) break w. -> 4.
%    S2: n=6, l=2, w=17, tree 1-2(3),2-3(5),2-4(13),1-5(6),5-6(9) -> 9.
%        Top centroid 1 claims 6 pairs; recursion on {2,3,4} claims 2 and on
%        {5,6} claims 1;  6+2+1 = 9.  Both matched by hand.
%  The dynamic nearest-marked pattern (Xenia and Tree) is taught here in PROSE;
%  its URL is OWNED by Ch.12 12.3 (/342/E), so it is named, never re-linked.
%  Hand-verified on the running tree T (unweighted, hops): mark {5,13}, query
%  node 6 -> the centroid-ancestor min gives 2 (nearest marked is 5).
% =====================================================================
\section{The Centroid Tree as a Distance Index}\label{sec:c18-index}

\begin{previouslybox}[Ch.~16 answers a distance between two nodes; this is a distance to a \emph{set}]
Chapter~16 (\chref{c11}) turned \(\dist(u,v)\) into an \(\bigO{1}\) lookup once
depths and \(\lca\) are precomputed. But a pairwise oracle answers \emph{one} pair per
call. Two tasks recur and are not pairwise: ``\emph{mark} a node, then ask the distance
to the \emph{nearest} marked node'' (an online set that grows), and ``\emph{count} the
pairs within a distance budget'' (a whole-component sweep). Chapter~22 (\chref{c16})
handles a fixed path's \emph{range}; it does not index a dynamic set and cannot
enumerate all pairs. The centroid tree does: each centroid \(c\) keeps a distance table
for its whole component, and a node \(x\) sits in exactly the \(\bigO{\log n}\) tables
of its centroid ancestors (\(\thmref{c18-height}\)) --- so \emph{touching \(x\) once
touches \(\log n\) component-wide tables}, which is the index this section builds.
\end{previouslybox}

\begin{topicintro}
Root the view not at a vertex of \(T\) but at a centroid \(c\) and record, for every
node \(x\) of \(c\)'s live component, the single number \(d_c(x)=\dist(c,x)\). Do that
at every centroid of the decomposition and you own a family of distance tables with a
special property: for any node \(x\), the tables containing \(x\) are precisely those of
\(x\)'s ancestors in the centroid tree --- at most \(\lceil\log_2 n\rceil\) of them.
That is enough to answer two very different questions. \emph{Online:} to mark \(x\),
write \(d_c(x)\) into each ancestor table's running best; to ask the nearest marked node
to \(v\), take \(\min_c\big(\dist(v,c)+\text{best}[c]\big)\) over \(v\)'s ancestors ---
correct because the highest shared centroid on the path to the true nearest witness
realises equality in the triangle inequality (\(\thmref{c18-nn}\)). \emph{Offline:} to
count pairs at distance within a budget, sweep each centroid's table and use
\(\dist(u,v)=d_c(u)+d_c(v)\) for the pairs its arms straddle, then recurse --- this is
CF~293E, where the budget is \emph{two}-sided (\(\le l\) edges and \(\le w\) weight),
turning the \(\secref{c18-count}\) histogram into a sorted table plus a Fenwick over
weights.

\textbf{Prerequisites.} \(\secref{c18-tree}\) (the decomposition, \(\mathit{par}_c\),
\(\mathit{cd}\), and \(\thmref{c18-height}\)'s \(\bigO{\log n}\) bound) and
\(\secref{c18-count}\) (the top-centroid claim and the same-arm correction). The
nearest-marked query is the pattern the classic \emph{Xenia and Tree} problem needs; its
statement lives in Chapter~12 (\chref{c09}) and is named here, not re-linked.

\textbf{After this section you can}
\begin{itemize}[nosep]
  \item store, for each node, its list of centroid ancestors with hop-distances, in \(\bigO{n\log n}\); \hfill\lv{4}
  \item answer \emph{mark} and \emph{nearest-marked} online in \(\bigO{\log n}\) (each a
        scan of one node's centroid-ancestor list), and prove the min is exact (\(\thmref{c18-nn}\)); \hfill\lv{5}
  \item count pairs under a two-sided distance budget by sorting one metric and a Fenwick
        on the other, inside each centroid bucket; \hfill\lv{5}
  \item solve CF~293E end to end and hand-check it against both official samples. \hfill\lv{5}
\end{itemize}
\end{topicintro}

% ---------------------------------------------------------------
\subsection{One distance table per centroid}
% ---------------------------------------------------------------

\begin{definition}{The centroid distance index}{c18-index}
Run the decomposition of \(\secref{c18-tree}\). For each centroid \(c\), the
\emph{distance table} of \(c\) is the map \(x \mapsto d_c(x)=\dist(c,x)\) defined for
every \(x\) in \(c\)'s live component at the moment \(c\) was chosen. Store it
implicitly: for every node \(x\), keep the list
\(\mathit{up}(x)=\big[(c_0,d_{c_0}(x)),(c_1,d_{c_1}(x)),\dots\big]\) of \(x\)'s centroid
ancestors \(c_0=x\)'s own centroid, up to the top centroid, each paired with its
hop-distance. Because a centroid is deleted before recursion, \(\mathit{up}(x)\) has one
entry per centroid-tree level, so \(|\mathit{up}(x)|\le\lceil\log_2 n\rceil+1\).
\end{definition}

\noindent\textbf{In plain words.} Every node remembers \emph{which} centroids owned it
and \emph{how far} it sat from each. That short list --- logarithmic, fixed at build
time --- is the entire index. Reading it lets one node talk to whole components at once;
it is what turns ``distance from a set'' or ``count within a radius'' from \(O(n)\) per
call into \(O(\log n)\).

\begin{insightbox}[Why the top shared centroid is the one that answers]
Fix a query node \(v\) and any target node \(m\). Let \(z\) be the highest centroid
(smallest \(\mathit{cd}\)) lying on the path \(v\text{--}m\). At the moment \(z\) was
picked, both \(v\) and \(m\) were in its live component, so \(z\in\mathit{up}(v)\cap
\mathit{up}(m)\) and \(z\) is on the \(v\text{--}m\) path, giving
\(\dist(v,m)=\dist(v,z)+\dist(z,m)=d_z(v)+d_z(m)\) \emph{exactly}. For every other
centroid \(c\) that is an ancestor of \(v\), the triangle inequality gives
\(d_c(v)+d_c(m)\ge\dist(v,m)\). So a \(\min\) over \(v\)'s ancestors of
\(d_c(v)+\text{best}[c]\) (where \(\text{best}[c]\le d_c(m)\)) is never \emph{below} the
true nearest distance and is \emph{attained} at \(z\): the min is exact. This is the whole
magic --- the same list that indexes distances also witnesses them.
\end{insightbox}

\begin{theorem}{Nearest-marked via centroid ancestors is exact}{c18-nn}
Maintain \(\text{best}[c]=\min\{d_c(m): m \text{ marked}\}\) (initially \(+\infty\)) for
every centroid \(c\), updating it on a \emph{mark} at \(m\) by scanning \(\mathit{up}(m)\).
Then \(\min_{(c,\,d_c(v))\in\mathit{up}(v)}\big(d_c(v)+\text{best}[c]\big)\) equals
\(\dist(v,m^\ast)\) for the nearest marked node \(m^\ast\).
\end{theorem}
\begin{proofbox}
Every value in the min is \(\ge\dist(v,m^\ast)\) because \(\text{best}[c]=d_c(m)\) for
some marked \(m\), and \(d_c(v)+d_c(m)\ge\dist(v,m)\ge\dist(v,m^\ast)\). Conversely let
\(z\) be the highest centroid on the path \(v\text{--}m^\ast\). Marking \(m^\ast\) set
\(\text{best}[z]\le d_z(m^\ast)\), and \(z\in\mathit{up}(v)\); the insight box shows the
entry at \(z\) is \(d_z(v)+d_z(m^\ast)=\dist(v,m^\ast)\). The min is therefore both
\(\ge\) and \(\le\dist(v,m^\ast)\), hence equal. \(\checkmark\)
\end{proofbox}

\begin{lstlisting}[caption={Build the centroid distance index and answer mark / nearest-marked online. Each operation scans one node's centroid-ancestor list, so it costs O(log n). This is the Xenia-and-Tree pattern; dfs size and find centroid are exactly lst:c18-build. ASCII only}, label={lst:c18-index}]
const int MAXN = 100005, INF = 1e9;
vector<int> adj[MAXN];
vector<pair<int,int>> up[MAXN];    // (centroid ancestor, hop distance node -> it)
bool dead[MAXN]; int sz[MAXN];
int best[MAXN];                    // best[c] = min dist from centroid c to a marked node

int dfs_size(int u, int p){ sz[u]=1;
    for(int v: adj[u]) if(v!=p && !dead[v]) sz[u]+=dfs_size(v,u);
    return sz[u]; }
int find_centroid(int u, int p, int tot){        // Ch.12 walk-to-heavy, lst:c18-build
    bool moved=true;
    while(moved){ moved=false;
        for(int v: adj[u]) if(v!=p && !dead[v] && sz[v]>tot/2){ p=u; u=v; moved=true; break; } }
    return u; }
void collect_anc(int u, int p, int c, int d){    // record centroid c at distance d in u
    up[u].push_back({c, d});
    for(int v: adj[u]) if(v!=p && !dead[v]) collect_anc(v, u, c, d+1);
}
void build(int entry){
    int tot = dfs_size(entry, -1);
    int c   = find_centroid(entry, -1, tot);
    collect_anc(c, -1, c, 0);                    // distances from c to its live component
    dead[c] = true;                              // wall off c before recursing
    for(int v: adj[c]) if(!dead[v]) build(v);
}
// call  build(1);  then set best[c]=INF for all c.
void mark(int v){                                // O(log n): touch v's ancestor tables
    for(auto &a : up[v]) best[a.first] = min(best[a.first], a.second);
}
int query(int v){                                // O(log n): nearest marked distance
    int r = INF;
    for(auto &a : up[v]) r = min(r, best[a.first] + a.second);
    return r;                                    // exact by thm:c18-nn
}
\end{lstlisting}

\begin{tracebox}[The index on the running tree \(T\) -- mark \(\{5,13\}\), then query node \(6\)]
Ancestor lists (centroid, hop-distance) from the decomposition of \(T\):
\(\mathit{up}(5)=\{(4,2),(1,1),(5,0)\}\), \(\mathit{up}(13)=\{(4,3),(3,2),(9,1),(13,0)\}\),
\(\mathit{up}(6)=\{(4,1),(1,1),(6,0)\}\).\\
\textbf{mark 5:} \(\text{best}[4]=2,\ \text{best}[1]=1,\ \text{best}[5]=0\).\\
\textbf{mark 13:} \(3\not<2\), so \(\text{best}[4]\) stays \(2\); set \(\text{best}[3]=2,\ 
\text{best}[9]=1,\ \text{best}[13]=0\).\\
\textbf{query 6} reads \(\mathit{up}(6)\): at \(4\), \(1+\text{best}[4]=1+2=3\); at \(1\),
\(1+\text{best}[1]=1+1=\mathbf{2}\); at \(6\), \(\text{best}[6]=\infty\). The min is
\(\mathbf{2}\). Ground truth: \(\dist(6,5)=2\) via \(6\text{--}1\text{--}5\), and node
\(5\) is marked, so the nearest marked node to \(6\) is \(5\) at distance \(2\). The tight
entry is centroid \(1\), the top centroid on the path \(6\text{--}5\) --- exactly
\(\thmref{c18-nn}\)'s witness. \(\checkmark\)
\end{tracebox}

% ---------------------------------------------------------------
\subsection{Problem Lab: CF 293E -- counting pairs within a two-sided budget}
% ---------------------------------------------------------------

\solvedproblem{c18-293e}{Close Vertices}{CF 293E}
  {https://codeforces.com/problemset/problem/293/E}{CF 2700 L5}
  {centroid decomposition; a per-centroid (edges, weight) table; sort-by-edges with a Fenwick over weights; same-arm subtraction}

\begin{problemstatement}[Close Vertices (statement and samples verified against the live problem)]
A weighted tree on \(n\) vertices (\(1\le n\le 10^5\)) has a non-negative weight on each
edge (\(\le 10^4\)). The \emph{length} of a path is its edge \emph{count}; its
\emph{weight} is the sum of edge weights. Two vertices are \emph{close} if their unique
path has length at most \(l\) (\(1\le l\le n\)) and weight at most \(w\)
(\(0\le w\le 10^9\)). Count unordered pairs \((u,v)\), \(u<v\), that are close. Edges
are given parent-first: line \(i\) gives \(p_i, w_i\), the edge joining vertex \(i+1\) to
\(p_i\) with weight \(w_i\).

\textbf{Examples (official).}
\begin{verbatim}
Input:            Output:
4 4 6             4
1 3
1 4
1 3

Input:            Output:
6 2 17            9
1 3
2 5
2 13
1 6
5 9
\end{verbatim}
\end{problemstatement}

\begin{attemptgate}[Attempt it yourself first -- \lv{5}, give yourself 40 minutes]
Everything is still the \(\secref{c18-count}\) claim: charge each pair to its top
centroid, where the path splits into two arms. \peek{what two numbers per node does a
top-centroid pair combine?} \nudge{for a node \(x\) at the centroid record
\((a_x,b_x)=(\text{edges from }c\text{ to }x,\ \text{weight from }c\text{ to }x)\); a
straddling pair is close iff \(a_x+a_y\le l\) and \(b_x+b_y\le w\).}
\fullhint{Count, over the whole component's table, the pairs meeting BOTH bounds, then
subtract the same count restricted to each single arm (the same-arm correction of
\(\secref{c18-count}\)). The two-sided count of \((a,b)\) pairs with \(a_i+a_j\le l\) and
\(b_i+b_j\le w\) is a 2-D dominance problem: sort by \(a\), and answer with a Fenwick over
compressed \(b\) values. Prune any arm-walk the moment \(a>l\) or \(b>w\). Recurse; total
\(\bigO{n\log^2 n}\).}
\end{attemptgate}

\solstep{Understanding}
Two independent metrics bound the same path: a hop count and a weight sum. Chapter~22's
HLD counts pairs on \emph{one} chosen path; here we range over \emph{all} pairs, and each
pair must clear \emph{both} thresholds. That kills the plain distance histogram of
\(\secref{c18-count}\) --- it only tracked one integer. The fix keeps the centroid
framing (top-centroid claim, same-arm correction) but replaces the \(O(1)\) histogram
lookup with an ordered structure that answers ``how many stored \(b\) values are \(\le
B\)''.

\solstep{Planning}
At centroid \(c\), for every live node \(x\) store \((a_x,b_x)\). A pair is claimed by
\(c\) iff it straddles two different arms (or one endpoint is \(c\)); then
\(\dist=c\)-sums apply: \(a_x+a_y\le l\), \(b_x+b_y\le w\). By inclusion--exclusion,
\(\text{close pairs at }c=\mathit{count}(\text{all})-\sum_{\text{arm}}\mathit{count}(\text{arm})\),
where \(\mathit{count}(S)\) counts unordered pairs from the multiset \(S\) clearing both
bounds. Then mark \(c\) dead and recurse.

\solstep{Implementing}
\(\mathit{count}(S)\): sort \(S\) by \(a\). Each element \(i\) asks ``how many \(j\in S\)
have \(a_j\le l-a_i\) and \(b_j\le w-b_i\)'' --- a prefix in \(a\) combined with a prefix
in \(b\), i.e.\ 2-D dominance. Process these queries in increasing \(l-a_i\), inserting
each point's \(b\) into a Fenwick as its \(a\) becomes eligible; a query reads a Fenwick
prefix sum. This counts \emph{ordered} pairs including self-pairs; subtract the self-count
and halve. Compress the \(b\)-values of \(S\) (weights up to \(10^9\)).

\begin{lstlisting}[caption={CF 293E: count close pairs. At each centroid, build a per-node (edges, weight) table, count pairs clearing both budgets with a sort plus a Fenwick over compressed weights, and subtract the same-arm counts. dfs size and find centroid match lst:c18-build. ASCII only}, label={lst:c18-293e}]
#include <bits/stdc++.h>
using namespace std;
typedef long long ll;
const int MAXN = 100005;
int n; ll L, W;                       // L = max edges, W = max weight
vector<pair<int,int>> adj[MAXN];      // (neighbour, weight)
bool dead[MAXN]; int sz[MAXN];
ll ans = 0;

// unordered pairs i<j in pts with a_i+a_j<=L and b_i+b_j<=W
ll countPairs(vector<pair<ll,ll>>& pts){
    int m = pts.size(); if(m<=1) return 0;
    vector<ll> bs; for(auto&t:pts) bs.push_back(t.second);
    sort(bs.begin(),bs.end()); bs.erase(unique(bs.begin(),bs.end()),bs.end());
    int B = bs.size(); vector<int> bit(B+1,0);
    auto add=[&](int i,int v){ for(;i<=B;i+=i&-i) bit[i]+=v; };
    auto qry=[&](int i){ int s=0; for(;i>0;i-=i&-i)s+=bit[i]; return s; };
    sort(pts.begin(),pts.end());                       // by a ascending
    vector<pair<ll,int>> q(m);                         // (L - a_i, i)
    for(int i=0;i<m;i++) q[i]={L-pts[i].first, i};
    sort(q.begin(),q.end());
    ll ordered=0; int p=0;
    for(auto&qu:q){
        while(p<m && pts[p].first<=qu.first){          // a_j <= L - a_i
            int bidx=lower_bound(bs.begin(),bs.end(),pts[p].second)-bs.begin()+1;
            add(bidx,1); p++;
        }
        int bi=upper_bound(bs.begin(),bs.end(),W-pts[qu.second].second)-bs.begin();
        ordered += qry(bi);                            // b_j <= W - b_i
    }
    ll self=0;                                         // remove (i,i) matches
    for(auto&t:pts) if(2*t.first<=L && 2*t.second<=W) self++;
    return (ordered-self)/2;                           // ordered -> unordered
}
void collect(int u,int p,ll a,ll b,vector<pair<ll,ll>>& arm,vector<pair<ll,ll>>& all){
    if(a>L || b>W) return;                             // cannot clear a bound through c
    arm.push_back({a,b}); all.push_back({a,b});
    for(auto&e:adj[u]) if(e.first!=p && !dead[e.first])
        collect(e.first,u,a+1,b+e.second,arm,all);
}
ll process(int c){
    vector<pair<ll,ll>> all; all.push_back({0,0});     // the centroid itself
    vector<vector<pair<ll,ll>>> arms;
    for(auto&e:adj[c]){ if(dead[e.first]) continue;
        vector<pair<ll,ll>> arm;
        collect(e.first,c,1,e.second,arm,all);
        arms.push_back(arm);
    }
    ll tot = countPairs(all);
    for(auto&arm:arms) tot -= countPairs(arm);         // same-arm correction
    return tot;
}
int dfs_size(int u,int p){ sz[u]=1;
    for(auto&e:adj[u]) if(e.first!=p && !dead[e.first]) sz[u]+=dfs_size(e.first,u);
    return sz[u]; }
int find_centroid(int u,int p,int tot){
    bool moved=true;
    while(moved){ moved=false;
        for(auto&e:adj[u]) if(e.first!=p && !dead[e.first] && sz[e.first]>tot/2){ p=u;u=e.first;moved=true;break; } }
    return u; }
void decompose(int entry){
    int tot=dfs_size(entry,-1); int c=find_centroid(entry,-1,tot);
    ans += process(c);
    dead[c]=true;
    for(auto&e:adj[c]) if(!dead[e.first]) decompose(e.first);
}
int main(){
    scanf("%d %lld %lld",&n,&L,&W);
    for(int i=2;i<=n;i++){ int p,w; scanf("%d %d",&p,&w);
        adj[i].push_back({p,w}); adj[p].push_back({i,w}); }
    decompose(1);
    printf("%lld\n",ans);
}
\end{lstlisting}

\begin{tracebox}[Sample 2 by hand -- \(n=6,\ l=2,\ w=17\), tree \(1\!-\!2(3),2\!-\!3(5),2\!-\!4(13),1\!-\!5(6),5\!-\!6(9)\)]
\textbf{Top centroid \(1\)} (removing \(1\) gives pieces \(\{2,3,4\}\) size 3 and
\(\{5,6\}\) size 2, max \(3\le 3\)). Table \((a,b)\) at \(1\): \(1(0,0)\), \(2(1,3)\),
\(3(2,8)\), \(4(2,16)\), \(5(1,6)\), \(6(2,15)\). Count \(\mathit{all}\) with \(a_i+a_j\le2\),
\(b_i+b_j\le17\): with \(1\): \((1,2),(1,3),(1,4),(1,5),(1,6)\) all clear (5); across arms
\(2\text{--}5\): \(a\,1{+}1{=}2,\ b\,3{+}6{=}9\) clear (1); \(2\text{--}6\) and \(3,4\)
with \(5,6\) fail the edge bound (\(a>2\)); \(3\text{--}4\): \(b\,8{+}16{=}24\) fails.
So \(\mathit{count}(\text{all})=6\). Arm \(\{2,3,4\}\): every pair has \(a\ge1{+}2=3>2\)
\(\Rightarrow0\). Arm \(\{5,6\}\): \(a\,1{+}2=3>2\Rightarrow0\). Contribution \(6-0-0=6\).\\
\textbf{Recurse arm \(\{2,3,4\}\),} centroid \(2\): table \(2(0,0),3(1,5),4(1,13)\);
\(\mathit{count}(\text{all})\): \((2,3)\) ok, \((2,4)\) ok, \((3,4)\) \(b\,18>17\) fails
\(\Rightarrow2\); arms singleton \(\Rightarrow0\). Contribution \(2\).\\
\textbf{Recurse arm \(\{5,6\}\),} centroid \(5\): \(5(0,0),6(1,9)\); pair clears
\(\Rightarrow1\). Contribution \(1\).\\
\textbf{Total} \(6+2+1=9\). \(\checkmark\) matches the official output.
\end{tracebox}

\begin{complexitybox}[Two-sided distance counting]
Pruned by \(a\le l\) and \(b\le w\), each centroid's table is at most its component size;
\(\mathit{count}(\cdot)\) is \(\bigO{s\log s}\) (a sort and a Fenwick sweep) and runs on
\(\mathit{all}\) plus every arm, so per centroid \(\bigO{s\log s}\). A node is in
\(\bigO{\log n}\) components (\(\thmref{c18-height}\)), giving \(\bigO{n\log^2 n}\) time
overall and \(\bigO{n}\) extra space beyond the tables. The single \(\log n\) of
\(\secref{c18-count}\) became \(\log^2 n\) because the second metric needs an ordered
structure rather than an array index --- the price of a two-sided constraint.
\end{complexitybox}

\begin{warningbox}[Three ways this goes wrong]
\begin{itemize}[nosep]
  \item \textbf{Forgetting the same-arm subtraction.} \(\mathit{count}(\text{all})\) also
        counts pairs that live in one child arm, whose real top centroid is deeper. Those
        pairs are counted here \emph{and} down there --- double count. The
        \(\sum_{\text{arm}}\mathit{count}(\text{arm})\) line is not decoration.
  \item \textbf{Self-pairs.} \(\mathit{count}(\text{all})\) treats \(u=v\): \((c,c)\) and
        any node with \(2a\le l,\ 2b\le w\) sneak in. The \(\mathit{self}\) correction
        removes them before halving; skip it and every answer is too large.
  \item \textbf{Ordering vs unordered.} The dominance sweep counts ordered pairs
        \((i,j)\). Halving is only valid after self-pairs are gone; halve first and you
        floor away real pairs.
\end{itemize}
\end{warningbox}

\begin{hardeningbox}[Where this skeleton generalises]
\begin{itemize}[nosep]
  \item \textbf{One budget, count all.} Drop \(b\) and use an array histogram: this is
        exactly \(\secref{c18-count}\) (\(\coderef{c18-count}\)) --- the Fenwick is the
        histogram upgraded to a second, unordered metric.
  \item \textbf{Farthest instead of nearest.} Replace \(\min\) in \(\coderef{c18-index}\)'s
        query with \(\max\); the same ancestor list reports the farthest marked node (and
        the tree's diameter via \(\chref{c09}\)'s farthest baseline, cited not recomputed).
  \item \textbf{Sum over all, not count.} Change \(\mathit{count}(\cdot)\) to a Fenwick
        that accumulates \(b\)-sums as well as counts and you get ``sum of path weights
        over close pairs'' for free --- the Summit (\(\probref{c18-summit}\)) leans on this.
\end{itemize}
\end{hardeningbox}

\begin{recallbox}[24.3 in one breath]
\emph{Index the component, not the pair.} Each centroid owns a distance table; a node
lives in \(\bigO{\log n}\) of them (\(\defref{c18-index}\)). Mark/write and
nearest/query read that short list, exact by \(\thmref{c18-nn}\)'s top-shared-centroid
witness. Offline counting keeps the top-centroid claim but swaps the histogram for a
sort-plus-Fenwick so a second metric fits --- CF~293E counts pairs within \(\le l\) edges
\emph{and} \(\le w\) weight in \(\bigO{n\log^2 n}\) (\(\coderef{c18-293e}\)).
\end{recallbox}

% =====================================================================
%  24.4  Divide and conquer on trees beyond path sums -- an ordered
%  structure inside the centroid bucket.
%  Exports: sec:c18-beyond  def:c18-beyond  lst:c18-beyond
%           prob:c18-833d  lst:c18-833d
%  Solved CF 833D "Red-Black Cobweb" (2800), live-verified via the problemset
%  API (tags data structures / divide and conquer / implementation / trees).
%  Statement: tree, each edge a colour (0 red / 1 black) and a clamminess x.
%  A pair's (unique) path is kept when #(red) and #(black) "differ at most
%  twice", i.e. max(r,b) <= 2*min(r,b).  Jelliness = product over kept
%  unordered pairs of (product of x on the path), mod 1e9+7; empty -> 1.
%  KEY reduction (verified): with r=#red, b=#black on a path, the filter
%  r<=2b and b<=2r is TWO additive features u=2r-b>=0 and v=2b-r>=0; along a
%  centroid's two arms these sum, so a kept pair is a 2-D dominance pair --
%  the same shape as 293E, but the aggregate is a PRODUCT, folded as
%  prod_x w_x^{#partners(x)} and combined by G(all) / prod_arm G(arm).
%  Official sample 1 replayed by hand at centroid 2 -> 1265625 (pairs (1,3),
%  (1,5),(3,4),(4,5) with products 45,45,25,25).  Sample 2 -> 452841614.
%  Back-cites Ch.18 Fenwick (c14); the BIT is used, never re-taught.
% =====================================================================
\section{Beyond Path Sums: A Structure in the Bucket}\label{sec:c18-beyond}

\begin{previouslybox}[\(\secref{c18-count}\) and \(\secref{c18-index}\) both counted; a path \emph{product} is new]
Two sections back, the centroid bucket held a small integer array (a distance
histogram, \(\coderef{c18-count}\)); one section back it held a sorted table and a
Fenwick (a two-sided budget, \(\coderef{c18-293e}\)). Both \emph{counted} qualifying
pairs. Chapter~18 (\chref{c14}) owns the Fenwick itself; here we only bolt one into a
bucket. The step up in this section is the \emph{aggregate}: some questions are not
"how many" but "combine a per-path value over all qualifying paths" --- a product, a
sum of products, a maximum --- and the fold is not additive. CF~833D asks for a
\emph{product} over paths of a \emph{product} along each path, so we need the general
recipe: reduce a path property to per-node \emph{features} that sum across the two arms
at a centroid, decide acceptance with a dominance predicate on those features, and carry
whatever monoid the aggregate needs in the bucket.
\end{previouslybox}

\begin{topicintro}
A centroid decomposition is a divide-and-conquer, and like any one it is only as
general as the three choices you plug into it: the \emph{feature} each node contributes,
the \emph{acceptance test} for a straddling pair, and the \emph{accumulator} folded over
accepted pairs. The first two sections fixed all three (feature = distance, test = a
threshold, accumulator = count). Here we free them. The recurring trick: pick per-node
quantities that are \emph{additive along the path}, so that at the top centroid \(c\) a
pair \((x,y)\) has path-value \(f(x)+f(y)\); acceptance then says "\(\ge 0\) in two
directions," a 2-D dominance query answered by sorting one direction and sweeping a
Fenwick over the other --- the identical skeleton as \(\coderef{c18-293e}\), now with the
accumulator promoted from a count to a modular product. The same-arm correction survives
too, because a product over a set difference of pairs is a quotient of the two products
(\(\defref{c18-beyond}\)).

\textbf{Prerequisites.} \(\secref{c18-count}\) (the top-centroid claim and same-arm
subtraction), \(\secref{c18-index}\) (the 2-D dominance sweep of \(\coderef{c18-293e}\)),
and the Fenwick from Chapter~18 (\chref{c14}). Multiplicative inverses modulo a prime
(Fermat) are used to "subtract" a product; that is arithmetic, not a new data structure.

\textbf{After this section you can}
\begin{itemize}[nosep]
  \item state the general centroid D\&C template and name its three plug-ins (\(\defref{c18-beyond}\)); \hfill\lv{5}
  \item turn "the red and black counts differ by a factor of at most two" into two additive
        dominance features; \hfill\lv{5}--\lv{6}
  \item fold a path \emph{product} with \(\prod_x w_x^{\#\text{partners}(x)}\) and combine arms by
        a quotient of products; \hfill\lv{6}
  \item solve CF~833D and hand-check it against the official sample. \hfill\lv{6}
\end{itemize}
\end{topicintro}

% ---------------------------------------------------------------
\subsection{The template, with the plugs pulled out}
% ---------------------------------------------------------------

\begin{definition}{Centroid divide-and-conquer, in general}{c18-beyond}
Fix, for the whole algorithm: a per-node \emph{feature map} \(x\mapsto \phi_c(x)\)
computed at centroid \(c\) that is \emph{additive} in the sense
\(\phi_c(\text{path } x\text{--}y)=\phi_c(x)+\phi_c(y)\) for a pair straddling \(c\); an
\emph{acceptance predicate} on \(\phi_c(x)+\phi_c(y)\) that is a conjunction of
inequalities each monotone in one coordinate; and a commutative \emph{accumulator}
\(\oplus\) over per-pair values with identity \(e\). Then
\(\mathit{solve}(S)\): choose centroid \(c\); gather each arm's node features; combine
arms through an ordered structure to fold \(\oplus\) over accepted straddling pairs
\emph{minus} the pairs inside a single arm; recurse on the pieces, multiplying the
answers. \(\bigO{\log n}\) levels make the total the per-centroid combine times
\(\log n\).
\end{definition}

\noindent\textbf{In plain words.} Find the balance node, describe a path by numbers that
\emph{add} when you glue its two halves at the centroid, keep the pairs whose summed
numbers pass a "both coordinates non-negative" test, and fold whatever you are asked to
fold. Everything in \(\secref{c18-count}\) and \(\secref{c18-index}\) was one instance;
CF~833D is the instance where the folded value is a product and the feature is a colour
imbalance.

\begin{insightbox}[Why a product still permits the same-arm subtraction]
For a count, \(\#\text{cross}=\#(\text{all})-\sum\#(\text{arm})\) because the qualifying
pairs split into "inside one arm" and "straddling." For a \emph{product} the pairs
partition the same way, so the product over the union is the product of the parts:
\(G(\text{all})=G_{\text{cross}}\cdot\prod_{\text{arm}}G(\text{arm})\). Divide out the arm
terms and you have the straddling product,
\(G_{\text{cross}}=G(\text{all})\cdot\prod_{\text{arm}}G(\text{arm})^{-1}\), where the
inverse is modular (Fermat, since the modulus \(10^9{+}7\) is prime). Subtraction becomes
division --- the only change a multiplicative accumulator forces.
\end{insightbox}

\begin{lstlisting}[caption={The reusable centroid divide-and-conquer skeleton with its three plug-ins named. Feature (additive per node), accept (2-D dominance), and the accumulator fold are the only things that change between 24.2, 24.3 and 24.4. ASCII only}, label={lst:c18-beyond}]
// template -- instantiate FEATURE / ACCEPT / FOLD for each problem.
void solve_comp(int entry){
    int tot = dfs_size(entry,-1);
    int c   = find_centroid(entry,-1,tot);        // lst:c18-build
    // 1. FEATURE: walk each arm, compute phi_c(x) for every node.
    vector<Feat> all = {(0,0, ELEM_IDENTITY)};    // the centroid: neutral feature
    for (v in adj[c]) if (!dead[v]) {
        Feat cur = feature_from(c, v);            // additive along c->v
        collect(cur, v, c);                        // DFS, skipping dead[]
    }
    // 2. ACCEPT: keep pair {x,y} iff phi(x)+phi(y) passes a dominance test.
    // 3. FOLD:    acc = fold over accepted cross-arm pairs;  subtract by the
    //             same-arm trick -- counts: acc -= fold(arm);  products: acc /= fold(arm).
    result_c = combine(all) MINUS_or_DIVIDE sum over arms of combine(arm);
    dead[c] = true;                               // wall off c, then recurse
    for (v in adj[c]) if (!dead[v]) solve_comp(v);
    combine_into_global(result_c);                // '+' for counts, '*' for products
}
// 24.2: Feat = depth,            ACCEPT = depth sum == K,   FOLD = count,  subtract.
// 24.3: Feat = (edges, weight),  ACCEPT = both <= budget,    FOLD = count,  subtract.
// 24.4: Feat = (2r-b, 2b-r, w),  ACCEPT = both >= 0,         FOLD = product, divide.
\end{lstlisting}

\begin{observation}[Reading a new statement as an instance]
Before reaching for code, ask the three plug-in questions. \emph{Feature:} is the path
property a \(\sum\) of per-edge contributions that can be regrouped as (arm1)+(arm2)?
\emph{Accept:} is the filter a monotone test on those sums (a threshold, a congruence, a
two-sided bound)? \emph{Fold:} is the answer count / sum / product / max over accepted
pairs? If any answer is "no," the centroid bucket will not directly carry it --- go to
\(\secref{c18-boundary}\) and pick a different engine.
\end{observation}

% ---------------------------------------------------------------
\subsection{Problem Lab: CF 833D -- a colour-balance product over all paths}
% ---------------------------------------------------------------

\solvedproblem{c18-833d}{Red-Black Cobweb}{CF 833D}
  {https://codeforces.com/problemset/problem/833/D}{CF 2800 L5}
  {centroid decomposition; colour imbalance as two additive dominance features; product accumulator folded as prod of w to the partner-count; same-arm subtraction by modular division}

\begin{problemstatement}[Red-Black Cobweb (statement and samples verified against the live problem)]
A tree on \(n\) nodes (\(2\le n\le 10^5\)); each edge has a \emph{clamminess}
\(x\) (\(1\le x\le 10^9{+}6\)) and a colour \(c\in\{0,1\}\) (0 = red, 1 = black). For an
unordered pair of nodes, look at its unique path: let \(r,b\) be the counts of red and
black edges on it. The pair is \emph{good} if the two counts "differ at most twice,"
i.e.\ \(\max(r,b)\le 2\min(r,b)\). The \emph{jelliness} is the product, over all good
pairs, of the product of clamminess along that path, modulo \(10^9{+}7\) (print \(1\) if
no good pair exists).

\textbf{Examples (official).}
\begin{verbatim}
Input:            Output:
5                 1265625
1 2 9 0
2 3 5 1
2 4 5 0
2 5 5 1

Input:            Output:
8                 452841614
1 2 7 1
2 3 4 1
3 4 19 1
5 1 2 0
6 2 3 0
7 3 3 0
8 4 4 0
\end{verbatim}
\end{problemstatement}

\begin{attemptgate}[Attempt it yourself first -- \lv{6}, give yourself 50 minutes]
Translate "differ at most twice" into algebra before touching the tree.
\peek{if \(r\) reds and \(b\) blacks, what two inequalities say \(\max(r,b)\le 2\min(r,b)\)?}
\nudge{\(r\le 2b\) and \(b\le 2r\), i.e.\ \(2r-b\ge 0\) and \(2b-r\ge 0\). Both left sides are
\emph{sums} of a per-edge \(+2/-1\) value, so they add along a path split at a centroid.}
\fullhint{Give each node at centroid \(c\) the feature \((u,v,w)\) where \(u=2r-b\),
\(v=2b-r\) (\(r,b\) from \(c\) to the node) and \(w\) is the clamminess product from \(c\)
to the node. A straddling pair \(\{x,y\}\) is good iff \(u_x+u_y\ge 0\) and \(v_x+v_y\ge 0\)
--- a 2-D dominance count, swept like \(\coderef{c18-293e}\). The path product is
\(w_x w_y\), so the fold is \(G=\prod_x w_x^{\#\text{partners}(x)}\). Compute \(G(\text{all})\)
and divide by \(\prod_{\text{arm}}G(\text{arm})\) (modular inverse); \(c\) enters as
\((0,0,1)\). Multiply \(G\) into the answer, recurse. \(\bigO{n\log^2 n}\).}
\end{attemptgate}

\solstep{Understanding}
The filter is not a distance bound; it is a \emph{ratio} bound between two colour counts.
Write it as the pair of linear inequalities \(2r\ge b\) and \(2b\ge r\). Define the
imbalance features \(u=2r-b\) and \(v=2b-r\). Because \(r\) and \(b\) count edges of the
path, and edges of a straddling path split into the two arms at \(c\), both \(u\) and
\(v\) are additive: \(u(x\text{--}y)=u_x+u_y\). So "good" is exactly
"\(u_x+u_y\ge0\) and \(v_x+v_y\ge0\)." The \(w\) (clamminess product) is multiplicative,
but \(\log\)-wise additive; we keep it as a plain product.

\solstep{Planning}
Per centroid: list every node of the component as \((u,v,w)\) (plus the centroid as
\((0,0,1)\)). Count, for each node \(x\), its partners \(y\ne x\) in \emph{any} arm with
\(u_x+u_y\ge0,v_x+v_y\ge0\) --- a dominance query (sort by \(u\), Fenwick over compressed
\(v\)). Fold \(G=\prod_x w_x^{\mathrm{cnt}_x}\). Do this for the whole component and for
each single arm, then combine as \(G(\text{all})/\prod_{\text{arm}}G(\text{arm})\) to keep
only cross-arm (and with-centroid) pairs.

\solstep{Implementing}
The dominance sweep returns \(\mathrm{cnt}_x\) per node, but \(x\) counts \emph{itself}
when it is inserted with \(u_x\ge0,v_x\ge0\); subtract that self-match (it is not a pair).
Modular exponentiation turns a partner count into the product \(w_x^{\mathrm{cnt}_x}\).

\begin{lstlisting}[caption={CF 833D: centroid decomposition. Each node carries the additive colour-imbalance features u=2r-b, v=2b-r and the clamminess product w. A 2-D dominance sweep counts partners; the accumulator is a product, combined across arms by modular division. dfs size and find centroid match lst:c18-build. ASCII only}, label={lst:c18-833d}]
#include <bits/stdc++.h>
using namespace std;
typedef long long ll;
const int MAXN = 100005;
const ll MOD = 1000000007LL;
int n;
struct Edge{ int to, col; ll x; };
vector<Edge> adj[MAXN];
bool dead[MAXN]; int sz[MAXN];
ll ans = 1;

ll modpow(ll b, ll e){ ll r=1%MOD; b%=MOD;
    while(e){ if(e&1) r=r*b%MOD; b=b*b%MOD; e>>=1; } return r; }

struct Feat{ ll u, v, w; };            // u=2r-b, v=2b-r, w=product along c->node

// G = product over accepted unordered pairs {x,y} of (w_x * w_y)
ll merge(const vector<Feat>& pts){
    int m = pts.size(); if(m <= 1) return 1;
    vector<ll> vs; for(auto&t:pts) vs.push_back(t.v);
    sort(vs.begin(),vs.end()); vs.erase(unique(vs.begin(),vs.end()),vs.end());
    int B = vs.size(); vector<int> bit(B+1,0);
    auto add=[&](int i,int d){ for(;i<=B;i+=i&-i) bit[i]+=d; };
    auto qry=[&](int i){ int s=0; for(;i>0;i-=i&-i) s+=bit[i]; return s; };
    vector<int> ins(m), ord(m); for(int i=0;i<m;i++) ins[i]=ord[i]=i;
    sort(ins.begin(),ins.end(),[&](int a,int b){return pts[a].u > pts[b].u;});   // by u desc
    sort(ord.begin(),ord.end(),[&](int a,int b){return pts[a].u < pts[b].u;});   // queries: T=-u desc
    ll res=1; int p=0;
    for(int qi=0; qi<m; qi++){
        int x=ord[qi]; ll T=-pts[x].u;                      // partners need u_y >= T
        while(p<m && pts[ins[p]].u >= T){                  // insert eligible u_y
            int bi=lower_bound(vs.begin(),vs.end(),pts[ins[p]].v)-vs.begin()+1;
            add(bi,1); p++;
        }
        int below = lower_bound(vs.begin(),vs.end(),-pts[x].v)-vs.begin(); // distinct v < -v_x
        ll cnt = p - qry(below);                           // inserted with v_y >= -v_x
        if(pts[x].u >= 0 && pts[x].v >= 0) cnt--;          // drop the self-match
        res = res * modpow(pts[x].w, cnt) % MOD;           // fold w_x^{cnt_x}
    }
    return res;
}
void collect(int u,int p,int r,int b,ll w,vector<Feat>& arm,vector<Feat>& all){
    arm.push_back({2*r-b, 2*b-r, w}); all.push_back({2*r-b, 2*b-r, w});
    for(auto&e:adj[u]) if(e.to!=p && !dead[e.to])
        collect(e.to,u, r+(e.col==0), b+(e.col==1), w*e.x%MOD, arm, all);
}
void process(int c){
    vector<Feat> all; all.push_back({0,0,1});               // centroid itself
    vector<vector<Feat>> arms;
    for(auto&e:adj[c]) if(!dead[e.to]){
        vector<Feat> arm;
        collect(e.to,c, e.col==0, e.col==1, e.x%MOD, arm, all);
        arms.push_back(arm);
    }
    ll tot = merge(all);
    for(auto&arm:arms) tot = tot * modpow(merge(arm), MOD-2) % MOD;   // divide same-arm
    ans = ans * tot % MOD;
}
int dfs_size(int u,int p){ sz[u]=1;
    for(auto&e:adj[u]) if(e.to!=p && !dead[e.to]) sz[u]+=dfs_size(e.to,u);
    return sz[u]; }
int find_centroid(int u,int p,int tot){
    bool moved=true;
    while(moved){ moved=false;
        for(auto&e:adj[u]) if(e.to!=p && !dead[e.to] && sz[e.to]>tot/2){ p=u;u=e.to;moved=true;break; } }
    return u; }
void decompose(int entry){
    int tot=dfs_size(entry,-1); int c=find_centroid(entry,-1,tot);
    process(c);
    dead[c]=true;
    for(auto&e:adj[c]) if(!dead[e.to]) decompose(e.to);
}
int main(){
    scanf("%d",&n);
    for(int i=1;i<n;i++){ int u,v,c; ll x; scanf("%d %d %lld %d",&u,&v,&x,&c);
        adj[u].push_back({v,c,x}); adj[v].push_back({u,c,x}); }
    decompose(1);
    printf("%lld\n", ans);
}
\end{lstlisting}

\begin{tracebox}[Sample 1 by hand -- centroid \(2\), features \((u,v,w)\) with \(u=2r-b,\ v=2b-r\)]
Edges: \(1\!-\!2\) red \(x9\); \(2\!-\!3\) black \(x5\); \(2\!-\!4\) red \(x5\);
\(2\!-\!5\) black \(x5\). At centroid \(2\) every other node is its own arm. Features:
\(c=2(0,0,1)\); \(1:(r1,b0)\Rightarrow(2,-1,9)\); \(3:(0,1)\Rightarrow(-1,2,5)\);
\(4:(1,0)\Rightarrow(2,-1,5)\); \(5:(0,1)\Rightarrow(-1,2,5)\). A straddling pair is good
iff \(u\)-sum \(\ge0\) and \(v\)-sum \(\ge0\):
\((1,3):2{-}1{=}1,\ {-}1{+}2{=}1\) ok; \((1,4):4,\ {-}2\) \(\times\); \((1,5):1,1\) ok;
\((3,4):1,1\) ok; \((3,5):{-}2,\ \times\); \((4,5):1,1\) ok; centroid pairs
\((2,\cdot)\) all fail (a one-colour path has \(2b\ge r\) or \(2r\ge b\) violated). Good
pairs: \(\{1,3\},\{1,5\},\{3,4\},\{4,5\}\) with path products \(45,45,25,25\). Partner
counts \(1{:}2,\,3{:}2,\,4{:}2,\,5{:}2,\,2{:}0\), so
\(G=9^2\!\cdot5^2\!\cdot5^2\!\cdot5^2 = 81\cdot625\cdot\ldots\). Grouping by pair
\(\prod w_x w_y=45\!\cdot\!45\!\cdot\!25\!\cdot\!25=1265625\). Each arm is a single node,
\(G(\text{arm})=1\). Recursion adds nothing. Answer \(1265625\). \(\checkmark\)
\end{tracebox}

\begin{complexitybox}[The product fold]
\(\mathit{merge}\) is a sort and a dominance sweep, \(\bigO{s\log s}\) on a size-\(s\)
list; it runs on \(\mathit{all}\) once and on each arm, and each modular inverse is
\(\bigO{\log \text{MOD}}\). A node lies in \(\bigO{\log n}\) components
(\(\thmref{c18-height}\)), so the total is \(\bigO{n\log^2 n + n\log n\log\text{MOD}}\),
comfortably inside 6 s for \(n=10^5\). The extra \(\log\) versus \(\secref{c18-count}\) is
the dominance structure the two-feature filter demands; the extra modular inverses are the
multiplicative accumulator's price.
\end{complexitybox}

\begin{warningbox}[What breaks the product version]
\begin{itemize}[nosep]
  \item \textbf{Dividing by zero.} \(G(\text{arm})\) is a product of \(w\)-powers; a
        single-node arm gives \(G=1\) (identity), never \(0\), since clamminess
        \(x\ge1\) and the modulus is prime. If you let a \(w\equiv0\pmod{\text{MOD}}\)
        slip in, the inverse vanishes --- but \(1\le x\le 10^9{+}6<\text{MOD}\) forbids it.
  \item \textbf{The self-match.} The sweep inserts node \(x\) and lets it dominate itself;
        forgetting the \(\mathit{cnt}--\) for \(u_x,v_x\ge0\) inflates an exponent and
        squares a \(w_x\) that should appear once per genuine partner.
  \item \textbf{Ratio vs difference.} "Differ at most twice" is a \emph{ratio}
        (\(\max\le2\min\)), not \(|r-b|\le2\). Misread it as a difference and the features
        are wrong; the sample's single-edge paths (diff \(1\)) would then be counted, but
        they are \emph{not} good here.
\end{itemize}
\end{warningbox}

\begin{hardeningbox}[Other problems on the same three plugs]
\begin{itemize}[nosep]
  \item \textbf{Sum of path products, not product.} Swap the accumulator to \(+\) and fold
        \(w_x\cdot(\text{partner weight-sum})\); the dominance sweep needs a prefix of
        \emph{products} rather than counts in the Fenwick --- still one structure per bucket.
  \item \textbf{Congruence filters.} A path divisible-by-\(m\) test (\(\secref{c18-count}\)'s
        CF~715C) is an \emph{equality} on the feature, answered with a hash map instead of a
        Fenwick: the template is the same, only the ordered structure changes to unordered.
  \item \textbf{Farthest-with-property.} Max-accumulator instances keep the centroid claim
        and let the bucket be a single running maximum --- no dominance needed when the
        feature is one-dimensional, which is why those problems read as easier.
\end{itemize}
\end{hardeningbox}

\begin{recallbox}[24.4 in one breath]
\emph{Pull the plugs out.} A centroid bucket can carry an \emph{additive} per-node feature,
accept pairs by a \emph{dominance} test, and fold \emph{any} commutative accumulator --- even
a product, whose same-arm subtraction becomes modular \emph{division}
(\(\defref{c18-beyond}\)). CF~833D reduces "colours differ by a factor of two" to
\(2r-b\ge0,\ 2b-r\ge0\) and answers the jelliness in \(\bigO{n\log^2 n}\)
(\(\coderef{c18-833d}\)). You now own the full divide-and-conquer-on-trees template.
\end{recallbox}

% =====================================================================
%  24.5  The Engine Boundary: when centroid decomposition is the WRONG tool.
%  WORKED EXAMPLES ONLY -- NO \solvedproblem (in-section count stays 4).
%  Exports: sec:c18-boundary  def:c18-map  lst:c18-map
%  DISCHARGES (prose + \chref/\secref only, NO re-link):
%    * Ch.22 (c16) \secref{c16-boundary}: "path COUNTING / distance stats is a
%      centroid job, not HLD" -> answered by \secref{c18-count}/\secref{c18-index}.
%    * Ch.23 (c29) "over ALL paths -> centroid" forward-cite -> \secref{c18-beyond}.
%    * Ch.16 (c11) LCA: the pairwise oracle the centroid index generalises.
%  Stays OUT (name + forward-cite, no code): single fixed-path RANGE updates ->
%    HLD Ch.22 (c16); subtree multiset unions -> DSU-on-tree Ch.17 (c28); a state
%    per node -> tree DP / rerooting Ch.13 (c10) / Ch.20 (c20); offline reorderable
%    node-set queries -> Mo Ch.27 (c21); dynamic topology -> link-cut Ch.25 (c23).
%  Do NOT \term anything already indexed ("centroid", "centroid tree" -- Ch.12 /
%    this chapter's 00/01).  This section uses ZERO \term.
% =====================================================================
\section{The Engine Boundary: When Centroid Is the Wrong Tool}\label{sec:c18-boundary}

\begin{previouslybox}[Four sections built one engine; this one says when to leave it]
Sections 24.1--24.4 made the centroid decomposition do a great deal: build the
\(\bigO{\log n}\)-deep centroid tree, count pairs by distance, index a component for
nearest-marked queries, and fold a product over colour-balanced paths. It is seductive
for exactly that reason --- once you can ``split at the balance and recurse,'' every
tree problem with a path in it starts to look centroid-shaped. But Chapter~22 (\chref{c16})
owns path \emph{ranges}, Chapter~17 (\chref{c28}) owns subtree-set \emph{unions},
Chapter~13 (\chref{c10}) owns per-node \emph{DP states}, and Chapter~25 (\chref{c23})
owns \emph{changing} trees --- and each of those beats centroid decomposition at its own
job, often by a full \(\log n\). A boundary section earns its place because the wrong
pick here is rarely a wrong \emph{answer}; it is a \(\tle\) or a needlessly \(n\log^2 n\)
where an \(n\log n\) or even \(n\) was available.
\end{previouslybox}

\begin{topicintro}
Three questions put any tree problem on the map below. \emph{What is the unit of work?}
A single rooted subtree is a DP (\(\chref{c10}\)); a chosen path is a range (HLD,
\(\chref{c16}\)); a whole set of marked nodes is an index (centroid, \(\secref{c18-index}\));
\emph{all} simple paths is centroid territory (\(\secref{c18-count}\), \(\secref{c18-beyond}\)).
\emph{Is the tree fixed?} Centroid decomposition is built once; if edges appear and
vanish between queries, it is invalid --- that is link-cut (\(\chref{c23}\)). \emph{Do you
merge child sets or split the component?} Bottom-up multiset merging is DSU-on-tree /
sack (\(\chref{c28}\)); top-down splitting by balance is centroid. The tell is directional:
sack \emph{accumulates} children into a parent, centroid \emph{dissolves} a component at
its middle.

\textbf{Prerequisites.} The whole of 24.1--24.4; Chapter~22 (\(\secref{c16-boundary}\)) and
Chapter~23 (\(\secref{c29-boundary}\)) for the two neighbouring boundary maps that
forward-cited here; Chapter~16 (\chref{c11}) for the pairwise \(\dist\) oracle.
\end{topicintro}

% ---------------------------------------------------------------
\subsection{The axis that decides the engine: what the query ranges over}
% ---------------------------------------------------------------

\begin{definition}{The tree-path engine map}{c18-map}
Classify by the \emph{set the query quantifies over} and whether the \emph{topology is
fixed}. \(\{\text{all simple paths},\ \text{property additive along path},\ \text{fixed tree}\}\)
is centroid decomposition (this chapter). \(\{\text{one chosen path},\ \text{range
update/query},\ \text{fixed tree}\}\) is heavy-light (\(\chref{c16}\)).
\(\{\text{subtree},\ \text{merge child multisets}\}\) is DSU-on-tree (\(\chref{c28}\)).
\(\{\text{per node},\ \text{local state recurrence}\}\) is tree DP / rerooting
(\(\chref{c10}\), \(\chref{c20}\)). \(\{\text{arbitrary node sets},\ \text{offline,
reorderable}\}\) is Mo on trees (\(\chref{c21}\)). \(\{\text{path or subtree},\ \text{edges
change}\}\) is link-cut (\(\chref{c23}\)). The cells below name the tell that selects one.
\end{definition}

\begin{center}
\small
\renewcommand{\arraystretch}{1.26}
\setlength{\tabcolsep}{4pt}
\begin{tabular}{@{}p{0.30\linewidth} p{0.24\linewidth} >{\raggedright\arraybackslash}p{0.24\linewidth} >{\centering\arraybackslash}p{0.14\linewidth}@{}}
\toprule
\textbf{Symptom in the spec} & \textbf{Right engine} & \textbf{Why centroid loses} & \textbf{Where}\\
\midrule
count/aggregate over \emph{all} paths, value additive along the path & centroid decomposition & --- it wins; this is its home & \(\secref{c18-count}\)/\(\secref{c18-beyond}\)\\
nearest / farthest marked node, online set, fixed tree & centroid distance index & --- it wins; \(\bigO{\log n}\) per touch & \(\secref{c18-index}\)\\
update and query a \emph{single} path, op is a monoid & heavy-light + segtree & centroid reprocesses every node per query, \(\bigO{n}\); no per-path laziness & Ch.~22 \(\chref{c16}\)\\
answer for each \emph{subtree} a set statistic (distinct colours, \(k\)-th) & DSU-on-tree / sack & centroid splits by balance, not by subtree; sack keeps \(\bigO{n\log n}\) merge & Ch.~17 \(\chref{c28}\)\\
a \emph{state per node} from a local recurrence (independent set, matching) & tree DP / rerooting & no distance / all-pairs dimension to exploit; a fold is enough & Ch.~13 \(\chref{c10}\)\\
many \emph{offline} queries over node sets, order free, no inverse & Mo on trees & centroid is online but has no handle on an arbitrary node subset & Ch.~27 \(\chref{c21}\)\\
path/subtree queries but edges are \emph{linked and cut} between ops & link-cut trees & the centroid tree is built once and dies when topology moves & Ch.~25 \(\chref{c23}\)\\
over all paths, but the per-node value is a line envelope & centroid bucket \emph{hosting} Li Chao & plain centroid folds scalars; combine the two & Ch.~23 \(\chref{c29}\)\\
\bottomrule
\end{tabular}
\end{center}

\begin{lstlisting}[language={}, basicstyle=\ttfamily\small, numbers=none, frame=single,
  caption={Decision ladder for tree-path problems. Read top to bottom and stop at the first branch that matches. Hands off to Chapters 13/17/22/25/27 whenever a different engine already owns the cell; centroid decomposition is the LAST resort for "all paths," not the first for "a path."}, label={lst:c18-map}]
  What does the query RANGE over?

  +-- A SINGLE chosen path, and you UPDATE + QUERY it (sum/min/assign on a subpath):
  |     -> HEAVY-LIGHT DECOMPOSITION + a segment tree (Ch.22, c16).
  |        Centroid is wrong: it has no notion of "the nodes on path u--v only"
  |        and would touch O(n) buckets to answer one path.

  +-- A SUBTREE (everything under some root), usually a SET statistic:
  |     -> DSU-ON-TREE / SACK (Ch.17, c28) or Euler-tour + Fenwick (Ch.18, c14).
  |        Centroid splits the tree at its CENTRE, top-down; sack merges children
  |        bottom-up.  "merge the children's sets" is sack's signature, not centroid's.

  +-- A STATE PER NODE from a LOCAL recurrence (max independent set, tree matching):
  |     -> TREE DP (Ch.13, c10); changing the root -> REROOTING (Ch.20, c20).
  |        There is no all-pairs / distance term, so halving buys nothing -- a
  |        single O(n) DFS fold already solves it.

  +-- A WHOLE SET of marked nodes and you ask DISTANCE-TO-SET, online:
  |     -> CENTROID DISTANCE INDEX (this ch., 24.3).  Mark/query walk the O(log n)
  |        centroid ancestors of one node.  (Fixed tree only -- see the OUT box.)

  +-- ALL simple paths, with a value that ADDS along the path and a filter that is
  |   a threshold / congruence / dominance on that sum:
  |     -> CENTROID PAIR COUNTING (24.2) or the general D&C-on-trees template with
  |        an ordered structure in the bucket (24.4).  Claim each path at its top
  |        centroid; a rooted DP cannot see paths that cross its root.

  +-- MANY OFFLINE queries over arbitrary node sets, you may REORDER them, aggregate
  |   has no inverse:
  |     -> MO ON TREES (Ch.27, c21).  Centroid answers online but cannot index an
  |        unconstrained subset of nodes chosen by the query.

  STILL OUT OF SCOPE for centroid decomposition (forward-cites only, no code):
    * edges LINKED / CUT between operations -> LINK-CUT TREES (Ch.25, c23); the
      centroid tree is computed once on a FIXED topology and is invalid after an edit.
    * the aggregate over a path is an ENVELOPE of lines queried at a variable x ->
      LI CHAO on a tree (Ch.23, c29); you may still host it in a centroid bucket, but
      the engine doing the work is Li Chao, not the balance split.
\end{lstlisting}

% ---------------------------------------------------------------
\subsection{Worked boundary calls (read a spec, name the engine)}
% ---------------------------------------------------------------

\begin{comparisonbox}[Five specifications, five different answers -- none solved here]
\textbf{(A) ``\(n\le10^5\), \(K\le n\); count unordered pairs at distance exactly \(K\).''}
The unit of work is \emph{all} pairs and the value (distance) is additive along the path.
Chapter~13's \(\bigO{nK}\) subtree DP times out because \(K\) is not small. This is
\(\secref{c18-count}\): claim each pair at its top centroid and read a depth histogram ---
\(\bigO{n\log n}\), independent of \(K\). Centroid is the right call.

\textbf{(B) ``fixed tree, node values; `add \(v\) to every node on path \(u\)--\(w\)', `query the sum on path \(p\)--\(q\)','' }
A single chosen path is \emph{updated and queried}, and the operation (add, sum) is a
monoid that composes into a lazy tag. Heavy-light plus a lazy segment tree
(Ch.~22 \(\chref{c16}\), over Ch.~20 \(\chref{c20}\)) is the answer. Centroid is
\emph{wrong}: it cannot restrict its work to the nodes of one path, so it degrades to
touching \(\bigO{n}\) buckets per query with no laziness to amortise them.

\textbf{(C) ``\(n\le10^5\) coloured tree; \(q\) offline queries `how many distinct colours in the subtree of \(v\)'.''}
The domain is a \emph{subtree} and the statistic merges the children's colour sets
bottom-up --- sack's signature (Ch.~17 \(\chref{c28}\)), or an Euler-tour interval plus a
Fenwick (Ch.~18 \(\chref{c14}\)). Centroid splits at the centre, not at a subtree
boundary, so it has no natural handle on "the descendants of \(v\)." Wrong tool.

\textbf{(D) ``operations `paint node black' and `distance from \(v\) to the nearest black node', online, fixed tree.''}
This is the \emph{nearest-marked} pattern (the classic \emph{Xenia and Tree}; its
statement is Ch.~12's, \chref{c09}). A per-centroid distance table (\(\secref{c18-index}\))
answers it in \(\bigO{\log n}\) per operation --- centroid is exactly right. Note the
contrast with (B): both are "online on a fixed tree," yet (B) is a path \emph{range} (HLD)
while (D) is a distance-to-\emph{set} (centroid). The unit of work, not the online-ness,
picks the engine.

\textbf{(E) ``\(n\le10^5\); `link \(u\)--\(v\)', `cut the edge at \(v\)', `query path aggregate', interleaved.''}
The topology moves. Any decomposition built once --- centroid \emph{or} heavy-light --- is
invalid after an edit. This is link-cut trees (Ch.~25 \(\chref{c23}\)). It is the sharpest
boundary this section draws: centroid decomposition assumes a \emph{fixed} tree, full stop.
\end{comparisonbox}

\begin{observation}[The all-pairs tell, and why a fixed root blinds a DP]
A rooted tree DP at \(r\) can fold anything \emph{inside} a subtree, but a path whose two
endpoints lie in \emph{different} child-subtrees of \(r\) passes through \(r\) and is never
seen by a single bottom-up pass --- and if it does not pass through \(r\), it lives in one
child and the DP misses it entirely unless it recurses anyway. Centroid decomposition
fixes the view instead of rooting it: the top-centroid claim (\(\defref{c18-claim}\)) makes
every pair appear exactly once, at the centroid that separates it. That is the whole reason
"quantifies over all paths" maps here and "quantifies over one subtree" maps to Ch.~13.
\end{observation}

\begin{notebox}[Centroid versus heavy-light: the honest tie-break]
Both \(\bigO{\log n}\)-decompose a tree, and beginners conflate them. They answer
\emph{orthogonal} questions. Heavy-light cuts a tree into root-to-leaf \emph{chains} so
that any single path \(u\)--\(v\) becomes \(O(\log n)\) contiguous array ranges --- it is a
\emph{path-as-range} engine, perfect for update-and-query on one path with a lazy
structure (Ch.~22). Centroid cuts at the \emph{balance point} so that any node becomes
\(O(\log n)\) ancestors across \emph{whole components} --- a \emph{component-as-index}
engine, perfect for all-pairs counts and distance-to-a-set. If the spec says ``the path
\(u\)--\(v\),'',' think heavy-light; if it says ``all pairs'' or ``nearest such node,'','
think centroid. Neither replaces the other, and (B) versus (D) above is the exercise that
proves it.
\end{notebox}

\begin{stuckbox}[If you reach for centroid but the tree has \(n\le 2000\)]
Small \(n\) is a licence to brute force. An all-pairs problem on \(n\le2000\) is a
\(\bigO{n^2}\) double loop with \(\dist\) from Ch.~16's \(\bigO{1}\) LCA oracle --- no
decomposition, no histogram, far less to get wrong. Reach for the centroid engine only when
the \(\bigO{n^2}\) or \(\bigO{nK}\) rival genuinely does not finish; a \(\bigO{n\log^2 n}\)
template that you mis-remember is worse than the correct quadratic you can write in three
lines. The cost of the technique is real code (deleted markers, per-arm subtraction, a
bucket structure), and it is only amortised by scale.
\end{stuckbox}

\begin{recallbox}[24.5 in one breath]
\emph{Pick by what the query ranges over, not by how ``tree-y'' it feels:} one chosen path
with updates \(\to\) heavy-light (Ch.~22); a subtree's set \(\to\) DSU-on-tree (Ch.~17) or
Euler-tour (Ch.~18); a per-node local state \(\to\) tree DP (Ch.~13); arbitrary offline node
sets \(\to\) Mo (Ch.~27); moving edges \(\to\) link-cut (Ch.~25). Centroid owns the two cells
nobody else can: \emph{all} paths with an additive value (24.2/24.4) and distance-to-a-set on
a fixed tree (24.3). Discharges the forward-cites from Ch.~22 (\(\secref{c16-boundary}\))
and Ch.~23 (\(\secref{c29-boundary}\)) that punted "all-pairs / counting paths" to this
chapter --- those promises now sit in \(\secref{c18-count}\), \(\secref{c18-index}\) and
\(\secref{c18-beyond}\).
\selftest{\item A query asks the maximum on path \(u\)--\(v\) after point updates.  Centroid
or heavy-light? \emph{(Heavy-light: a single path is being queried, and the max is a composable
monoid; centroid has no per-path laziness.)}
\item Same tree, but the query is ``distance from \(v\) to the closest marked node,'' marks
change online.  Why does the answer switch? \emph{(Now the set of marked nodes is the object,
not a path --- the centroid distance index of 24.3 answers in \(\bigO{\log n}\) per touch.)}
\item Edges are added and removed between queries.  Why is centroid decomposition disqualified
entirely? \emph{(It is built once on a fixed topology; any link/cut invalidates every
\(\mathit{par}_c\) table --- that is link-cut territory, Ch.~25.)}}
\end{recallbox}

% =====================================================================
%  Chapter 24 review -- the 14-station revision contract.
%  Revision-only: tables, one-liners, pointers back via \secref/\thmref/
%  \defref/\figref/\coderef/\probref.  No re-derivations.
%  Summit: CF 757G "Can Bash Save the Day?" (3400, L6) -- a weighted tree plus a
%  permutation a[1..n].  Query 1 l r v = sum over i in [l,r] of dist(a_i, v);
%  query 2 x = swap a_x and a_{x+1}; queries are XOR-encoded against the last
%  answer (online).  Centroid decomposition with a Fenwick / segment tree per
%  centroid over POSITIONS, storing dist(a_i, centroid), plus the child-side
%  "sub" subtraction so each pair is claimed at its top shared centroid.  Both
%  official samples hand-replayed (XOR-decoded): answers 23, 37, 28.
%  Exports prob:c18-summit, lst:c18-kit, lst:c18-summit.
%  Mixed set is EXACTLY 10, all live-verified, un-owned, sorted ASCENDING by
%  effective difficulty; each hint names its reduction and cites a c18 listing.
%  CF 342E (Ch.12), CF 161D (Ch.13), CF 914D/1601C/1743F and the whole
%  Ch.19--23 rosters are OWNED elsewhere -- never re-linked here.
% =====================================================================

\chaptersection{Chapter 24 Summary}

One sentence: \emph{a rooted DP can only see paths that pass through its own
root, so it is blind to the pairs that cross it; centroid decomposition fixes the
blindness by splitting every tree at its \emph{balance point}, claiming each path
at the single highest centroid it passes through, solving the pieces independently,
and merging a per-component distance table in \(\bigO{\log n}\) levels
(\(\secref{c18-tree}\), \(\secref{c18-count}\)).} Chapter~12 (\chref{c09}) proved a
centroid exists and finds one in linear time; this chapter starts at ``remove it
and recurse'' and ends with a live index and a product fold. Each row below is a
pointer, not a proof.

\begin{center}
\renewcommand{\arraystretch}{1.25}
\begin{tabular}{@{}p{0.24\linewidth} p{0.48\linewidth} p{0.20\linewidth}@{}}
\toprule
\textbf{Result} & \textbf{Statement} & \textbf{Where} \\
\midrule
Decomposition & Remove the centroid, mark it dead, recurse on each surviving piece; every
piece has \(\le n/2\) nodes so there are \(\bigO{\log n}\) levels and \(\bigO{n\log n}\) total
work. & \defref{c18-decomp}, \coderef{c18-build} \\
Centroid tree height & Parent \(\mathit{par\_c}\) links each centroid to the one that split
its piece; the tree has height \(\le \lceil\log_2 n\rceil\), so a node touches \(\bigO{\log n}\)
centroids. & \thmref{c18-height}, \figref{c18-centtree} \\
Top-centroid claim & Every simple path is claimed by exactly one centroid -- the highest one it
passes through -- so counting over all paths reduces to a per-centroid merge plus recursion. & \defref{c18-claim}, \secref{c18-count} \\
Distance histogram count & Collect each arm's depths, query-then-merge a histogram, and
\emph{subtract} the same-arm pairs; \(\bigO{n\log n}\) for all pairs at distance \(K\)
(\(\le K\)). & \coderef{c18-count}, \coderef{c18-715c} \\
Centroid distance index & One distance table per centroid; a node lives in its \(\bigO{\log n}\)
centroid-ancestor tables, so mark / nearest / within-budget queries are \(\bigO{\log n}\). & \defref{c18-index}, \thmref{c18-nn} \\
Beyond paths & Feature (additive along the arm), accept (a dominance / threshold test), and fold
(accumulator) are the only plugs that change from counting to products. & \defref{c18-beyond}, \coderef{c18-833d} \\
Engine map & Path \emph{ranges} are HLD, subtree \emph{sets} are sack, per-node \emph{states}
are DP; centroid owns \emph{all paths} and \emph{distance-to-set}. & \defref{c18-map}, \secref{c18-boundary} \\
\bottomrule
\end{tabular}
\end{center}

\medskip
{\small\textbf{Algorithms from this chapter.}}
\begin{center}
\renewcommand{\arraystretch}{1.2}
\begin{tabular}{@{}p{0.36\linewidth} l l >{\raggedright\arraybackslash}p{0.20\linewidth}@{}}
\toprule
\textbf{Algorithm} & \textbf{Prep} & \textbf{Per op} & \textbf{Code} \\
\midrule
Build centroid tree (find, mark dead, recurse) & \(\bigO{n\log n}\) & -- & \coderef{c18-build} \\
Count pairs at distance \(=K\) / \(\le K\) & \(\bigO{n\log n}\) & one pass & \coderef{c18-count} \\
CF 715C (Digit Tree) -- decimal-hash pair count & \(\bigO{n\log n}\) & one pass & \coderef{c18-715c} \\
Centroid distance index (mark / nearest) & \(\bigO{n\log n}\) & \(\bigO{\log n}\) & \coderef{c18-index} \\
CF 293E (Close Vertices) -- 2-D dominance count & \(\bigO{n\log^2 n}\) & one pass & \coderef{c18-293e} \\
CF 833D (Red-Black Cobweb) -- product fold & \(\bigO{n\log^2 n}\) & one pass & \coderef{c18-833d} \\
Engine pick table (which tree engine) & -- & -- & \coderef{c18-map} \\
Centroid decomposition kit (one listing) & -- & -- & \coderef{c18-kit} \\
CF 757G (Summit) -- range sum of distances & \(\bigO{n\log n}\) & \(\bigO{\log^2 n}\) & \coderef{c18-summit} \\
\bottomrule
\end{tabular}
\end{center}

\chaptersection{Master Decision Flowchart}

Post-study ladder: walk top to bottom, stop at the first branch that matches. It decides
\emph{whether the problem quantifies over all paths} (centroid's home) and, if so,
\emph{which per-centroid bookkeeping} the answer needs -- and hands back every cell that
Chapters~13/17/22/25/27 already own. A \emph{single} chosen path with range updates goes to
heavy-light (\(\chref{c16}\)); a subtree set statistic goes to sack (\(\chref{c28}\)); a
changing topology goes to link-cut (\(\chref{c23}\)).

\begin{lstlisting}[language={}, basicstyle=\ttfamily\small, numbers=none, frame=single]
  Does the problem ask about PATHS in a FIXED tree?
  |
  |-- ONE chosen path that you UPDATE and QUERY (sum/min/assign on a subpath)?
  |     +-- YES: HEAVY-LIGHT + a segtree (Ch.22, c16).  NOT centroid -- centroid
  |         cannot restrict work to one path's nodes (def:c18-map, secref:c18-boundary).
  |
  |-- A SUBTREE (everything under some node), usually a SET statistic (distinct
  |   colours, k-th largest)?
  |     +-- YES: DSU-ON-TREE / sack (Ch.17, c28) or Euler-tour + Fenwick (Ch.18, c14).
  |         Sack MERGES children bottom-up; centroid SPLITS at the centre top-down.
  |
  |-- A STATE PER NODE from a LOCAL recurrence (independent set, matching)?
  |     +-- YES: TREE DP (Ch.13, c10); change the root -> REROOTING (Ch.20, c20).
  |         No distance / all-pairs dimension, so halving buys nothing.
  |
  |-- ALL simple paths (or a distance-to-a-SET), with a value ADDITIVE along the
  |   path -- THIS chapter.  First BUILD the centroid tree (24.1):
  |     decompose(entry,pc): tot=size; c=find_centroid (Ch.12 thm:c09-centfind);
  |     par_c[c]=pc; gather dist(x,c) into x's O(log n) ancestor tables; dead[c]=1;
  |     recurse on each surviving piece.  lst:c18-build.
  |
  |     |-- COUNT / SUM over pairs by DISTANCE?  (24.2)
  |     |     At each centroid claim the paths THROUGH it: build one arm's depth
  |     |     histogram, QUERY-then-MERGE so same-arm pairs are never formed here,
  |     |     then SUBTRACT each arm re-counted.  CF 161D (=K) is Ch.13's DP --
  |     |     the centroid upgrade is prose only.  lst:c18-count, lst:c18-715c.
  |     |
  |     |-- DYNAMIC nearest/farthest marked node, online, FIXED tree?  (24.3)
  |     |     One distance table per centroid; mark/query walk a node's O(log n)
  |     |     centroid ancestors taking the min.  CF 342E Xenia is Ch.12's URL.
  |     |     lst:c18-index.
  |     |
  |     |-- Pair count with a TWO-sided budget (<= L edges AND <= W weight)?  (24.3)
  |     |     Sort the (edges,weight) table on edges, sweep, Fenwick over compressed
  |     |     weight for the partner count; ordered minus self, then halve.  CF 293E.
  |     |     lst:c18-293e.
  |     |
  |     |-- Aggregate over all paths a SINGLE TAG cannot fold (product, ratio)?  (24.4)
  |     |     General plug-in template: FEATURE (additive per node), ACCEPT (2-D
  |     |     dominance), FOLD (accumulator).  CF 833D folds a product via modular
  |     |     division for the same-arm subtraction.  lst:c18-beyond, lst:c18-833d.
  |     |
  |     |-- SUM of dist(a_i, v) over a POSITION RANGE of a permutation, online?
  |            -> SUMMIT CF 757G: Fenwick per centroid over POSITIONS + the child-
  |               "sub" subtraction (all - sub[child]).  lst:c18-summit.
  |
  |-- NONE of the above -- you probably do NOT need centroid decomposition:
        a per-vertex 0/1 knapsack over a tree is NOT all-pairs; a MOVING topology is
        link-cut (Ch.25, c23); an OFFLINE batch of node-set queries is Mo (Ch.27, c21).
\end{lstlisting}

\chaptersection{Common Mistakes Checklist}

\begin{itemize}[nosep, leftmargin=1.4em]
  \item \textbf{Reusing subtree sizes across levels.} The size array is only valid inside the
        \emph{current} undead component; recompute \(\mathit{dfs\_size}\) at every
        \(\mathit{decompose}\) call with a fresh entry. A stale \(\mathit{sz}\) picks the wrong
        centroid and the height bound \(\thmref{c18-height}\) collapses to a chain. \(\coderef{c18-build}\).
  \item \textbf{Forgetting the \code{dead[]} marker.} Once a centroid is claimed it must be
        skipped by every downstream DFS. Enter a removed node and the arm collection walks
        \emph{up} into an already-processed component -- double counting or infinite recursion.
        \(\defref{c18-decomp}\), \(\coderef{c18-build}\).
  \item \textbf{Missing the same-arm subtraction.} At a centroid, a pair whose two endpoints
        live in the \emph{same} child arm is not a path through this centroid -- it belongs to a
        lower one. Count the merged histogram and \emph{subtract} each arm re-counted, or every
        same-arm pair is claimed twice. \(\defref{c18-claim}\), \(\coderef{c18-count}\).
  \item \textbf{The parent-side size \(n-\mathit{sz}[u]\) bug.} The centroid-finder walks into
        the child with \(\mathit{sz}> \mathit{tot}/2\), but the component ``above'' a candidate
        has size \(\mathit{tot}-\mathit{sz}[u]\), not \(\mathit{sz}[u]\). Ignore it and you accept
        a node that is not a centroid. Use \(\thmref{c09-centfind}\) exactly. \(\coderef{c18-build}\).
  \item \textbf{Query-then-merge order flipped.} In the histogram count you must \emph{query}
        the accumulated arms for a new arm's depths \emph{before} merging that arm in; merge
        first and a pair from the same arm matches against itself. \(\coderef{c18-count}\).
  \item \textbf{Indexing the distance table with an array \emph{position} that is not the
        node id.} The centroid tables are keyed by node; the Summit's Fenwick is keyed by
        \emph{array position} and a node moves between positions on a swap. Conflating the two
        silently reads the wrong bucket. \(\coderef{c18-summit}\).
  \item \textbf{Not dividing the ordered count by two (or forgetting the self-match).} A
        two-pointer / dominance sweep counts \((u,v)\) and \((v,u)\); the diagonal \(u=v\) can
        also sneak in when \(2a\le L\) and \(2b\le W\). Subtract self first, then halve. \(\coderef{c18-293e}\).
  \item \textbf{Product fold with integer division.} The same-arm correction on a
        \emph{multiplicative} accumulator is division modulo the prime -- use a modular inverse,
        not \code{/}. Integer division truncates and destroys the residue. \(\coderef{c18-833d}\).
  \item \textbf{Off-by-one in the centroid-depth badge.} CF 321C prints \code{'A'+depth}; if the
        top centroid is depth \(0\) the root prints \(\mathtt{A}\). Start depth at \(0\), not
        \(1\), or every letter shifts. \(\coderef{c18-321c}\).
\end{itemize}

\chaptersection{WA Clinic}

Verdicts are symptoms; the causes are the short list this chapter drilled. Diagnose top-down --
every leaf names a listing.

\begin{figure}[H]
\centering
\begin{tikzpicture}[diagram,
    v/.style={boxnode, fill=primaryblue!12, draw=primaryblue, text width=22mm, align=center},
    b/.style={boxnode, fill=difforange!12, draw=difforange, text width=12mm},
    l/.style={boxnode, fill=successgreen!10, draw=successgreen, text width=48mm,
              font=\scriptsize, align=left}]
  \node[v] (verdict) {\textbf{Verdict on the judge}};
  \node[b, below left=8mm and 5mm of verdict] (wa) {\textbf{WA}};
  \node[b, below right=8mm and 5mm of verdict] (tle) {\textbf{TLE / RE}};
  \node[l, below right=1mm and 2mm of wa.south] (wa1)
    {the same-arm pairs were never subtracted -- a pair inside one child arm is claimed at this
     centroid and again below (\(\defref{c18-claim}\), \(\coderef{c18-count}\))};
  \node[l, below left=1mm and 2mm of wa.south] (wa2)
    {an index written under the array position instead of the node id, or the ordered count
     halved before the self-match was removed (\(\coderef{c18-index}\), \(\coderef{c18-293e}\))};
  \node[l, below right=1mm and 2mm of tle.south] (tle1)
    {a stale size array picked a non-centroid: levels degrade from \(\log n\) to \(n\), turning
     \(O(n\log n)\) into \(O(n^2)\) (\(\thmref{c18-height}\), \(\coderef{c18-build}\))};
  \node[l, below left=1mm and 2mm of tle.south] (tle2)
    {reusing \code{dead[]} DFS without recomputing per piece, or a product fold using integer
     division instead of a modular inverse (\(\coderef{c18-build}\), \(\coderef{c18-833d}\))};
  \draw[flow] (verdict) -- (wa);
  \draw[flow] (verdict) -- (tle);
  \draw[flow] (wa) -- (wa1);
  \draw[flow] (wa) -- (wa2);
  \draw[flow] (tle) -- (tle1);
  \draw[flow] (tle) -- (tle2);
\end{tikzpicture}
\caption*{WA: the same-arm correction is missing, or a distance table is indexed by the wrong key
(position vs node) or halved before removing the diagonal. TLE or RE: a stale size array picked a
non-centroid so the decomposition stops halving, or the dead marker is reused instead of recomputed.}
\end{figure}

\chaptersection{Quick Reference Card}

\begin{quickref}[Chapter 24 at a glance]
\textbf{When a query quantifies over ALL paths (or a distance-to-a-set) on a fixed tree, split the
tree at its centroid, claim each path at the highest centroid on it, recurse, and merge a
per-component distance table -- \(\bigO{\log n}\) levels. }

\medskip
\textbf{Build (24.1).} \(\mathit{decompose}(\mathit{entry},\mathit{pc})\): recompute the
component size, \(\mathit{find\_centroid}\) (Ch.~12's \(\thmref{c09-centfind}\) as a black box),
set \(\mathit{par\_c}\), gather distances, mark \code{dead}, recurse on each piece. Height
\(\le\lceil\log_2 n\rceil\) (\(\defref{c18-decomp}\), \(\thmref{c18-height}\), \(\coderef{c18-build}\)).

\medskip
\textbf{Count (24.2).} Claim each path at its top centroid: collect every child arm's depths,
\emph{query} the running histogram then \emph{merge}, and \emph{subtract} each arm re-counted.
\(\bigO{n\log n}\) for \(\{=K,\le K,\)\text{divisible-by-}\(m\}\)
(\(\defref{c18-claim}\), \(\coderef{c18-count}\), \(\coderef{c18-715c}\)).

\medskip
\textbf{Index (24.3).} One distance table per centroid; a node lives in its \(\bigO{\log n}\)
centroid-ancestor entries. Mark in \(\bigO{\log n}\); nearest-marked = min over ancestors
(exact, \(\thmref{c18-nn}\)); a two-sided budget becomes 2-D dominance (sort on edges, Fenwick on
weight) (\(\defref{c18-index}\), \(\coderef{c18-index}\), \(\coderef{c18-293e}\)).

\medskip
\textbf{Beyond (24.4).} The template is three plugs: \emph{feature} additive along the arm,
\emph{accept} a dominance / threshold test, \emph{fold} an accumulator (sum, gcd, product via
modular inverse). Only the plugs change (\(\defref{c18-beyond}\), \(\coderef{c18-beyond}\), \(\coderef{c18-833d}\)).

\medskip
\textbf{The boundary (24.5).} Single path range \(\to\) HLD (\(\chref{c16}\)); subtree set
\(\to\) sack (\(\chref{c28}\)); per-node state \(\to\) DP/reroot (\(\chref{c10}\), \(\chref{c20}\));
offline node-sets \(\to\) Mo (\(\chref{c21}\)); moving edges \(\to\) link-cut (\(\chref{c23}\));
\emph{all paths} / distance-to-set \(\to\) centroid (\(\defref{c18-map}\), \(\secref{c18-boundary}\)).
\end{quickref}

\begin{lstlisting}[caption={Chapter 24 kit: the centroid find and the decompose-and-recurse build, the query-then-merge arm count, the per-centroid distance index with its nearest-marked query, the general D-C-on-trees skeleton, the Summit all-minus-sub range accumulator, and a Fenwick -- one compact skeleton each, in one place.}, label={lst:c18-kit}]
// ===== CENTROID FIND + DECOMPOSE (24.1) =====
int dfsSize(int u,int p){ sz[u]=1; for(int v:adj[u]) if(v!=p&&!dead[v]) sz[u]+=dfsSize(v,u); return sz[u]; }
int findCentroid(int u,int p,int tot){
    for(int v:adj[u]) if(v!=p&&!dead[v]&&sz[v]>tot/2) return findCentroid(v,u,tot);
    return u;                                  // no child > tot/2 and tot-sz[u] <= tot/2  (Ch.12)
}
void decompose(int entry,int pc){
    int tot=dfsSize(entry,0), c=findCentroid(entry,0,tot);
    par_c[c]=pc;                               // centroid-tree parent
    dead[c]=true;                              // CLAIM c: no downstream DFS may re-enter
    for(int v:adj[c]) if(!dead[v]) decompose(v,c);
}

// ===== COUNT pairs at distance K through the top centroid (24.2) =====
ll count(int c,int K){
    vector<int> hist(K+1,0); ll res=0;
    for(int arm: children(c)){
        vector<int> d=depths(arm,c);           // distances from c into this arm
        for(int x: d) if(x<=K) res+=hist[K-x]; // QUERY accumulated arms BEFORE merge
        for(int x: d) if(x<=K) hist[x]++;      // MERGE
    }
    return res;                                // pairs THROUGH c; same-arm handled by recursion
}

// ===== DISTANCE INDEX: nearest marked (24.3) =====
void build(int c,int entry){ dfsDist(entry,-1,0,c); }   // push (c,dist(x,c)) into anc[x]
void mark(int x){ for(auto&[c,d]:anc[x]) best[c]=min(best[c],d); }
ll query(int x){ ll r=INF; for(auto&[c,d]:anc[x]) r=min(r,best[c]+d); return r; }

// ===== GENERAL D&C-ON-TREES skeleton (24.4) =====
//   FEAT arm(x) -> feature (additive along path) ; ACCEPT(p,q) -> dominance ; FOLD(acc,val)
void solveComp(int c){
    vector<Feat> all; ll acc=NEUTRAL;
    for(int arm: children(c)){ vector<Feat> A=collect(arm,c);
        for(auto&p:A) for(auto&q:all) if(ACCEPT(p,q)) acc=FOLD(acc, combine(p,q));
        insertAll(all,A); }                    // merge arms; subtract re-counts per arm
    recurse(pieces(c));
}

// ===== SUMMIT accumulator: all - sub[child], keyed by POSITION (24.4 / 757G) =====
void addNode(int pos,int u,int sgn){ for(k=0;k<anc[u].size();k++){ int c=anc[u][k].first; ll W=anc[u][k].second;
        all[c].add(pos,sgn,sgn*W);             // store dist(u,c) at array index pos
        if(k>0) sub[anc[u][k].c].add(pos,sgn,sgn*anc[u][k-1].second); } }  // store dist(u,par_c[c])
ll queryRange(int L,int R,int v){ ll res=0; int m=anc[v].size();
    for(int k=0;k<m;k++){ int c=anc[v][k].c; ll W=anc[v][k].d;
        ll cnt=all[c].cntQ(L,R), sm=all[c].sumQ(L,R);
        if(k+1<m){ int ch=anc[v][k+1].c; cnt-=sub[ch].cntQ(L,R); sm-=sub[ch].sumQ(L,R); }
        res+=cnt*W+sm; }                        // paths through c to v, minus same-child side
    return res; }

// ===== FENWICK (Ch.18 owned; carried forward) =====
struct Fenwick{ int n; vector<ll> t; Fenwick(int n):n(n),t(n+1,0){}
    void add(int i,ll d){ for(;i<=n;i+=i&-i) t[i]+=d; }
    ll pref(int i){ ll s=0; for(;i>0;i-=i&-i) s+=t[i]; return s; }
    ll range(int l,int r){ return pref(r)-pref(l-1); } };
\end{lstlisting}

\chaptersection{Write-From-Memory Drills}

Five signature skeletons, each lifted from a listing already in the chapter. Cover the marked line,
reproduce from memory, then check the answer one line below.

\begin{lstlisting}[caption={Drill 1 -- decompose-and-recurse: find the centroid, claim it, recurse on each piece (from lst:c18-build)}]
void decompose(int entry,int pc){
    int tot=dfsSize(entry,0);                    // size of THIS undead component
    int c=findCentroid(entry,0,tot);             // Ch.12 black box (thm:c09-centfind)
    par_c[c]=pc;                                 // link into the centroid tree
    dead[c]=true;                                // CLAIM c before recursing
    for(int v:adj[c]) if(!dead[v]) decompose(v,c);
}
\end{lstlisting}
\noindent{\scriptsize Key: recompute \(\mathit{dfs\_size}\) from the \emph{entry}, mark the centroid
\emph{dead} before descending, and recurse only into still-alive neighbours -- every piece then has
\(\le n/2\) nodes, giving height \(\le\lceil\log_2 n\rceil\) (\(\thmref{c18-height}\), \(\coderef{c18-build}\)).}

\begin{lstlisting}[caption={Drill 2 -- the arm histogram count: query accumulated arms, then merge (from lst:c18-count)}]
for each child arm of centroid c:
    d = depths of nodes in this arm (distance from c)
    for x in d:  res += hist[K - x]              // QUERY first: partner is in an EARLIER arm
    for x in d:  hist[x]++                       // MERGE this arm into the histogram
\end{lstlisting}
\noindent{\scriptsize Key: \emph{query-then-merge} is what makes same-arm pairs impossible at this
centroid -- a pair whose two endpoints share an arm is claimed lower down. Query after merging would
let an arm match itself (\(\defref{c18-claim}\), \(\coderef{c18-count}\)).}

\begin{lstlisting}[caption={Drill 3 -- the distance index: mark and nearest over centroid ancestors (from lst:c18-index)}]
void mark(int x){ for((c,d) in anc[x]) best[c]=min(best[c], d); }   // d = dist(x,c)
ll nearest(int x){ ll r=INF;
    for((c,d) in anc[x]) r=min(r, best[c]+d); return r; }
\end{lstlisting}
\noindent{\scriptsize Key: a node touches only its \(\bigO{\log n}\) centroid-ancestor tables. The
min over \(\mathit{best}[c]+\dist(x,c)\) is \emph{exact} because the top shared centroid of \(x\) and
its nearest marked node is one of \(x\)'s ancestors (\(\thmref{c18-nn}\), \(\coderef{c18-index}\)).}

\begin{lstlisting}[caption={Drill 4 -- 2-D dominance pair count: sort on one axis, Fenwick the other (from lst:c18-293e)}]
sort points by a ascending                        // a = edges to the centroid
Fenwick over compressed b                         // b = weight to the centroid
for each point q in sorted order:                 // sweep the EDGE budget
    partners = number of earlier points with b <= W - q.b   // Fenwick prefix on weight
    cnt += partners                               // pairs clearing BOTH budgets with q
self = (2*q.a <= L && 2*q.b <= W) ? 1 : 0         // did q match itself?
ordered = cnt - self                              // drop the diagonal FIRST
answer  = ordered / 2                             // unordered pairs: then halve
\end{lstlisting}
\noindent{\scriptsize Key: fix one budget by the sweep order, count the other with a Fenwick; subtract
the diagonal \((u=u)\) \emph{before} halving \((u,v)/(v,u)\) -- then re-count the same arm and divide
out to claim each pair once (\(\coderef{c18-293e}\)).}

\begin{lstlisting}[caption={Drill 5 -- product fold with same-arm correction (from lst:c18-833d)}]
tot = FOLD over all arms merged                 // whole component through c
for arm: tot = tot * inverse(FOLD over that arm) // modular DIVISION removes same-arm pairs
\end{lstlisting}
\noindent{\scriptsize Key: the accumulator is a \emph{product}, so subtracting a re-counted arm is not
a minus but a multiply-by-modular-inverse \(\mathit{mod}\;\mathit{MOD}\). Integer division truncates
and wrecks the residue -- this is the classic red on CF 833D (\(\defref{c18-beyond}\), \(\coderef{c18-833d}\)).}

\chaptersection{Complete Problem Checklist}

Every problem this chapter offers -- four fully solved in-section, worked entries, and one Summit.
Tick \emph{only} after an unaided solve; a hinted attempt counts as half.

\begin{center}
\footnotesize
\setlength{\tabcolsep}{4pt}
\renewcommand{\arraystretch}{1.25}
\begin{tabular}{@{}c p{0.16\linewidth} l p{0.20\linewidth} >{\raggedright\arraybackslash}p{0.26\linewidth} c@{}}
\toprule
\textbf{\#} & \textbf{Problem} & \textbf{Source} & \textbf{Key technique} & \textbf{What to look for} & \textbf{Done?} \\
\midrule
\multicolumn{6}{l}{\textbf{Section 24.1 -- building the centroid tree}} \\
W24.1 & decompose the running tree T (15 nodes) & worked & \lv{4} find, mark dead, recurse & top centroid 4; par\_c chain; depth \(0,1,2,3\le\lceil\log_2 15\rceil\) & \(\square\) \\
W24.1b & Ciel the Commander & CF 321C & \lv{4} print \code{'A'+cd\_depth} & the \(\mathit{decompose}\) build \emph{is} the solve; rank \(0=\) A & \(\square\) \\
\midrule
\multicolumn{6}{l}{\textbf{Section 24.2 -- counting paths by distance}} \\
W24.2 & distance-2 pairs on T & worked & \lv{4} histogram, query-then-merge & 17 pairs total; top centroid 4 claims 9, recursion 8 & \(\square\) \\
W24.2b & Digit Tree & CF 715C & \lv{5} decimal-hash pair count & forward/back hash, \(\gcd(M,10)=1\) inverse; short-circuit \(M=1\) & \(\square\) \\
\midrule
\multicolumn{6}{l}{\textbf{Section 24.3 -- the centroid tree as an index}} \\
W24.3 & nearest-marked walk on T & worked & \lv{4} ancestor table min & mark \(\{5,13\}\), query 6 \(\to\) 2 (nearest marked is node 5) & \(\square\) \\
W24.3b & Close Vertices & CF 293E & \lv{5} 2-D dominance pair count & \(\le l\) edges AND \(\le w\) weight; sort+Fenwick, subtract self then halve & \(\square\) \\
\midrule
\multicolumn{6}{l}{\textbf{Section 24.4 -- beyond path sums}} \\
W24.4 & colour-balanced product on T & worked & \lv{5} feature/accept/fold & max\((r,b)\le 2\min(r,b)\) = two dominance features; arm modular-inverse & \(\square\) \\
W24.4b & Red-Black Cobweb & CF 833D & \lv{5} product fold via inverse & clamminess product; sample \(\to 1265625\) & \(\square\) \\
\midrule
\multicolumn{6}{l}{\textbf{Section 24.5 -- the engine boundary (worked only)}} \\
W24.5 & place five specs on the tree-path map & worked & \lv{5} which engine & path range/scope: HLD/sack/DP/centroid/link-cut -- pick the cell & \(\square\) \\
\midrule
\multicolumn{6}{l}{\textbf{Summit (chapter end)}} \\
W24 & Can Bash Save the Day? & CF 757G & \lv{6} centroid + Fenwick, all-sub & sum dist\((a_i,v)\) over position range under swaps; XOR-decoded; \(23,37,28\) & \(\square\) \\
\bottomrule
\end{tabular}
\end{center}

\chaptersection{Recommended Study Order}

Five focused days. Day~1 internalises the centroid-tree build and CF 321C; day~2 the
query-then-merge distance count and CF 715C; day~3 the distance index and CF 293E; day~4 the
general template and CF 833D plus the engine boundary; day~5 the Summit and the mixed set.

\begin{center}
\small
\renewcommand{\arraystretch}{1.25}
\begin{tabular}{@{}llp{0.55\linewidth}@{}}
\toprule
\textbf{Day} & \textbf{Focus} & \textbf{Tasks} \\
\midrule
1 & Read + do 24.1 & decompose T by hand to \(\mathit{par\_c}\) and depths (\(\figref{c18-centtree}\)); write \(\coderef{c18-build}\) cold; then CF 321C printing \code{'A'+depth} (\(\coderef{c18-321c}\)) \\
2 & Read + do 24.2 & hand-count distance-2 pairs on T (\(17=9+8\)); name the query-then-merge and same-arm steps (\(\coderef{c18-count}\)); then CF 715C with the forward/backward decimal hash (\(\coderef{c18-715c}\)) \\
3 & Read + do 24.3 & trace the nearest-marked walk (mark \(\{5,13\}\), query 6 \(\to\) 2) (\(\coderef{c18-index}\)); write the ancestor index cold; then CF 293E's sort-plus-Fenwick dominance count (\(\coderef{c18-293e}\)) \\
4 & Read + do 24.4 + 24.5 & name the three plugs (feature/accept/fold) on CF 833D and replay sample 1 to \(1265625\) (\(\coderef{c18-833d}\)); place five specs on the engine map (\(\defref{c18-map}\)) \\
5 & Summit + mixed & derive the all-minus-sub accumulator and replay the decoded samples \(23,37,28\) (\(\probref{c18-summit}\)); then the first 3 mixed entries \\
\bottomrule
\end{tabular}
\end{center}

\chaptersection{Difficulty Progression}

Tiers are competency gates -- each row names the \emph{skill} the tier unlocks.

\begin{center}
\small
\renewcommand{\arraystretch}{1.3}
\begin{tabular}{@{}lp{0.32\linewidth}p{0.40\linewidth}@{}}
\toprule
\textbf{Tier} & \textbf{Problems} & \textbf{What passing this tier proves} \\
\midrule
\lv{3} \ \(\bigstar\bigstar\bigstar\)
  & the \(\mathit{find\_centroid}\) walk; a single arm's depth collection
  & you recompute sizes per component and pick the true balance point by hand \\
\lv{4} \ \(\bigstar\bigstar\bigstar\bigstar\)
  & CF 321C; the distance-\(K\) histogram on T; the nearest-marked trace
  & you build the centroid tree, keep the query-then-merge order, and walk ancestor tables \\
\lv{5} \ \(\bigstar\bigstar\bigstar\bigstar\bigstar\)
  & CF 715C, CF 293E, CF 833D, mixed at CF~2400--2900
  & you bolt an ordered structure (hash, Fenwick, product) onto each centroid bucket \\
\lv{6} \ \(\bigstar\bigstar\bigstar\bigstar\bigstar\bigstar\)
  & CF 757G (Summit), the CF~3100--3400 mixed set
  & you keep \emph{two} accumulators (all and sub) keyed by position and fold them correctly \\
\bottomrule
\end{tabular}
\end{center}

\chaptersection{Reduction Gallery}

Every entry is ``one all-paths shape \(\to\) one centroid move.'' Each line collapses a problem
family onto a listing already in the chapter.

\begin{itemize}[nosep, leftmargin=1.4em]
  \item ``output the centroid hierarchy / rank every node by the level that removed it'' (CF 321C)
        \(\to\) the build \emph{is} the solve: \(\mathit{decompose}\) writes \(\mathit{par\_c}\) and
        depth, print \code{'A'+depth} (\(\coderef{c18-build}\), \(\coderef{c18-321c}\)).
  \item ``count unordered pairs at distance \(=K\) / \(\le K\)'' (CF 161D, Ch.~13's DP upgraded) \(\to\)
        claim each pair at its top centroid, query-then-merge a depth histogram, subtract the same
        arm (\(\defref{c18-claim}\), \(\coderef{c18-count}\)).
  \item ``count pairs whose path-as-a-number is divisible by \(m\)'' (CF 715C) \(\to\) the same
        centroid count with a forward/backward decimal hash and \(\gcd(m,10)=1\) inverse
        (\(\coderef{c18-715c}\)).
  \item ``distance-to-a-set / nearest marked node, online, fixed tree'' (CF 342E Xenia, Ch.~12's URL)
        \(\to\) one distance table per centroid, min over a node's \(\bigO{\log n}\) ancestors
        (\(\thmref{c18-nn}\), \(\coderef{c18-index}\)).
  \item ``pairs within TWO budgets (\(\le L\) edges AND \(\le W\) weight)'' (CF 293E) \(\to\) 2-D
        dominance in the bucket: sort on edges, Fenwick over compressed weight, subtract self then
        halve (\(\coderef{c18-293e}\)).
  \item ``aggregate over all paths a single tag cannot fold (product, gcd, ratio)'' (CF 833D) \(\to\)
        feature/accept/fold plugs; the same-arm correction is a modular inverse on a product
        (\(\defref{c18-beyond}\), \(\coderef{c18-833d}\)).
  \item ``sum of distances over a \emph{position range} of a permutation, under swaps'' (CF 757G,
        Summit) \(\to\) a Fenwick per centroid keyed by position, with the child \emph{sub}
        table so each pair is claimed once (\(\coderef{c18-summit}\)).
  \item ``is this even a centroid job?'' \(\to\) check the engine map: single path range \(\to\) HLD,
        subtree set \(\to\) sack, per-node state \(\to\) DP, moving edges \(\to\) link-cut
        (\(\defref{c18-map}\), \(\secref{c18-boundary}\)).
\end{itemize}

\chaptersection{Mixed Problem Set}

Ten entries, each live-verified against the current problem sets (Codeforces tags
\emph{divide and conquer} plus \emph{trees}; CF has no dedicated centroid tag) and absent from
every other roster in the book (the owned IDs -- CF 342E in Ch.~12, CF 161D in Ch.~13, CF
914D/1601C/1743F, and this chapter's own solved 321C / 715C / 293E / 833D -- are excluded). None
says ``use centroid decomposition'' -- you recognise it from \emph{what the query quantifies over}
(all paths, or a distance-to-a-set) and \emph{whether the topology is fixed}. Each hint names its
reduction and the Chapter~24 listing it leans on, ordered ascending by effective difficulty.

\begin{mixedset}
  \practice{Fib-tree}{CF 1491E}
    {https://codeforces.com/problemset/problem/1491/E}{CF 2400 L5}
    {Cut a tree into two connected parts repeatedly so the part sizes follow a Fibonacci-like
     rule. Each cut is ``split the current component at a node,'' so the natural candidate is the
     balance split of \(\secref{c18-tree}\): the centroid is a node whose removal leaves pieces of
     size \(\le n/2\), and a greedy centroid-first cut is exactly the move a Fibonacci decomposition
     needs. Reuse \(\coderef{c18-build}\)'s recompute-sizes-per-component discipline; the difference
     is that you cut at a \emph{node} (edge removal), validating sizes, not merely ranking.}
  \practice{New Year Tree}{CF 379F}
    {https://codeforces.com/problemset/problem/379/F}{CF 2400 L5}
    {The tree only \emph{grows} -- you attach a new leaf -- and queries ask for the distance between
     two existing nodes. A growing tree is a fixed tree at every instant, so \(\dist(u,v)=
     \mathit{depth}_w[u]+\mathit{depth}_w[v]-2\mathit{depth}_w[\mathrm{lca}]\) with binary lifting is
     the cheapest route; but see why centroid decomposition is the WRONG pick here
     (\(\defref{c18-map}\)): the centroid tree would have to be rebuilt after every attachment, so a
     one-shot LCA wins. Recognise the boundary, not the engine.}
  \practice{Magic Matrix}{CF 632F}
    {https://codeforces.com/problemset/problem/632/F}{CF 2400 L5}
    {A ``magic matrix'' is defined by a nearest-neighbour property you check with a \emph{divide and
     conquer on the tree}: root arbitrarily and verify, for each node, that the closest marked node
     can be found by recursion. That is the distance-index pattern -- a per-centroid nearest
     structure (\(\thmref{c18-nn}\), \(\coderef{c18-index}\)) answers each ``who is nearest?'' check
     in \(\bigO{\log n}\). Port the \(\mathit{best}[c]+d\) min from the index listing verbatim.}
  \practice{Palindromes in a Tree}{CF 914E}
    {https://codeforces.com/problemset/problem/914/E}{CF 2400 L5}
    {For every node, count the paths \emph{through} it whose label multiset is palindrome-able (at
     most one odd colour). The path is claimed at its top centroid, and you merge per-arm colour
     parities with the query-then-merge histogram of \(\coderef{c18-count}\): parity is the
     \emph{feature}, the at-most-one-odd rule is the \emph{accept}, and the count is the
     \emph{fold}. The same-arm subtraction is the one thing that keeps it correct.}
  \practice{GCD Counting}{CF 990G}
    {https://codeforces.com/problemset/problem/990/G}{CF 2400 L5}
    {Count paths whose node-value gcd is at least \(2\), by prime factor. Fix a prime \(p\) and keep
     the component of \(p\)-divisible nodes -- inside it, count connected paths, a task that a
     centroid pass over each surviving component handles with the arm merge of
     \(\coderef{c18-count}\). Mobius inversion cleans the double-count across primes; the centroid
     machinery is the ``all paths, additive along path'' engine
     (\(\defref{c18-beyond}\)).}
  \practice{Vertex Pairs}{CF 2042E}
    {https://codeforces.com/problemset/problem/2042/E}{CF 2900 L6}
    {Count pairs \((u,v)\) whose \emph{maximum} edge on the path is at most a per-node threshold --
     a two-sided condition on the path exactly like CF 293E. Claim each pair at its top centroid,
     collect \((\mathit{maxEdge}, \mathit{depth})\) features per arm, and count with the sort-plus-Fenwick
     dominance sweep of \(\coderef{c18-293e}\). The reduction is ``within a budget over all
     pairs \(\to\) 2-D dominance in the centroid bucket.''}
  \practice{Logistical Questions}{CF 566C}
    {https://codeforces.com/problemset/problem/566/C}{CF 3000 L6}
    {Find the node minimising a weighted sum of distances to all others. That node is a
     \emph{weighted centroid} (the balance point under weights), which \(\secref{c18-tree}\) and
     Ch.~12's \(\thmref{c09-centfind}\) locate by walking to the heavy side -- here, the side whose
     weighted mass exceeds half. Binary-search the answer value using the distance-table walk of
     \(\coderef{c18-index}\) to check a candidate. Centroid \emph{is} the objective, not just the tool.}
  \practice{Tavas on the Path}{CF 536E}
    {https://codeforces.com/problemset/problem/536/E}{CF 3100 L6}
    {A fixed tree with path queries whose answers must be precomputed efficiently -- but read the
     engine map first: a \emph{single} queried path with an associative combine is a heavy-light or
     static-path problem (\(\chref{c16}\)), \emph{not} an all-paths problem. Use this entry to
     practise REJECTING centroid: the query ranges over one chosen path, so \(\defref{c18-map}\)
     hands it to Ch.~22. Contrast with \(\secref{c18-count}\).}
  \practice{Statistics on Tree}{CF 2222G}
    {https://codeforces.com/problemset/problem/2222/G}{CF 3100 L6}
    {An aggregate over all root-to-leaf paths of a tree, where each path is really the unique path to
     a query node -- i.e.\ an accumulation that a per-node DP cannot scope to \emph{every} path at
     once. Claim each path at its top centroid and fold with the feature/accept/fold template of
     \(\coderef{c18-beyond}\); the bucket aggregate is a plain sum (a composable fold), so no ordered
     structure is needed -- the win is the \(\bigO{\log n}\)-level coverage of all paths.}
  \practice{Iqea}{CF 936E}
    {https://codeforces.com/problemset/problem/936/E}{CF 3400 L6}
    {The hardest all-paths entry: aggregate a value over every simple path with a filter that is a
     dominance/threshold on an additive feature -- the exact home of centroid decomposition. Run
     \(\coderef{c18-beyond}\) with the pair of additive features that encode the constraint (as CF
     833D used \(2r-b,2b-r\)), a Fenwick over the compressed second feature inside each bucket
     (\(\coderef{c18-293e}\)), and the same-arm correction. Everything from 24.2--24.4 composes here.}
\end{mixedset}


\chaptersection{The Summit -- CF 757G Can Bash Save the Day?}

\begin{notebox}[Why this is the chapter top rope]
Every in-section engine here answers one kind of question about a \emph{set of paths}; the Summit
makes you answer a question about a \emph{set of array positions} whose node contents \emph{move},
and do it online. You need \emph{two} of the chapter's ideas at once: the per-centroid distance
table that lets a query walk a node's \(\bigO{\log n}\) ancestors (\(\secref{c18-index}\)), and the
top-centroid claim that lets you \emph{subtract} the child-side so each path is counted once
(\(\defref{c18-claim}\)). Bolt a Fenwick / segment tree keyed by \emph{position} onto each centroid
bucket (Chapter~18's \(\chref{c14}\) carried forward) and hold two accumulators per centroid -- the
\texttt{all} table and the child \texttt{sub} table. Adjacent swaps are then just delete-and-reinsert
a node's \(\bigO{\log n}\) entries. It is the honest summit rather than the highest badge: it is the
hardest problem this chapter \emph{fully derives} and hand-replays to the judge's output.
\end{notebox}

\solvedproblem{c18-summit}{Can Bash Save the Day?}{CF 757G}
  {https://codeforces.com/problemset/problem/757/G}{CF 3400 L6}
  {Centroid decomposition with a Fenwick or segment tree per centroid over array positions storing each node's distance to the centroid, plus a child-side subtraction table; a range-sum-of-distances query walks the query node's O(log n) centroid ancestors combining all minus sub, and an adjacent swap re-inserts two nodes' ancestor entries; queries are XOR-encoded so the structure stays online}

\begin{problemstatement}[CF 757G Can Bash Save the Day?]
You are given a weighted tree on \(n\) nodes (edge lengths are positive integers up to
\(10^6\)) and a permutation \(a_1,\dots,a_n\) of \(1,\dots,n\). Let \(\dist(u,v)\) be the total
edge weight on the unique path between \(u\) and \(v\). Answer \(q\) online queries of two types:

\textbf{Type 1,} written \texttt{1 l r v}: report \(\sum_{i=l}^{r} \dist(a_i, v)\) -- the sum of
distances from node \(v\) to the nodes occupying array positions \(l, l+1, \dots, r\).\\
\textbf{Type 2,} written \texttt{2 x}: swap \(a_x\) and \(a_{x+1}\); the new permutation is used by
all later queries.

The queries are \textbf{XOR-encoded} against the previous answer (so you cannot reorder them). Read
the raw integers \(a,b,c\) for a type-1 query and \(a\) for a type-2 query, and decode with the
running answer \(\mathit{last}\) (initialised to \(0\), and \emph{unchanged} by a type-2 query):
\[
  l = a \oplus \mathit{last},\quad r = b \oplus \mathit{last},\quad v = c \oplus \mathit{last}
  \qquad(\text{type }1),\qquad x = a \oplus \mathit{last}\qquad(\text{type }2),
\]
where \(\oplus\) is bitwise exclusive-or. The guarantees hold after decoding:
\(1\le l\le r\le n\), \(1\le v\le n\), \(1\le x\le n-1\).

\textbf{Input.} Line~1: \(n, q\) (\(1\le n,q\le 2\cdot10^5\)). Line~2: the permutation
\(a_1,\dots,a_n\). Next \(n-1\) lines: \(u, v, w\), an edge of length \(w\) (\(1\le w\le 10^6\)).
Then \(q\) blocks: a line \texttt{t}, and a line of the (encoded) query parameters.

\textbf{Output.} For each type-1 query, print the sum on its own line.

\textbf{Sample.} \(n{=}5,q{=}5\), \(a=(4,5,1,3,2)\), edges \((4,2,4),(1,3,9),(4,1,4),(4,5,2)\);
after XOR-decoding the queries are \(\{1\ 1\ 5\ 4\}, \{1\ 1\ 3\ 3\}, \{2\ 3\}, \{2\ 2\},
\{1\ 1\ 3\ 3\}\), and the three type-1 answers are \texttt{23}, \texttt{37}, \texttt{28}.
\end{problemstatement}

\begin{attemptgate}[Attempt it yourself first -- \lv{6}, budget a real afternoon]
Summing \(\dist(a_i,v)\) one position at a time is \(\bigO{n}\) per query, i.e.\ \(\bigO{nq}\)
overall -- \tle at \(2\cdot10^5\), and each \(\dist\) would itself need an LCA. The XOR encoding
means you cannot batch offline.
\peek{a distance \(\dist(u,v)\) splits cleanly once you fix a centroid that lies \emph{on} the path
\(u\)--\(v\) -- which centroid is that, and what does \(\dist(u,v)\) become?}
\nudge{Take \(c\) to be the top shared centroid of \(u\) and \(v\): then
\(\dist(u,v)=\dist(u,c)+\dist(v,c)\). So \(\sum_{i\in[l,r]}\dist(a_i,v)=\sum_c\big(\dist(v,c)\cdot
\#\{i:[l,r],\,c\text{ top for }a_i\}+\sum \dist(a_i,c)\big)\) -- but you must claim each \(a_i\) at
\emph{exactly one} centroid \(c\).}
\fullhint{For each centroid \(c\) keep two structures keyed by array \emph{position}: \texttt{all[c]}
holds \(\dist(a_i,c)\) and a count for every position whose node lies in \(c\)'s component, and
\texttt{sub[c]} holds \(\dist(a_i,\mathit{par\_c}[c])\) for those same positions. To answer a query
for \(v\), walk \(v\)'s centroid ancestors \(c\) (root toward \(v\)); at level \(c\) the child \(d\)
of \(c\) on the path to \(v\) owns exactly the positions whose top shared centroid with \(v\) is
\emph{below} \(c\), so subtract \texttt{sub[d]}: the contribution is
\((\mathit{cnt}_{\mathit{all}[c]}-\mathit{cnt}_{\mathit{sub}[d]})\cdot\dist(v,c)+
(\mathit{sum}_{\mathit{all}[c]}-\mathit{sum}_{\mathit{sub}[d]})\). Each position lives in
\(\bigO{\log n}\) centroid tables, so a swap re-inserts two nodes in \(\bigO{\log^2 n}\) and a query
is \(\bigO{\log^2 n}\). A per-centroid Fenwick or dynamic segment tree over positions supports the
point updates and range reads; use 64-bit -- sums reach \(O(n\cdot 10^{11})\).}
\end{attemptgate}

\solstep{Understanding}
Fix a query \((l,r,v)\). For each position \(i\in[l,r]\) we must add \(\dist(a_i,v)\). The key
identity (\(\defref{c18-claim}\)) is that every node pair has a \emph{unique top shared centroid}
\(c\) lying on the path between them, and for that \(c\), \(\dist(a_i,v)=\dist(a_i,c)+\dist(v,c)\).
Group the range's positions by that top centroid \(c\). As \(c\) ranges over \(v\)'s own centroid
ancestors, the positions whose top shared centroid with \(v\) is exactly \(c\) are: those in \(c\)'s
component \emph{minus} those in the child-of-\(c\) component that also contains \(v\). For those
positions the distance-to-\(c\) values are exactly what a table at \(c\) stores, and \(\dist(v,c)\) is
a constant of the walk. So the whole sum is a difference of two per-centroid aggregates: the
\texttt{all[c]} table minus the \texttt{sub[d]} table of the child \(d\) toward \(v\). Positions move
only by adjacent swaps, which are just delete-and-reinsert on \(\bigO{\log n}\) centroid tables per
node.

\solstep{Key Observations}
\begin{observation}
Each node \(x\) is stored in exactly its \(\bigO{\log n}\) centroid-ancestor tables (the identity
behind \(\thmref{c18-nn}\)): for every centroid \(c\) on \(x\)'s path to the centroid-tree root,
\(x\in\texttt{all[c]}\) with value \(\dist(x,c)\), and \(x\in\texttt{sub[c]}\) with value
\(\dist(x,\mathit{par\_c}[c])\). That bound is what makes both updates and queries polylogarithmic.
\end{observation}
\begin{observation}
When walking \(v\)'s ancestors root\(\to v\), the ``child toward \(v\)'' at level \(c\) is simply the
\emph{next} centroid in the walk. Subtracting \(\texttt{sub}[\text{next}]\) removes precisely the
positions that will be (and must be) claimed lower down, so no pair is double-counted -- this is the
same-arm correction of \(\secref{c18-count}\), re-derived for a position range.
\end{observation}
\begin{observation}
The array index (position) is the query's range axis, but the tree node is the membership axis. The
two are linked by the permutation \(a\), which a swap changes at two positions -- so a node's
contribution must be \emph{moved} between positions, not between centroid tables (whose contents are
topological and never change).
\end{observation}

\solstep{Think It With Me}
\begin{thinkbox}
\deadend{Answer each type-1 query by looping \(i=l..r\) and summing \(\dist(a_i,v)\) via LCA}
{That is \(\bigO{n}\) per query, \(\bigO{nq}\) total -- \tle at \(2\cdot10^5\). A single \(\dist\) is
cheap, but the \emph{range sum} of \(\bigO{n}\) of them is the wall. You need the range aggregate to
be answerable without touching every position.}
\deadend{One global Fenwick over positions storing \(\dist(a_i,v)\)}
{The stored value would have to depend on the queried \(v\), which changes every query -- you cannot
precompute \(\dist(a_i,v)\) for a \(v\) you have not seen. The distance must be \emph{decomposed}
into per-centroid pieces indexed only by position.}
\deadend{Centroid \texttt{all[c]} tables but no \texttt{sub} table}
{Summing \(\texttt{all[c]}\) over \(v\)'s ancestors counts a position \(i\) once for \emph{every}
centroid that is an ancestor of both \(a_i\) and \(v\), not once at their top shared centroid --
massive over-count, \textbf{WA}. The child \texttt{sub[d]} subtraction is exactly what restricts the
claim to one centroid (\(\defref{c18-claim}\)).}
\aha{Keep \emph{two} position-indexed accumulators per centroid -- \texttt{all[c]} with
\(\dist(x,c)\) and \texttt{sub[c]} with \(\dist(x,\mathit{par\_c}[c])\) -- and answer a query by
walking \(v\)'s centroid ancestors combining \(\texttt{all[c]}-\texttt{sub[child]}\). Adjacent swaps
delete and re-insert a node's \(\bigO{\log n}\) entries. Two chapter engines (the distance index and
the top-centroid claim) plus the Fenwick compose into one \(\bigO{\log^2 n}\) query.}
\end{thinkbox}

\solstep{The Approach Ladder}
\begin{center}\renewcommand{\arraystretch}{1.25}
\begin{tabular}{@{}c p{0.24\linewidth} p{0.30\linewidth} l c@{}}
\toprule
\textbf{\#} & \textbf{Approach} & \textbf{Idea} & \textbf{Time} & \textbf{Verdict}\\
\midrule
1 & \lv{2}\ loop the range, LCA each & sum \(\dist(a_i,v)\) position by position & \(\bigO{nq\log n}\) & \tle \\
2 & \lv{4}\ global position Fenwick & one table of distances & \(\bigO{\log n}\) & \textcolor{warningred}{\ding{55}\ \textbf{WA}} \\
3 & \lv{5}\ centroid \texttt{all} only & per-centroid tables, no child subtraction & \(\bigO{\log^2 n}\) & \textcolor{warningred}{\ding{55}\ \textbf{WA}} \\
4 & \lv{6}\ \textbf{centroid all \(−\) sub, by position} & walk \(v\)'s ancestors, \(\texttt{all[c]}-\texttt{sub[child]}\) & \(\bigO{\log^2 n}\) & \ac \\
\bottomrule
\end{tabular}\end{center}
Rung~4 is the Summit: rung~3 has the right tables but the wrong \emph{scope} -- without the child
\(\texttt{sub}\) subtraction a position is claimed at several centroids at once, and the over-count
is the whole error.

\solstep{Solution}
Build the centroid decomposition once (\(\coderef{c18-build}\)): as each centroid \(c\) is chosen,
gather \(\dist(x,c)\) into every node \(x\) of its component, appending \((c,\dist(x,c))\) to
\(\mathit{anc}[x]\) in root\(\to\)leaf order. Insert every position \(i\) into the structures for its
node \(a_i\): \(\texttt{all}[c]\) gets \(\dist(a_i,c)\) and \(\texttt{sub}[c]\) gets
\(\dist(a_i,\mathit{par\_c}[c])\). A type-1 query walks \(\mathit{anc}[v]\), combining
\(\texttt{all}[c]-\texttt{sub}[\text{next}]\) over the position range. A type-2 swap deletes the two
nodes at their old positions and re-inserts them at the swapped positions. Read and write 64-bit
throughout.

\solstep{Proof of Correctness}
\begin{proofbox}
\textbf{Distance decomposition.} For any nodes \(u,v\), let \(c^\star\) be their lowest common node
in the centroid tree (equivalently the first centroid whose removal separated them, or one of them).
Then \(c^\star\) lies on the path \(u\)--\(v\), so \(\dist(u,v)=\dist(u,c^\star)+\dist(v,c^\star)\)
(\(\defref{c18-claim}\)). \textbf{Single claim.} Fix \(v\). As \(c\) ranges over \(\mathit{anc}[v]\)
(root\(\to v\)) with \(d\) the next centroid toward \(v\), the set \(\{\text{nodes in }c\text{'s
component}\}\setminus\{\text{nodes in }d\text{'s component}\}\) is exactly the nodes whose \(c^\star\)
with \(v\) equals \(c\): \(u\) is claimed at \(c\) iff \(u\) shares the ancestor \(c\) with \(v\) but
does not fall into the same child branch \(d\) (else its \(c^\star\) is lower). These differences
partition all nodes, so each position \(i\) is counted \emph{once}, at its true \(c^\star\).
\textbf{Value read.} At that level the algorithm adds
\((\mathit{cnt}_{\mathit{all}[c]}-\mathit{cnt}_{\mathit{sub}[d]})\cdot\dist(v,c)+
(\mathit{sum}_{\mathit{all}[c]}-\mathit{sum}_{\mathit{sub}[d]})
=\sum_{i\in[l,r],\,c^\star(i,v)=c}(\dist(v,c)+\dist(a_i,c))
=\sum \dist(a_i,v)\) by the decomposition. \textbf{Updates.} A swap only permutes which node sits at
which position; deleting the node's \(\bigO{\log n}\) entries and re-inserting them at the new
position keeps every \(\texttt{all}/\texttt{sub}\) table equal to the current position-assignment, so
every later range read is exact. Summing over levels gives \(\sum_{i=l}^r\dist(a_i,v)\). \(\Box\)
\end{proofbox}

\solstep{Implementation}
\begin{lstlisting}[caption={CF 757G Summit: centroid decomposition; per centroid, two dynamic segment trees over POSITIONS -- all[c] stores dist(node,c) and sub[c] stores dist(node,parc[c]). A query for v walks v's centroid ancestors combining all minus the child's sub; a swap deletes and re-inserts two nodes' O(log n) entries. XOR-decode against the last answer. ASCII only}, label={lst:c18-summit}]
#include <bits/stdc++.h>
using namespace std;
typedef long long ll;
const int MAXN = 200005;
const int MAXP = 8000000;              // segment-tree pool, O(n log^2 n) nodes
int n, q;
vector<pair<int,int> > adj[MAXN];
int a[MAXN];                            // a[position] = node currently there
bool dead[MAXN]; int sz[MAXN], parc[MAXN];
vector<pair<int,ll> > anc[MAXN];        // (centroid c, dist(node,c)) in ROOT->node order
int dfsSize(int u,int p){ sz[u]=1; for(size_t k=0;k<adj[u].size();k++){ int v=adj[u][k].first;
    if(v!=p&&!dead[v]) sz[u]+=dfsSize(v,u);} return sz[u]; }
int findCent(int u,int p,int tot){ for(size_t k=0;k<adj[u].size();k++){ int v=adj[u][k].first;
    if(v!=p&&!dead[v]&&sz[v]>tot/2) return findCent(v,u,tot);} return u; }   // Ch.12 black box
void gather(int u,int p,ll d,int c){ anc[u].push_back(make_pair(c,d));
    for(size_t k=0;k<adj[u].size();k++){ int v=adj[u][k].first, w=adj[u][k].second;
        if(v!=p&&!dead[v]) gather(v,u,d+w,c);} }
void decompose(int entry,int pc){ int tot=dfsSize(entry,0); int c=findCent(entry,0,tot);
    parc[c]=pc; gather(c,0,0,c); dead[c]=true;                 // claim c, record dist(x,c)
    for(size_t k=0;k<adj[c].size();k++){ int v=adj[c][k].first; if(!dead[v]) decompose(v,c);} }
// ---- dynamic segment tree over positions, storing (count, sum of distances) ----
struct Seg{ int l,r; ll cnt,sum; } pool[MAXP]; int Npool=0;
int allR[MAXN], subR[MAXN];
int newNode(){ ++Npool; pool[Npool].l=pool[Npool].r=0; pool[Npool].cnt=pool[Npool].sum=0; return Npool; }
void upd(int &t,int lo,int hi,int pos,ll dc,ll ds){ if(!t) t=newNode();
    pool[t].cnt+=dc; pool[t].sum+=ds; if(lo==hi) return; int mid=(lo+hi)>>1;
    if(pos<=mid) upd(pool[t].l,lo,mid,pos,dc,ds); else upd(pool[t].r,mid+1,hi,pos,dc,ds); }
pair<ll,ll> qry(int t,int lo,int hi,int L,int R){ if(!t||R<lo||hi<L) return make_pair(0LL,0LL);
    if(L<=lo&&hi<=R) return make_pair(pool[t].cnt,pool[t].sum); int mid=(lo+hi)>>1;
    pair<ll,ll> A=qry(pool[t].l,lo,mid,L,R), B=qry(pool[t].r,mid+1,hi,L,R);
    return make_pair(A.first+B.first,A.second+B.second); }
void addNode(int u,int pos,int sgn){                          // insert (sgn=+1) or delete (-1)
    for(size_t k=0;k<anc[u].size();k++){ int c=anc[u][k].first; ll W=anc[u][k].second;
        upd(allR[c],1,n,pos,sgn,(ll)sgn*W);                  // all[c] : dist(u,c) at position pos
        if(k>0) upd(subR[c],1,n,pos,sgn,(ll)sgn*anc[u][k-1].second); } } // sub[c] : dist(u,parc[c])
ll queryRange(int L,int R,int v){ ll res=0; int m=(int)anc[v].size();
    for(int k=0;k<m;k++){ int c=anc[v][k].first; ll W=anc[v][k].second;   // walk v's centroid ancestors
        pair<ll,ll> AC=qry(allR[c],1,n,L,R); ll cnt=AC.first, sm=AC.second;
        if(k+1<m){ int ch=anc[v][k+1].first;                          // child of c toward v
            pair<ll,ll> SD=qry(subR[ch],1,n,L,R); cnt-=SD.first; sm-=SD.second; }
        res+=cnt*W+sm; }                                              // claim paths at top centroid c
    return res; }
int main(){ scanf("%d%d",&n,&q); for(int i=1;i<=n;i++) scanf("%d",&a[i]);
    for(int i=1;i<n;i++){ int u,v,w; scanf("%d%d%d",&u,&v,&w);
        adj[u].push_back(make_pair(v,w)); adj[v].push_back(make_pair(u,w)); }
    decompose(1,0);
    for(int i=1;i<=n;i++) addNode(a[i],i,1);                 // seed position i with node a[i]
    ll last=0;
    for(int qi=0;qi<q;qi++){ int t; scanf("%d",&t);
        if(t==1){ ll A,B,C; scanf("%lld%lld%lld",&A,&B,&C);
            int l=(int)(A^last), r=(int)(B^last), v=(int)(C^last); if(l>r) swap(l,r);
            ll ans=queryRange(l,r,v); printf("%lld\n",ans); last=ans; }
        else { ll A; scanf("%lld",&A); int x=(int)(A^last); int u=a[x], w=a[x+1];
            addNode(u,x,-1); addNode(w,x,1); addNode(w,x+1,-1); addNode(u,x+1,1);
            a[x]=w; a[x+1]=u; } }
    return 0; }
// Two accumulators, one walk. all[c] keys dist(node,c) by ARRAY POSITION; sub[c] keys
// dist(node,parc[c]) so the child-side (claimed lower) can be subtracted. O(log n) tables per node
// => O(log^2 n) per query/swap. 64-bit: a distance sum can reach ~n*1e11, still inside long long.
\end{lstlisting}

\solstep{Trace on the Official Sample}
\begin{tracebox}[Sample: decompose the 5-node tree, then replay the three decoded type-1 queries]
Centroid of \(\{1,2,3,4,5\}\) is \(4\) (pieces \(\{2\},\{5\},\{1,3\}\)); recurse gives par\_c:
\(4\!\to\!0\), \(2,5,1\!\to\!4\), \(3\!\to\!1\). Ancestor tables (root\(\to\)leaf), each entry a
centroid and the distance to it: \(\mathit{anc}[4]=[(4,0)]\); \(\mathit{anc}[2]=[(4,4),(2,0)]\);
\(\mathit{anc}[5]=[(4,2),(5,0)]\); \(\mathit{anc}[1]=[(4,4),(1,0)]\);
\(\mathit{anc}[3]=[(4,13),(1,9),(3,0)]\). The array is \(a=(4,5,1,3,2)\) at positions \(1..5\).

\textbf{Query 1: \(l=1,r=5,v=4\).} Only level \(c=4\), \(W=0\), no child: \(\texttt{all}[4]\) spans
the whole tree so the range sum is \(\sum_i\dist(a_i,4)=0+2+4+13+4=23\), count \(5\), giving
\(5\cdot0+23=23\). \(\boxed{23}\) \(\checkmark\)

\textbf{Query 2: \(l=1,r=3,v=3\).} Positions \(1..3\) hold nodes \(4,5,1\).
Level \(c=4\), \(W=13\), child \(d=1\): \(\texttt{all}[4]\) over \([1,3]\) is count \(3\), sum
\(0+2+4=6\); \(\texttt{sub}[1]\) (nodes of component \(\{1,3\}\) stored with distance to \(4\)) over
\([1,3]\) catches node \(1\) at position \(3\): count \(1\), sum \(4\). Contribution
\((3-1)\cdot13+(6-4)=26+2=28\). Level \(c=1\), \(W=9\), child \(d=3\): \(\texttt{all}[1]\) over
\([1,3]\) is node \(1\) at position \(3\): count \(1\), sum \(\dist(1,1)=0\); \(\texttt{sub}[3]\)
(node \(3\), position \(4\)) is empty. Contribution \((1-0)\cdot9+(0-0)=9\). Level \(c=3\):
component \(\{3\}\) sits at position \(4\), outside the range, contributes \(0\). Total \(28+9=37\).
\(\boxed{37}\) \(\checkmark\)

\textbf{After the two swaps} (\(x=3\), then \(x=2\)) the array is \(a=(4,3,5,1,2)\). \textbf{Query 3:
\(l=1,r=3,v=3\)} now reads nodes \(4,3,5\): \(13+0+15=\boxed{28}\) \(\checkmark\) -- matching the
judge exactly.
\end{tracebox}

\solstep{Complexity}
\begin{complexitybox}
Decomposition is \(\bigO{n\log n}\) time (\(\thmref{c18-height}\)), storing \(\bigO{\log n}\)
ancestor entries per node. A type-1 query walks \(v\)'s \(\bigO{\log n}\) ancestors, doing two
range-reads (\(\bigO{\log n}\) each on the position segment tree) per level: \(\bigO{\log^2 n}\). A
type-2 swap re-inserts two nodes across \(\bigO{\log n}\) tables each, also \(\bigO{\log^2 n}\).
Total \(\bigO{(n+q)\log^2 n}\approx 4\cdot10^{5}\cdot 18^2 \approx 1.3\cdot10^{8}\) node touches --
AC in 5s with the small \(\log\) constant. Memory: \(\bigO{n\log^2 n}\) segment-tree nodes (the
all and sub trees) -- the reason the limit is a generous 768 MB; an implementation that shares
structure or uses per-centroid Fenwicks of touched positions is leaner. Use \texttt{long long} for
answers: a distance sum reaches \(O(n\cdot\!\max\!\dist)\approx 2\cdot10^{5}\cdot2\cdot10^{11}\) in
the theoretical worst case, though the test data keeps decoded answers inside \(2^{30}\).
\end{complexitybox}

\begin{hardeningbox}
\begin{itemize}[nosep,leftmargin=1.4em]
  \item \textbf{Never forget the \texttt{sub[child]} subtraction.} Summing \(\texttt{all[c]}\) over
        \(v\)'s ancestors claims a position at \emph{every} shared centroid, not just its top one --
        the over-count is the classic \textbf{WA}. The child table removes exactly the positions that
        belong to a lower claim (\(\defref{c18-claim}\)).
  \item \textbf{Index by POSITION, key distance by the table's centroid.} \(\texttt{all[c]}\) stores
        \(\dist(x,c)\) but is \emph{addressed} by \(x\)'s array position; \(\texttt{sub[c]}\) stores
        \(\dist(x,\mathit{par\_c}[c])\). Mixing the node id and the position -- or putting
        \(\dist(x,c)\) into \(\texttt{sub[c]}\) -- silently mis-scales every difference
        (\(\coderef{c18-summit}\)).
  \item \textbf{A swap moves a node between positions, not between components.} The centroid tables
        are topological and fixed; only \(a[\cdot]\) changes. Delete the node at its old position and
        re-insert at the new one across \emph{all} its \(\bigO{\log n}\) centroid tables, both
        accumulators -- missing the \texttt{sub} update leaves a stale distance in the parent bucket.
  \item \textbf{XOR-decode against the \emph{type-1} answer.} A type-2 query does not update
        \(\mathit{last}\); reading the raw \(a,b,c\) without decoding, or advancing \(\mathit{last}\)
        on a swap, scrambles every later query into \textbf{WA}.
  \item \textbf{64-bit everywhere, but the encoded values fit in \texttt{int}.} Parameters are
        \(<2^{30}\); the \emph{answer} needs \texttt{long long}. Compute \(\mathit{cnt}\cdot W + \mathit{sum}\)
        in 64-bit -- \(\mathit{cnt}\le 2\cdot10^5\) and a distance up to \(2\cdot10^{11}\) overflow
        \code{int} instantly.
\end{itemize}
\end{hardeningbox}

\begin{recallbox}[Chapter 24 in ten seconds]
A rooted DP is blind to paths that cross its root; centroid decomposition removes the blindness by
splitting at the balance point and claiming each path at its top centroid. Build it by
find-recurse-claim (\(\coderef{c18-build}\), height \(\le\lceil\log_2 n\rceil\), \(\thmref{c18-height}\));
count pairs at a distance by query-then-merging an arm histogram and subtracting the same arm
(\(\coderef{c18-count}\), CF 715C \(\coderef{c18-715c}\)); index a component with a per-centroid
distance table for nearest-marked and 2-D-budget queries (\(\thmref{c18-nn}\), \(\coderef{c18-index}\),
CF 293E \(\coderef{c18-293e}\)); fold anything a tag cannot -- products via modular inverse
(\(\coderef{c18-833d}\)). Single-path ranges are HLD (\(\chref{c16}\)), subtree sets are sack
(\(\chref{c28}\)), moving topologies are link-cut (\(\chref{c23}\)) (\(\defref{c18-map}\)). Summit CF
757G (\(\probref{c18-summit}\)) keeps \(\texttt{all}-\texttt{sub}\) Fenwicks keyed by position and
answers \(\sum_{i\in[l,r]}\dist(a_i,v)\) online.
\selftest{
  \item Why must \(\mathit{dfs\_size}\) be recomputed inside every \(\mathit{decompose}\) call, and
        what breaks if a stale \(\mathit{sz}\) is used? \hfill\lv{3}
  \item State the query-then-merge rule and explain why merging first double-counts. \hfill\lv{4}
  \item Why is the nearest-marked min over a node's centroid ancestors \emph{exact}? \hfill\lv{4}
  \item In CF 293E, in what order do you remove the self-match and halve, and why? \hfill\lv{5}
  \item In the Summit, what does \(\texttt{sub[c]}\) store, and why must it be keyed by position
        rather than by node id? \hfill\lv{6}
}
\end{recallbox}

\chaptersection{Tool Kit Inventory}

The machine assembled through Chapter~24. Earlier chapters' rows are carried forward; this
chapter's are appended. Pointers only.

\begin{center}
\small
\renewcommand{\arraystretch}{1.25}
\begin{tabular}{@{}p{0.23\linewidth} l c p{0.22\linewidth} >{\raggedright\arraybackslash\scriptsize}p{0.18\linewidth}@{}}
\toprule
\textbf{Tool} & \textbf{Since} & \textbf{Lv} & \textbf{Trigger phrase} & \textbf{Where} \\
\midrule
\multicolumn{5}{l}{\itshape Carried forward} \\
Centroid, existence, linear find & 12 & \lv{3} & ``the balance point this chapter recurses on'' & \thmref{c09-centfind} \\
LCA / \(\dist\) oracle & 16 & \lv{4} & ``the pairwise distance the index generalises to sets'' & \chref{c11} \\
Fenwick tree & 18 & \lv{3} & ``the bucket structure for dominance and range sums'' & \chref{c14} \\
Heavy-light decomposition & 22 & \lv{5} & ``a single path you UPDATE + QUERY, not all paths'' & \chref{c16} \\
DSU-on-tree / sack & 17 & \lv{5} & ``merge child sets of a subtree'' & \chref{c28} \\
Link-cut / Euler-tour trees & 25 & \lv{6} & ``the topology moves between operations'' & \chref{c23} \\
\midrule
\multicolumn{5}{l}{\itshape Chapter 24 -- new this chapter} \\
Decompose-and-recurse build & 24.1 & \lv{4} & ``remove the centroid, mark dead, recurse on pieces'' & \defref{c18-decomp}, \coderef{c18-build} \\
Centroid tree + log height & 24.1 & \lv{4} & ``every node touches \(\log n\) centroids'' & \thmref{c18-height} \\
Top-centroid pair count & 24.2 & \lv{4} & ``count/aggregate over all paths by distance'' & \defref{c18-claim}, \coderef{c18-count} \\
Distance-to-set index & 24.3 & \lv{4} & ``nearest / marked set, online, fixed tree'' & \thmref{c18-nn}, \coderef{c18-index} \\
2-D dominance in the bucket & 24.3 & \lv{5} & ``two budgets on the path (edges AND weight)'' & \coderef{c18-293e} \\
Feature/accept/fold template & 24.4 & \lv{5} & ``an aggregate a single tag cannot fold'' & \defref{c18-beyond}, \coderef{c18-833d} \\
The tree-path engine map & 24.5 & \lv{5} & ``which tree engine, and when to hand back'' & \defref{c18-map}, \coderef{c18-map} \\
Centroid decomposition kit (one listing) & 24 & \lv{5} & ``all the skeletons in one place'' & \coderef{c18-kit} \\
Summit: all\(-\)sub by position & Summit & \lv{6} & ``sum of distances over a position range'' & \coderef{c18-summit} \\
\bottomrule
\end{tabular}
\end{center}

\chaptersection{Related Topics}

\begin{itemize}[nosep, leftmargin=1.4em]
  \item \chref{c09} (Ch.~12, Height / Diameter / Centroid) --- backward, and the foundation: Ch.~12
        owns the centroid \emph{definition}, its existence, and the linear find
        (\(\thmref{c09-centfind}\)); \(\secref{c18-tree}\) reuses that finder as a black box and starts
        at ``remove it and recurse.'' CF 342E (Xenia and Tree) is Ch.~12's URL, named in prose only at
        \(\secref{c18-index}\).
  \item \chref{c10} (Ch.~13, Tree DP) --- backward/contrast: CF 161D (count pairs at distance \(K\))
        is Ch.~13's \(\bigO{nK}\) subtree DP; \(\secref{c18-count}\) is the centroid upgrade to
        \(\bigO{n\log n}\), re-treated in prose, never re-linked.
  \item \chref{c11} (Ch.~16, LCA) --- backward: the pairwise \(\dist(u,v)\) oracle a single query uses;
        the centroid distance index (\(\secref{c18-index}\)) generalises it to \emph{sets} of marked
        nodes.
  \item \chref{c14} (Ch.~18, Fenwick Trees) --- backward: the BIT is the per-bucket structure for the
        2-D dominance count (CF 293E, \(\coderef{c18-293e}\)) and the Summit's position-indexed
        accumulator (\(\coderef{c18-summit}\)); the BIT itself is never re-taught.
  \item \chref{c16} (Ch.~22, Heavy-Light) --- backward, and a promise \textbf{discharged}: Ch.~22's
        \(\secref{c16-boundary}\) handed path \emph{counting} and distance statistics forward to
        centroid; \(\secref{c18-count}\) and \(\secref{c18-boundary}\) pay it back. The line is the
        unit of work: HLD owns \emph{one chosen path}, centroid owns \emph{all paths}.
  \item \chref{c28} (Ch.~17, DSU on Tree / Sack) --- contrast: sack \emph{merges} child multisets
        bottom-up for a \emph{subtree} statistic; centroid \emph{dissolves} a component at its middle
        top-down. The direction of the merge is the tell (\(\defref{c18-map}\)).
  \item \chref{c29} (Ch.~23, Li Chao / beats) --- forward/contrast: a path whose aggregate is a
        \emph{line envelope} queried at a variable \(x\) is Li Chao's, though it can be hosted inside a
        centroid bucket (\(\secref{c18-boundary}\)'s last table row).
  \item \chref{c23} (Ch.~25, Link-Cut / Euler-Tour Trees) --- forward: centroid decomposition is built
        \emph{once} on a fixed topology and dies the moment edges are linked or cut; that is
        link-cut's job. \textbf{This is the Next chapter.}
  \item \chref{c21} (Ch.~27, Mo's Algorithm) --- forward: an \emph{offline} batch of queries over
        arbitrary node sets with no inverse is Mo-on-trees, not the online centroid index
        (\(\secref{c18-boundary}\)).
\end{itemize}

\begin{lastword}[Last Word --- Chapter 24]
\emph{``A rooted DP can only see the paths that cross its own root, so it is blind to every pair
that straddles it; centroid decomposition lifts that blindness by cutting the tree at its balance
point, claiming each path at the single highest centroid it passes through, solving the pieces
independently, and merging a per-component distance table -- \(\bigO{\log n}\) levels, every node
touched \(\bigO{\log n}\) times.''}

\medskip
Chapter~12 proved a centroid exists and finds one in linear time; this chapter began at ``remove it
and recurse.'' Sections 24.1--24.4 turned that single move into a family of engines: the centroid
tree and its logarithmic height (\(\thmref{c18-height}\)); the query-then-merge arm count that
counts every pair by distance (\(\secref{c18-count}\)); the per-centroid distance table that indexes
a whole component so a mark or a nearest query is \(\bigO{\log n}\) (\(\thmref{c18-nn}\)); and the
feature/accept/fold template that folds products and 2-D dominance the way 24.2 folds plain sums. The
Summit CF 757G (\(\probref{c18-summit}\)) took the index and the top-centroid claim and ran them
together over \emph{positions}, answering an online range-sum-of-distances with adjacent swaps.

\medskip
Carry three reflexes forward:
\begin{enumerate}[nosep, leftmargin=1.6em]
  \item \emph{Ask what the query ranges over.} One path \(\to\) HLD; a subtree \(\to\) sack; a
        per-node state \(\to\) DP; \emph{all paths} or a distance-to-\emph{set} \(\to\) centroid
        (\(\defref{c18-map}\)).
  \item \emph{Claim each path exactly once, at its top centroid, and subtract the same arm.} The
        histogram merge, the 2-D dominance sweep, and the Summit's \(\texttt{all}-\texttt{sub}\) are
        the \emph{same} correction wearing different hats (\(\defref{c18-claim}\)).
  \item \emph{Recompute sizes per component, and honour the dead marker.} The whole \(\bigO{\log n}\)
        bound rests on each piece being \(\le n/2\); a stale size or a re-entered centroid turns it
        into a chain (\(\coderef{c18-build}\)).
\end{enumerate}

\medskip
And three habits, drilled until boring: gather the per-centroid distance table once at build time so
every later query is a walk, not a fresh DFS; subtract the self-match \emph{before} halving an ordered
pair count; and when an accumulator is a product, the same-arm subtraction is a modular \emph{inverse},
not a division. You now own the divide-and-conquer on the tree's \emph{shape} --- the one engine
heavy-light and Li Chao could only hint at for paths that refuse to share a root.

\medskip
\textbf{Next:} \chref{c23} (Link-Cut / Euler-Tour Trees) drops the last assumption this chapter leaned
on everywhere -- that the tree is \emph{fixed}. When edges are linked and cut between operations, no
decomposition built once survives; the splay-based link-cut tree recomputes path aggregates on a
\emph{moving} forest, the dynamic companion to everything in Part~V.
\end{lastword}


\closechapter
